\quellseite{457}
\part{V. The Method of Religious Instruction}

\quellseite{459}
\chapter{§ 32. PRINCIPLES OF METHOD}

*G. v. Zezschwitz ii 2: \emph{Die Katechese oder die kirchl. Unterrichtsmethode}, 1882. K. Buchrucker, pp. 173–236. *Charles De Garmo, \emph{The Essentials of Method. A Discussion of the Essential Form of Right Methods of Teaching}, 1889. Ch. A. McMurry, \emph{Elements of General Method}, 1892. E. Sachse, pp. 368–381. *F. Zange, pp. 142–269. Chr. Burkstuemmer, pp. 175–197. J. Steinbeck, pp. 171–184. *J. Schmarje, \emph{Das katechetische Lehrverfahren}, ²1892. *F. W. Doerpfeld, Ges. Schriften iii 1: \emph{Religioeses und Religionsunterrichtliches}, ²1895. P. Staude, \emph{Zur Anwendung der Formalstufen im Religionsunterricht}, 1903. O. Schoenhuth, \emph{Methodenlehre f. d. Unterricht in Religion}, 1904. Chr. Ufer, \emph{Introduction to the Pedagogy of Herbart}, ed. by Ch. de Garmo, 1894. O. Messmer, \emph{Kritik d. Lehre v. d. Unterrichtsmethode}, 1905. O. Harnisch, \emph{Didaktik u. Methodik d. ev. Religionsunterrichts i. d. Volksschule}, 1906. Rein, Pickel, and Schneller, \emph{Theorie und Praxis d. Volksschulunterrichts nach Herbarts Grundsaetzen. Das I. Schuljahr}, ⁸1908. *K. Richter, \emph{Die Herbart-Zillerschen formalen Stufen}, ³1908. *W. Bittorf, \emph{Methodik des evang. Religionsunterrichts}, ²1908, pp. 1–31. E. Linde, \emph{Der darstellende Unterricht}, ²1910. W. Lawin, \emph{Methodik d. ev. Religionsunterrichts in d. Volksschule}, 1910. *Th. Wiget, \emph{Die formalen Stufen d. Unterrichts}. E. Einführung i. d. Schriften Zillers, ¹⁰1910. E. Thraendorf, \emph{Allgemeine Methodik d. Religionsunterrichts}, ⁵1912. *S. Ch. Parker, \emph{A Textbook in the History of Modern Elementary Education}, 1912, pp. 375–430. *G. Fankhauser, \emph{Die bibl. Geschichte in Sonntagsschule und Religionsstunde}, ³1915. Ch. W. Heathcote, \emph{The Essentials of Religious Education}, 1915. *Sneath, Hodges, and Tweedy, \emph{Religious Training}, 1917, pp. 35–70.

In the preceding parts we have sought to gain a clear understanding of the pupil who is to be taught and trained, of the aim which the pupil is to reach, and of the material that will serve to reach that aim. Now one question remains, viz., by what method the material is to be transferred into the pupil’s soul life in order that the aim in view may be reached. In attempting to give an answer to this question we refrain for the present from studying the method of
\quellseite{460}
instruction as applied to the several kinds of material and the successive stages of development in the pupil, and limit our attention to the principles of method, to things bearing upon the work of instruction as a whole.

From what has been said, it must be evident that only such a method can hope to be effective and successful which adjusts itself at the same time to the distinctive character of the catechetical material, and of the catechumen. The thorough welding of the two is the very object in view—a welding in which the trinity of the inner life of the pupil participates, and the divinely designed purposes of the material are realized. The former must be emphasized over against the exponents of a one-sided intellectualism which has in view merely the enrichment of the pupil’s intellect, as if the feelings and the will could safely be ignored; the latter requires emphasis over against those who, under the spell of the same intellectualism, look upon Christianity merely as a body of doctrines (cf. p. 275), and forget that it is also life—life which produces inner blessedness and strives after expression in word and deed. The harm which has been done and is being done by such one-sidedness, is appalling; and every method must be emphatically condemned that aims merely at imbedding and anchoring the material in the intellect and does not attempt to arouse an interest in the material and does not seek to stir the will to action. When such a method is followed, the pupil whose inner life has been created as a trinity, does not come into his own; and the material does not come into its own which aims to reach and to influence the whole inner life of the pupil.

The catechetical material has for its source the divine revelation; it is not a sum of truths and experiences gathered by man through his own efforts. This being true, the Socratic method which would draw forth from the child’s mind religious truth that was not first put into it, is precluded
\quellseite{461}
at the outset (cf. p. 142). Just as mankind needed the impartation of divine truth, so also the child must first be given the sacred material. \emph{Presentation} is, in consequence, the first element of a correct method. This is also quite in harmony with the nature of the catechumen; for if his inner life is to be set in motion, a stimulus from the world without must be exerted upon him and conducted by the nerves to his soul. Not until then can there be sensations, perceptions, intuitions, concepts — these fundamental prerequisites of the further development of the whole inner life (cf. pp. 203 ff.). The presentation should place the objects so vividly before the soul that they can be grasped and held, seen and examined; that they will be firmly anchored in the mind and easily recalled for the purpose of contemplation; and that they will arouse the feelings and move the will (cf. p. 205, 223, 225, 238). This result will be achieved in proportion as the presentation is fragmentary or complete, sketchlike or distinct, superficial or thorough, incomprehensible or intelligible, colorless or intuitional, abstract or concrete, and in proportion as it is pervaded by sympathy and warmth. — As a rule an explanation of terms will be connected with the presentation. This will keep the teacher from presenting anything to his pupils which they do not—at least in an elementary way—understand. The more the vocabulary of the children is limited, the more the explanation of terms is required of the catechist. He need not, of course, explain every single term, but he must ascertain that the child understands correctly such terms as do not circulate freely in his environment, or words that in local usage have acquired a meaning different from the accepted one. Always the catechist has to ask himself, “Do my children understand my words, so that they not only hear the sound but also connect a definite image and concept with every word?” Again and again he will have to reassure himself of this by addressing questions to
\quellseite{462}
his pupils. The teacher must be particularly careful in the \emph{explanation of terms} in the following three instances. (1) Foreign words require specific explanation even though they have become part and parcel of our language, e.g., evangelical, apostle, sacrament; (2) archaic words or locutions, e.g., take no thought (do not worry), conversation (conduct); (3) synonyms, e.g., tempt and test, anger and wrath. Unless the teacher without ever wearying explains such words, the strangest misconceptions take root in the soul. The rule must be to take as little for granted as possible. But do not devote too much time to such defining: simply circumscribe it, substitute a well known synonymous word for it, let the plain word take the place of a figure, or show wherever the need arises the impossibility of accepting the literal sense (e.g., “the world is gone after him”; “the lion that is of the tribe of Judah has overcome”).

The soul of the catechumen, however, has not only the faculty of intuition, but also that of conceiving (cf. pp. 205 ff.) and of thinking (cf. pp. 213 ff.); it strives to associate the material imbedded in it with mental possessions previously acquired; to analyze it by thinking, to examine it as to its inner relation, to combine the single truths, and thus to form a consistent and harmonious view of life (cf. pp. 214 ff.). Nothing is better calculated to meet this desire of the soul than the catechetical material; for it not only suffers penetration by thought, but even challenges it, and only thus its value can be recognized and can it be made a guide for life. Accordingly, the second step of method must be \emph{penetration}.

\emph{Sentence analysis} is one of the means to this end. At times it may be well to connect it with the “presentation”, but more frequently it becomes the apt starting-point of “penetration”. Such analysis may be either grammatical or logical. The grammatical analysis, in the case of a simple sentence, occupies itself with subject, predicate, object, and the
\quellseite{463}
several modifiers of these main parts. In the case of the compound sentence analysis is occupied with determining the principal and the subordinate clause and their modifiers. In this way the child will readily recognize that the Second Commandment does not speak of the lie as such, or that, in the explanation of the Second Article, everything turns upon the words, “I believe that Jesus Christ is my Lord.” The logical analysis of the sentence strives at making plain the inner connection of sentences, the relation between ground and result, cause and effect, and inquires about aim and purpose (cf. the explanation of the Second Article: whom, wherewith, whereto?, or of the Fifth Petition). Without such analysis of sentences, no exact comprehension of the verbal meaning is possible, for instance, in connection with the Catechism or the hymn; and yet, so very much depends upon such comprehension. In many cases sentence analysis is a good preparation for the penetration proper.

\emph{Development} is another essential element of penetration. The faculty of the soul to determine the inner relation of concepts and to combine them into a logical sequence of thought, is applied to good purpose in connection with this process of development; for it is the function of development to compare the new factors of the material with those already known; to recognize both in their essential characteristics; to associate the new with the old, and to educe the new from the old; to ascend from intuition to the concept; to draw deductions; to form judgments; to link together series of conceptions in order to bring them under the same point of view; and to recognize the idea which underlies the individual facts and combines them into a unit. In other words, it is advisable to follow the old didactic rules, grounded as they are in the inner life of the pupil: (1) From concept by way of intuition to the conception; (2) from the easy to the difficult; (3) from the simple to the complex; (4) from the
\quellseite{464}
near to the distant; (5) from the known to the unknown; (6) from the individual truth derived from the material and from its vivid combination with already acquired truths to a harmonious whole by means of reflection.

We distinguish between two forms of development, the analytic and the synthetic. When we take the conception and unfold by means of partition, the wealth of material contained in it, we follow the \emph{synthetic} method, to use the terminology as correctly restored by G. v. Zezschwitz (Katechetik ii 2 § 26); and when we start from the variety of given materials and, by means of reflection, combine them into a unity, into the conception comprising them all, we follow the \emph{analytical} method. The ancients correctly said that in the former case we advance \emph{a principiis ad principata}, in the latter \emph{a principiatis ad principia}. For instance, we develop the conception “altruism” according to the \emph{analytical} method when we proceed from the story of the Good Samaritan and by means of reflection discover all those individual features that go to make up the conception of “altruism”; we follow the \emph{synthetic} method when we proceed from the conception “altruism” and by way of partition and division arrive at the variety of features contained in this conception. Or suppose that the sentence, “The Holy Ghost has called me by the gospel” is to be developed according to the \emph{synthetic} method. In that case, proceeding from the conception “call”, we say: “To call” means to make one hear who is far away. When the father calls his child, he is not near but afar from him. When he calls him, he wants to show that he should come. Perhaps he wants to tell or give the child something. When the child follows the call, he comes to the father and is with him. Just so the Holy Spirit calls; He calls us men. Now that shows that we are far from Him. Where are we? Whence does He call us? Whither does He call us? Where is he who follows the call and
\quellseite{465}
what is his condition? Whom, for instance, has the Holy Spirit called in this way? When did He call you and bring you to Himself and to Jesus Christ?—The same sentence is treated according to the analytic method when one proceeds, e.g., from the call of Jesus’ disciples or from the parable of the Great Supper, shows what the call is and how it is issued, and finally demonstrates how the Holy Spirit has called us through Baptism and instruction. What, according to the synthetic method, is a mere example, is, according to the analytic method, the source from which everything is deducted and derived.

Although both forms are valuable and likely to be used by an apt catechist, the analytical form of development should be particularly cultivated, for not only does it, more than the other, stimulate the independent intellectual activity of the young; it also supplies the building-stones, as it were, and connects them into one structure. But not only this: it is in keeping with the psychological law that we should ascend always by way of intuition from the concept to the conception (cf. p. 205). A clear understanding and doctrinal firmness result more readily from the use of the analytical method though it is true that it requires more time than the synthetic method. Analytic and synthetic catecheses, if strictly arranged, may be recognized externally at once: the synthetic catechesis, like the synthetic sermon, places the truth to be taught, theme-like, at the very beginning (e.g., Jesus Christ is true God); the analytical catechesis may, indeed, announce a general statement of aim (e.g., Today let us learn what true faith is), but not a theme, least of all the “final theme” (Finalthema) because it is just the object of the catechesis: like the strictly analytical sermon, or Kunsthomilie, it would develop the individual truths from the material and gather these all up into a unity, the theme. The individual truths are more easily gathered up and combined
\quellseite{466}
and they are more firmly impressed upon the mind if the catechist during the process of developing the truths, consistently stresses the main points. It may even be necessary to make a pause now and then for the purpose of an occasional retrospect and preliminary recapitulation.—Finally it must be noted that the definition of a conception which is the goal of analysis or the starting point of synthesis, does not need be as exact as a strictly scientific definition.

\emph{Argument} is another constituent element of penetration. The argument should receive careful attention, inasmuch as a foundation of solid knowledge is to be laid, and the young Christians ought to be able, in later life, to give an account to everyone that shall ask them to give a reason concerning the hope that is in them (1 Pet. 3:15). The proof for the correctness of a statement is either direct, when it takes the form of a quotation from Scripture, or indirect, when it recurs to, and draws conclusions from, the amount of saving knowledge which the children have already gained from Scriptures; or, finally, it can be taken from Christian experience. The first method of argumentation should be plied with particular assiduity. The cultivation of the second has the advantage of training the catechumen, by the employment of his reflective powers, to make use of what he already knows, and of enabling him to see with increasing clearness the inner connection among the several Christian truths, while the catechist, at the same time, is again and again reminded of the duty of organically combining the new with the old. If the argument is drawn from Christian experience, it may be quite permissible to draw upon that of the children, but only with the greatest caution; it is much safer to draw upon that of mature Christians as this has found expression in our hymns and elsewhere; at times, but always with chaste modesty, the experience of the catechist himself may be used as a source of argument (cf. pp. 262, 364). In all cases these
\quellseite{467}
rules must be observed : (1) The argument must be conclusive; the proof dare not be obtained by surreptitious methods: the highest truth is here at stake, and, likewise, the training of the young for truthfulness, and the cultivating of the catechumen’s sense of truth (cf. p. 220). If, in later life, one proof after the other is found to be invalid, the very truth itself is in danger of being lost. (2) Argumenta non sunt numeranda, sed ponderanda.” Especial caution is required when central truths are to be proved, as “God protects me against all danger, etc.”, “Christ has redeemed, purchased, and won me from all sins, etc.”, “not by my own reason or strength”, and also such points as are under attack in the present. This, indeed, does not imply that the tenor of instruction should become polemical, or at least, pronouncedly apologetical. What the catechist will have to do, is to develop and inculcate the disputed truth, whether it is doctrinal or ethical, more thoroughly in a positive manner than commonly so that the viciousness of this or that false ethical principle or doctrine stands disclosed without special effort, and is readily seen by the children. This is the most effective method of making them immune. In the Young People’s Society, the pastor is given sufficient scope to endow his arguments with a distinctly apologetic tenor.

The third element of a correct catechetical method, additional to presentation and penetration, is \emph{application}. Failure to apply apprehended truth to practical life is a common fault with the mature congregation. Knowledge concerning the most important doctrinal and ethical truths is there; but it is not put to work in actual life so that there is often a marked contrast between knowing and doing. For this reason it becomes necessary, already when the young are taught, to inveigh with great seriousness and force against the assumption that all is well when the knowledge of God that has been acquired can be voiced in fine words; or even that
\quellseite{468}
this is the main thing. Rather should the fact be impressed upon the souls of the young as deeply as possible that instruction in religion has absolutely no value for them unless it effects a transformation of the whole life. Application, then, is intended as the path which leads from the schoolhouse into real life, and not only into real life in general, but into the life of the children in particular. Application is to show how that which has been learned can and should be applied in life; it should create a desire to do it, and move them to action. For instance, in exhibiting the characters of Biblical History, it will behoove the catechist to disclose the blessed results of godly action, the baneful effects of wicked conduct, the ultimate incentives and motives of their life. In this way the teacher should lead his pupils to take pleasure in righteousness, to feel repugnance for wickedness, and to long for the powers that give life. But the teacher will not be permitted to stop there; he must take them by the hand, as it were, and show them how, in their youthful life, they can as much as lies in them do that which is good and abstain from that which is evil; he will vividly bring home to them that it is their holy duty to do as they have learned, and that, as baptized disciples, they have received power to do so, and that in the Word and Sacrament they have the means for its daily renewal. Their consciences must clearly realize their duty in the light of Matt. 7:21. Stimulation of moral and religious interests is useless, and may even become positively perilous unless it is followed by action. It would be more profitable not to arouse at all a religious power if one is unwilling to put it to work in the service of God and man (cf. pp. 235, 239 f., 258, 261).

\begin{quote}
\emph{Fankhauser} (\emph{op. cit.}, pp. 169 f.) says on this subject: “Application may be compared to a large electric search-light. The Bible passage, the lesson from history, is the flame. With the reflector of application we catch the light and throw the bright ray wherever we please. We cast light upon the past life of the children, upon
\quellseite{469}
this or that incident. On this or that occasion, did you really have God before your eyes and in your heart? In that particular instance, did you love your neighbor as yourself? Stop to think whether you were not enticed by sinners! By such a method the conscience is awakened, self-knowledge is imparted, and by the grace of God, even contrition and desire for pardon. In the childhood stage there is, as a matter of fact, little profound recognition of sin. But does that imply that we should beware of arousing it? By no means! It will rather behoove us, through the Word of God and His Spirit, to open the eyes of the child to the fullest possible extent. The New Testament truth that even the most trifling wrong-doing is sin, must shine into the child’s heart. And recognition of sin drives one to Jesus. The more the child learns to love Jesus and to desire Him, the deeper will grow his recognition of sin; and in the same proportion will grow the desire for grace and the latter increase in importance in the child’s eyes. — Grace, however, is not merely a pardoning, but just as much a regenerating, dynamic, manifestation of divine power. Many self-satisfied saints, alas! take a rest when they have assured themselves of the forgiving grace of God. Let us never neglect to show that the real ultimate aim of grace is to make us obedient children of God! For this reason reverse the search-light and turn it upon the life lying ahead of the child. What would you do, should you be in this or that difficulty? How should you act if your comrade threw slurs at you? etc. We picture to the child conditions in which he may actually find himself some day, and induce him to determine for himself what he should do in such conditions, in order to remain obedient to the Lord Jesus. In this way application is intended to cast light upon both the past and the future of the child’s life: upon the past in order to effect self-knowledge and tenderness of conscience; upon the future in order to clarify the soul’s eye so that it may see how love to the Saviour should receive a practical demonstration.”
\end{quote}

Inasmuch as the aim of religious instruction is not the production of merely momentary impressions, but the impartation of permanent possessions, drilling in the subjects taught and apprehended must be named as the final element of a correct method. It is true that this element must not invariably follow the factors of method already named: beside the penetration of the particular lesson
\quellseite{470}
a special drill may be superfluous, but in the whole process of instruction it is indispensable; even the so-called immanent review (cf. p. 410) is not an adequate substitute for it. While it may not be necessary, in some branches, to conduct a regular drill in connection with every single lesson, it is indispensable even there for the purpose of recapitulation and for weaving the several lessons into an organic unity (cf. pp. 269, 298 ff.). Of course, the drill should not be allowed to degenerate to a mechanical quiz; the catechist should spare no effort, by means of free recapitulation, inversion, and transformation, to assure himself that the matter has been mentally assimilated, and thus to equip his pupil for mastery in the use of it. Whole series of individual lessons can be placed under a common view-point; for instance, Abraham’s life under that of faith (cf. p. 297); the activity of Jesus under that of the prophet and deliverer from sin, death, and devil (cf. p. 305), or under “faith as the only condition of entering His kingdom (cf. p. 304). In catechetical instruction so-called catechetical “excursions” are practicable, upon which important conceptions are traced through the whole Catechism, special attention being paid meanwhile to the shades of difference the conceptions receive from the connection in which they are found. Such excursions, naturally, can be successfully undertaken only where, after the completion of the daily lessons as well as larger sections, recapitulation and drill have already done their part. The rule to be followed is: “When a review has become a necessity, it is already too late.” — In connection with this whole chapter compare also chapters 34–36.

The Herbart-Ziller School advocates the theory of the \emph{Formal Steps}, according to which every lesson should be carried over five steps: (1) preparation (a. statement of aim, b. preparation proper), (2) presentation, (a. presentation of facts, b. penetration), (3) association or combination (comparison,
\quellseite{471}
abstraction), (4) generalization (i.e., formulating a statement of the general principles which have been worked up to in step three), (5) application. It must be conceded that these grades conform to the process of mental acquisition as it actually transpires; for here the faculties in operation are apperception (preparation and presentation) and abstraction (association and generalization) cf. p. 212. It will be highly necessary, therefore, for the beginner to ask himself whether he always prepares his pupils for the truth to be taught; whether he presents it well; whether he always combines the part with the whole; whether he carefully collects the elements gained in the preceding steps into a brief, pregnant, generalizing statement; whether he properly applies the truth derived to the life of the children. Notwithstanding, when the “theory of formal steps” is applied to every kind of material and every single lesson in all branches, the peculiar nature of the material is overlooked, usually with the result that the manner of teaching becomes stilted or wooden. Ziller himself did not hesitate to affirm that “partition according to the formal steps is inapplicable where the material is already of a ‘conceptional nature’, and has a trend preponderatingly doctrinal”. The formal steps are well adapted to historic material. But even here the question is in place whether we had not better content ourselves with the three steps preferred by Doerpfeld: (1) Presentation for the purpose of intuition; (2) Penetration by means of the thinking capacity; (3) Application to the life of the pupil, provided only that the true and necessary features intended to be brought out in the steps of “combination” and “generalization” are permitted to come to their own. These steps of Doerpfeld correspond to what has above been said concerning presentation, penetration, and application; for there is no doubt that these three are inseparable from instruction as a whole, though they need not be in evidence in every lesson.

\quellseite{472}
\chapter[§ 33. THE METHOD: ITS FORMAL ASPECTS]{§ 33. THE METHOD ACCORDING TO ITS FORMAL ASPECTS}

*G. v. Zezschwitz ii 2: \emph{Die Katechese oder die kirchliche Unterrichtsform}, 1872. *L. Kraussold, \emph{Katechetik}, pp. 358–421. J. Gottschick, pp. 163–175. G. H. Gerberding, \emph{The Lutheran Catechist}, pp. 103–109. J. H. Herzer, Katechetik, pp. 174–304. J. Schmarje, \emph{Das katechetische Lehrverfahren auf psychologischer Grundlage}, ²1892. O. Schoenhuth, \emph{Methodenlehre fuer den Unterricht in der Religion}, 1904. W. F. Young, \emph{The Art of Putting Questions}, 1895. J. G. Fitch, \emph{The Art of Questioning}, 1897. R. Stevens, \emph{The Question as a Means of Efficiency in Instruction}, 1912. *Th. Schmauk, \emph{How to Teach in Sunday School}, 1920.

When the principles governing the method of catechetical instruction have been laid down, nothing has as yet been settled in regard to its formal aspects; and yet, it is of the utmost importance that the right form of instruction shall have been selected if the results looked for are to be achieved. Didactic science in general distinguishes between the acroamatic, or lecture, form and the erotematic, or interrogatory, form. It was a mistake that the acroamatic form was exclusively used in religious instruction down to the days of rationalism (so far as the question was at all in evidence, it was strictly examinational and confessional; pietism added the analysing question which differed, however, essentially but little from the former; cf. p. 135). The “Socratic” Rationalists were quite as much in error—considered from the standpoint of principle even more so—when they discarded the acroamatic form altogether (p. 142), and confined themselves strictly to the erotematic form. Even a cursory examination of the nature of the catechumen and of the catechetical material—the only spheres that can contribute to the solution of the question—proves the necessity for both forms.

The \emph{acroamatic} form is required both by the nature of the material and that of the catechumen, and this for the same
\quellseite{473}
reasons which require presentation to be an element of the method. But the other elements of the method, too, require now and then this form, even while the questioning is in full swing. The \emph{erotematic} form is required, not so much by the peculiarity of the material as by that of the catechumen. The effort of following closely the unbroken lecture of the teacher would be too great for him; he would lose interest and attention, and consequently fail to be influenced psychically; but the erotematic form of teaching compels and enables the children to be attentive, encourages them to co-operate with the teacher, establishes a status of mental reciprocity between the catechist and the catechumens, and thereby helps them to lay firm foundations for a really Christian view of life. Moreover, the question put to the child and the expectation of an answer express the confidence that he is capable of forming correct judgments, and gives him an opportunity to furnish, in his own words, the evidence that he has mastered the subject, or that (in the case of confessional questions) he takes the right attitude toward it. As to the catechist, he can not control the mental co-operation of the children until he begins to ply them with questions. Every answer received serves as an incentive for him to descend more fully to the level of the children and enables him to teach and to hold the attention of children of various degrees of mental capacity; misunderstandings can be recognized and corrected at once; and the entire class can successfully be taught at the same time. A mistake though it is, accordingly, to find the erotematic form implied in the very word κατηχεῖν (p. 3), it is anything but a mistake to conclude that the nature of the catechumen and practical pedagogic considerations urgently require the erotematic form of teaching. Nor is it exclusively the examinational and its subordinate species, the confessional, question which is required, but rather the \emph{didactic question}.

\quellseite{474}
The proper combination of the acroamatic and erotematic forms of teaching is the true method. Both forms are likely to be found in nearly every catechesis although the one which happens to preponderate impresses its character upon the catechesis. Where a certain truth may be taken for granted, a mere reference to it being sufficient; where tasks are imposed upon the thinking faculty; and where a thought is to be independently developed by means of inferences—in other words, chiefly, though not exclusively in “penetration”—, there the erotematic form is to be employed; in other cases, the acroamatic form will predominate.

Since, then, catechetical instruction demands the use of the question, the \emph{question} must be considered as to its \emph{nature, its various kinds, and its essential attributes}.

It took almost a century to develop a clear and precise conception of the nature of the catechetical question and of the question in general. At first the question was defined as an incomplete sentence to be completed by the person addressed. Dinter went so far as to say: “The question arises when I leave out one or several parts of the sentence and ask the person addressed to supply the missing parts.” Unfortunately only too frequently such questions are formulated (e.g., For God so loved ——? God is ——?). This, however, does not make them real questions. No, the question is, grammatically, a complete sentence, containing all its essential elements, subject as well as predicate; for instance, Who has created the world? The truth was more closely approximated when it was said that the question is an incomplete judgment to be made complete by the person required to answer. This definition was based on the observation that the question has its own subject and predicate, occasionally also an object, etc.; but that it is devoid of content, which, when found, completes the judgment. But that the question becomes a judgment through the addition of
\quellseite{475}
the missing content, does not prove that it is a judgment before that or even an incomplete judgment. It is true, however, that the question is based upon a judgment—a fact that becomes evident as soon as an indefinite pronoun is substituted for the interrogative pronoun; for instance, “someone has made the world” for “who has made the world?” In the process of changing the judgment into a question it ceases to be such in form, having become a demand upon the person questioned to form a judgment for his part, and thereby more fully to define the indefinite judgment on which the question is based. If but a part of the judgment on which the question is based has been left indefinite, the object must be to more fully define that part; if the whole judgment has been left indefinite, the object must be to decide, that is, either to affirm or deny the judgment. And because a sentence is outwardly disclosed as a question through the tone or the arrangement of the words, the following definition is warranted: The question is a demand, expressed in the interrogatory sentence itself through the tone and the verbal arrangement, to form a judgment, and thus either more closely to define the indefinite part of the judgment upon which the question is based, or to affirm or deny the judgment as a whole. The judgment solicited from the catechumen (that is, the closer definition of the indefinite judgment upon which the question is based, or the decision as regards its correctness or incorrectness) is called the \emph{quaesitum quaestionis}, in distinction from the \emph{datum quaestionis}, whereby we understand the fulcrum supplied to the person questioned either in the question itself or through the development culminating in the question. This datum quaestionis forms, for the pupil’s thinking faculty, the basis from where to find the quaesitum. Inasmuch as the question calls for the formation of a judgment, the datum of the question is tantamount to the major premise of the syllogism.

\quellseite{476}
It must be clear, then, that much depends upon the right relation between the datum and the quaesitum of the question.

The correct definition of the question suggests its several kinds: When in the judgment underlying the question only a part of the whole has been left undefined, there arises the so-called \emph{defining question}, less aptly called also supplementing question. But if the whole judgment has been left indefinite, it remains to decide the validity of the judgment by affirming or denying it; thus arises the \emph{deciding question}, also called “yes and no” question. Besides these two kinds of questions, there are no others, as far as essence is concerned; but they may be variously applied. Thus we can speak of causal, final, relation, comparison, preparatory, auxiliary, questions etc.

The defining question can easily be distinguished from the deciding question: in the former the interrogative pronoun, in the latter the predicate or a part of the predicate begins the sentence. While the distinction here made applies to all questions, no matter what the sphere, it certainly applies also to the \emph{didactic question} which alone concerns us here. This question distinguishes itself from others that may be put during the process of catechetical instruction (e.g., the question for information, the recitation question, the rhetorical question, etc.), in that it is not put in the interest of the questioner but of the one questioned, whose knowledge and judgment are to be furthered.

Every species of didactic question, the defining question as well as the deciding question, is, from the aspect of its content, either \emph{examinational or developing}. The examinational question, to be carefully distinguished from the recitation question, will always arise when assurance is sought that something has really become the mental possession of the catechumen and may easily be recalled. This question is, therefore, in place in the drill and at the examination proper. The developing
\quellseite{477}
question, however, is in place especially, though not exclusively, in connection with penetration. This is the question which should by far predominate in catechetical instruction. The confessional question, which is sometimes treated as a distinct species, must be subordinated to the examination question.

But when should the question take the form of the defining question, and when the form of the deciding question? The deciding, or “yes and no” questions have a peculiar history. Both orthodoxism and pietism treated them as the normal catechetical questions; but the Socratic School rejected them as altogether useless, on the ground that they failed to stimulate mental effort. Both went to extremes. There is no doubt that the defining question should be employed with greatest frequency. Excessive use of the “yes and no” questions will result in obtuseness, rather than in stimulation and training for mental cooperation; but they cannot be dispensed with altogether. This species of question is in place, e. g., where the teacher needs a statement which is to serve as a premise for further conclusions and which he does not care to develop once more by a series of defining questions; still more, where a moral decision on the part of the catechumen is the object in view; where a confession or a vow is to be made; or where the catechumen is endowed with sufficient mental maturity at once to substantiate his yes or no with a proof. To reject it entirely would imply a failure to recognize the conversational character of catechization. The beginner should indeed avoid its use as much as possible lest he should use it at the wrong place. The defining questions, in this as in every case, are the best.

But, to be effective, a question must possess certain definite \emph{attributes}. Which are they? A reference to the purpose of the question will help to indicate them. Each question is put for the purpose of obtaining an appropriate response; but every question also aims to bring the pupil nearer to his goal.
\quellseite{478}
Unless the catechumen has understood his teacher correctly, no appropriate answer can be expected. On the one hand, then, the question should be clear, unambiguous, perspicuous. On the other hand, I cannot come one step nearer to the goal with my question unless it is appropriate to the plan of the whole and to the material immediately preceding. \emph{Perspicuity} and appropriateness are accordingly the most necessary attributes of a good question.

Perspicuity is required both in a grammatical and in a logical sense. In the grammatical, as in the formal sense in general, a question, in order to be perspicuous, should possess (1) conformity to linguistic usage; that is, the catechist must come down to the mental capacity and vocabulary of the children, and, at the same time be intent upon augmenting, raising and ennobling it. He should cultivate popular every-day language in a noble sense; but every-day language in an ignoble sense or slang, he should neither use nor permit. Above all, he should beware of the use of technical terms of which he is not sure that the children understand them; (2) correctness of speech as found in the proper combination and arrangement of words. The teacher will, e. g., refrain from putting the interrogative too often at the end of the sentence. Failure to formulate the question as a complete sentence is intolerable; for instance, Honor thy father and ——? That tends to degrade the children to parrots and the teacher to their trainer, all the more, if he should smooth the way to the answer by putting the initial letter of the expected word upon the lips of the children. On the same level with the impropriety just censured is the other, which cannot be criticized too severely, of interspersing certain words or phrases, or even inarticulate sounds, without sense and reason between syllables or clauses, as now, and, but, then, etc.; and (3) simplicity of expression. A question may have been formed in perfect accordance with the rules of grammar, and yet possess a form which prevents its
\quellseite{479}
being perspicuous to the child. Instead of being simple it is clumsy and heavy. The question should be simple in regard to syntax—a principle which rules out long-winded periods, in fact, periods had best be absent altogether. It should be simple also in point of content; that is, the question dare not be really a double question, requiring a double answer, or an answer that requires the simultaneous forming of two judgments; for instance, Who made the world out of what? If such a double question is at all permissible, it is only with more mature pupils for the sake of brevity. Perspicuous in regard to form, the question should be perspicuous also logically, or in point of content. The teacher himself must be perfectly clear concerning the substance of his question, not merely in general, but down to the last detail; and the question as formulated by him must express just what he meant to say. Nowhere is superficiality and partial mastery more mischievous than here. Many a man did not learn what is meant by a thorough and all-sided mastery of a subject and its precise formulation, until he became a catechist and daily had to grapple with this task in the schoolroom. Only a precise and perspicuous question will spare the child the trouble of choosing one of a number of appropriate answers and will lead him directly to the one and only possible answer.

This brings us to the second chief attribute of a good question, that of \emph{appropriateness}, which does not only mean the adaptation of the question to the capacity, the age, etc., of the catechumens, but also to the train of thought of which it is a member. Nothing alien should be introduced with the question into the steady development of the thought, as that would have the effect of deflecting the child’s attention. Nor should any question be asked for which the preceding discussion has not sufficiently prepared the child. In short, there must be no “jumps”. The preceding development must supply the fulcrum upon which the child, in his search for an answer, may
\quellseite{480}
take his position, and from which the correct answer may and will be found. It is quite likely that, in normal conditions, the larger part of wrong answers can be traced to the failure to adjust the quaesitum quaestionis and the datum quaestionis; that is, to inappropriateness. If consideration for the children should, at any time, have led to a digression, it is best, by a few comprehensive sentences, to re-establish the connection with the main thought. What we have said of the appropriateness does not apply with the same force to the examinational question: this must frequently omit subordinate points and must not again become a developing question.

Finally, it must be noted here that the question should, as a rule, be addressed to all the children. That arouses attention. But never should the answer be given by all the children in unison; such a procedure may lead to disorder. When, however, the catechist directs the whole class to repeat an answer just received from one child this danger is not incurred. When the subject is of special moment, such a measure is even heartily recommended. It is also advisable that the questioner, during the process of questioning, should occupy one and the same posture in sight of the children so that he may never lose their attention. Artificiality and stiffness, to be sure, are to be avoided.

Many sins are committed against the laws underlying the formulation of questions. Hence we add, for the purpose of critical study and also for training the eye to discover one’s own faults, a number of questions almost all of which are to be condemned as defective or altogether wrong: The First Commandment forbids what?—When one worships anything else but God, we call it what?—Since you yourself do not merit the grace of God, upon what, when you repent of your sins while desiring forgiveness of the same, do you hope to base your confidence?—What kind of being must he be who governs the world, which contains so many beings of which each one has a will of its own, and so many million heavenly bodies, many of which are much larger than our earth?—Who created the world, and who preserves it?—Who fulfills the Seventh Commandment better, he who steals or he who does not
\quellseite{481}
steal?—Which is the greater sin, to steal or to be poor?—When you are once in heaven, shall you, too, be blessed?—How may one guard against the sufferings that have come upon him?—Who died on Calvary?—What did Jesus redeem?—What has procured for us redemption?—What does the Christian do when he is in trouble?—Whom did God send into the world?—When did Christ die?—Luther was born where?—The Law was given when?—What does this lesson treat of? (Of Abraham.)—Jordan flows into…?—Capernaum is situated on…?—Then Joseph could no longer? (Refrain himself.)—What must the oath not? (Be broken.)—What must man not do wilfully? (Court danger.)—When did Alexander desire to make peace?—What did Lee, knowing the widespread discouragement in the Northern states and hoping that a brilliant stroke as near New York as he could get might terrify the Northern bankers, decide to do? (Attack General Meade.)—Where is Jerusalem situated?—What are we on earth?—What kind of faith did A. H. Francke have when he performed his works of love? (An active faith.)—When John lay in prison, he thought what?—When you think of God’s government of the world, He is what?—Judas, in leading the enemies to Christ, was guilty of be . . .?—Peter, declaring that he did not know Him, was guilty of de . . .?—What must everyone be that wants to be a Christian? (Baptized.)—What can a man not do for his life? (Ransom it.)—What should married people never? (Forsake each other and separate.)—What must this promise not be? (Broken.)—When you want to make children happy, you must? (Give them something.)—All the creatures of God, in what will He hold them? (In honor.)—If it were true that the world has not been created, what would it have to be? (Eternal.)—In what circumstances can one not reap when he is old?—When is it allowed to kill?—What are departed human beings, even animals, by the heathen? (Worshipped.)—Why are the good works of the heathen, no matter how good, no good works?—What do you call the man who is even more saving than the saving? (A miser.)—Which commandment makes a demand upon us regarding the honor of the neighbor?—What is the difference between the frugal man and the miser in regard to benevolence?—Whereby was the unfortunate man aided? (By the Samaritan.)—Since only those are true Christians who, in view of their Baptism, remain true to their faith, what does it mean to believe in God?—Who is subject to the penalty of decay? (The body.)—Whereby did God most gloriously reveal Himself? (Through Christ.)—Wherein should we take our refuge? (In God.)—What do we sin against when we steal? (Against God.)—How did the Israelites eat the quail? (Raw, immoderately.)—
\quellseite{482}
What is the real destiny of man?—What works contrition in man?—What does the sick man feel in regard to health? (Desire.)—What did Jesus do after being baptized in Jordan?—Who smiled upon Greece? (An ever blue sky.)—What kind of disposition should we possess if we wish to become like God?—On which day did Jesus die? (On a Friday.)—On what day did Jesus die? (On the fourteenth day of the month Nisan or on Good Friday.)—Why should we do good works? (To prove our faith therewith.)—Why should we do good works? (Because we owe God thanks for His mercies upon us.)—Therefore we should earnestly seek to remain true what?—Therewith Jesus means to tell His disciples that who stood bodily before them?—The devil goes about like a rrr....?—What did the sermon of Peter produce upon the visitors at the feast? (A deep impression.)—What, according to the Third Commandment, is to be kept holy?—Who especially has seriously sinned against Jesus?—What was Luther?—How many men are sinners?—What is the first thing we do, and what is sin?—What kind of people does Jesus not have in view when He speaks of the poor in spirit, and what kind of people does He have in view?—Which is the surest way to heaven?—What did Jesus give up when He hung upon the cross?—For which blessing have we reason to thank God every morning when we are privileged to rise hale and hearty? (For good health.)—Is the bad conscience an agreeable or a disagreeable feeling?—How many gods are there?—What man is not mortal?—Does prayer really give powers which we could not have obtained without?—Was the suffering of Job a retribution or a test or a testimony?—Would you have the confidence in yourselves to remain more steadfast than Peter in view of such a temptation?—If Jesus should ask you, “Will ye also go away?”, what would your answer be?

How is the following series of questions to be judged from the standpoint of the developing catechization? The question under consideration is: “What does it mean to fear God?” Acting under the presumption that the pupil cannot at once answer this main question, the catechist puts the following leading questions: “1. An obedient child does not have to be in fear of his father. How does such a child act toward his father’s command? He carries out the command.—2. But what is his attitude when the father forbids something? He keeps from doing what the father forbids.—3. That is the way an obedient child acts. But alas! there are children that do not do what the father wants them to do. What are such children called? They are called disobedient.—4. No good father can overlook the child’s disobedience. What is it that the child must expect in consequence
\quellseite{483}
of it? He has to expect punishment.—5. From this you see what a child has to look for from his father when he fears him. When a child fears his father, he looks for punishment from him.—6. Punishment for what? For his disobedience.—7. Now give me, in a complete sentence, an answer to the question: ‘What do we mean when we say: a disobedient child fears his father?’ We mean: He looks for punishment for his disobedience from his father.”

The center of gravity in the catechetical dialogue is found in the \emph{answers}. To obtain answers: correct answers—in agreement with the question in point of form and content; true answers—given not by accident, but the result of reflection; correctly worded answers, if at all possible—this is the purpose behind the question. Questions are intended to serve as evidence of the understanding and judgment of the child; and, here and there, they serve also the purpose of eliciting a confession. But inasmuch as answers are often wrong, untrue, or faulty, and sometimes, not forthcoming at all, there arises the question as to the \emph{treatment of answers}. Always to find the right method in this respect, presupposes skill, presence of mind, and a thorough mastery of the subject and of self.

If the answer is correct, let the questioning proceed forthwith, unless there is occasion for a pause because the results attained require summarization, or possibly, for a word of encouragement or exhortation. A word of commendation, interspersed here and there, is quite in place in the case of timidity, or when the catechist sees that he has previously reproached the pupil too severely for having given no answer or a faulty one. But too many words of commendation are to be avoided; for, if the question was clear and to the point, a correct answer should naturally be looked for—a fact which the child should be taught and trained to understand. If another method is chosen, there is danger of cultivating the false ambition slumbering in the child. While the individuality of each child should receive due consideration—
\quellseite{484}
and no catechist can afford to remain indifferent to such an important factor — this must remain the rule: The catechist should produce the conviction in the pupil that the question of the teacher is always appropriate; that it is never too exacting, and that it is his manifest duty to give a correct answer. In that case, to be sure, the teacher must be exacting toward himself and see to it that his questions are in every way perspicuous and appropriate. This demand becomes a challenge for self-examination, especially when the answer received is wrong or faulty, or when the question remains unanswered. If the resultant self-examination leads to the conclusion that the pupil’s failure is explained by that of the teacher, let the question be repeated in a perspicuous and appropriate manner; let the teacher even retrace his steps a little, in order to give the pupil the necessary basis. If even then no answer be given, a word of encouragement is in place, or the teacher should pass on to another pupil. Stubbornly to insist on squeezing an answer out of the child, is altogether wrong. But if no progress is made with the whole class, the catechist must once more find the fault in himself and make another attempt, by endeavoring to come down into the concept sphere of the child and to cultivate the simplest of language. This is true of the developing question in particular. If the developing question is in every way distinct and appropriate and again and again elicits a wrong answer, the fault is in most cases to be found in a lack of attention and concentration on the part of the pupil. In that case a short, earnest word of reproof is quite in place. On the occasion of a simple quiz or a drill, and especially in connection with the examinational question, failure to receive an acceptable answer to a question in every way distinct and appropriate, unless occasioned by a lesson of inordinate length or by a lack of proper presentation and direction of judgment, usually finds its explanation in a lack of attention or of diligent preparation
\quellseite{485}
While, even in that case, scolding and railing are out of place, a few brief and earnest words of reproof are certainly in place. There are times when the undisguised sadness of the teacher over his pupil’s shortcomings is even more effective than a reproof. Much, of course, depends on the individuality of the child. The teacher should also ascertain whether the disappointing results are not accounted for by the home life of the pupil rather than his own dereliction. Should the teacher’s suspicion prove founded, he must make due allowance for the new factor lest he become unjust in his treatment of the child. The best rule will be that the children must inform the teacher, before class, of their lack of preparation whenever circumstances have made it impossible. To treat with disdain—in religious instruction!—poorly endowed children or to ignore them altogether constitutes a grave sin against those unfortunates. A true catechist will, never as much as then, bring great affection and attention to bear upon them; will never weary of dealing with them, and strive to make them understand at least a minimum. Just such pupils as these often reward their teacher by a faithful attachment to his person as well as to the truth which they have learned at his feet. Disobedient and mischievous children may be won by a wise alternation of mildness and severity, and by redoubled diligence in preparation. Even then the teacher will not be spared grief (Hebr. 13:17).

\emph{Faulty or partially correct answers} are entitled to special attention. They dare not be summarily rejected nor accepted as correct. The former proceeding would be a wrong against the catechumens, the latter against the truth itself, with a bitter penalty as the possible outcome. The right way will be to establish the element of truth found in the answer; to make that, when found, the connecting link for further efforts, and thus to bring out what is correct through elimination
\quellseite{486}
of what is false. This might take the form of pointing out the evil consequences bound to come should the error be allowed to stand. If the answer is correct in point of content and objectionable merely in point of form, the latter should be corrected by the teacher himself or by another child; and this as briefly as possible; for the religious period should not become a lesson in language. Only when the same faulty expression occurs again and again, the child should be earnestly held to correct his fault. Sometimes ludicrous answers are given; but not all of them are wilful or intentionally frivolous. Should this be the case nevertheless, they must be repudiated with rigor, though never without dignity. Should they have been the result of inadvertence, self-control is absolutely called for, lest the teacher should laugh with the whole class or even incite them to laughter. If the catechist is master of the situation, the ludicrous answer may serve as an occasion for a fruitful remark; and the occasion for mirth will at once disappear.

\quellseite{487}
\chapter[§ 34. THE METHOD IN BIBLICAL HISTORY]{§ 34. THE METHOD OF INSTRUCTION IN BIBLICAL HISTORY}

Compare the literature in § 25 and 32. A. Eckert, pp. 79–118. J. Gottschick, pp. 175–181. O. Harnisch, pp. 63–64. *W. Bittorf, \emph{Methodik}, ²1908, pp. 32–118. J. Berndt, pp. 78–95. R. Kabisch, pp. 182–219. J. Steinbeck, pp. 185–193. *F. W. Doerpfeld, \emph{Ges. Schriften iii: Der Religionsunterricht}, ²1895. P. Staude, \emph{Die formalen Stufen i. bibl. Geschichtsunterricht} in Kehr’s Paedag. Blaetter, 1884. L. Hohmann \emph{Methodik d. ev. Religionsunterrichts}, ²1904. *F. Gansberg, \emph{Schaffensfreude}, ³1909. *H. Scharrelmann, \emph{Weg zur Kraft}, ⁷1910. *O. and E. Zurhellen, \emph{Wie erzaehlen wir den Kindern die biblischen Geschichten}, ²1910. Max Paul, \emph{Fuer Herz und Gemuet der Kleinen}, ⁶1911. M. Reu, \emph{Grundsaetze z. Herstellung von Sonntagsschulliteratur} in Kirchl. Ztschr., 1911. *M. Reu, \emph{How I Tell the Bible Stories to My Sunday School}, 2 vols., ²1926. \textsc{Teaching Helps for All Grades:} K. Francke, \emph{Weide meine Laemmer. Die Heilige Geschichte der Jugend erzaehlt}, 1897. O. Zuck, \emph{Der gesamte Religionsunterricht}, 3 parts, 5 vols., ⁵1902 ff. R. Staude, \emph{Praeparationen zu der bibl. Gesch. d. Alten u. Neuen Tests.}, 3 vols., ⁴1905 ff. E. Thraendorf and H. Meltzer, \emph{Der Religionsunterricht}, 3 parts, ²1905 ff. *Reukauf and Heyn, \emph{Evang. Religionsunterricht}, vol. iii to ix, ⁴1907 ff. Kessel and Spanuth, \emph{Praepar. f. d. ev. Religionsunterricht}, 3 parts, ²1910 f. *G. Fankhauser, \emph{Die bibl. Gesch. in Sonntagsschule u. Religionsstunde}, ³1915. G. Fankhauser, \emph{Christ d. Retter. Gesch. aus d. Leben Jesu. kleinen u. grossen Kindern erzaehlt}, 1925. *G. Fankhauser, \emph{Der Weg zum Kinde. Monatsschrift f. christl. Erziehung u. bibl. Unterricht}, 1915 ff. \textsc{Teaching Helps for the Instruction of Beginners:} G. L. W. Koehnlein, \emph{Bibl. Geschichten f. Kinder von 4 bis 9 Jahren}, 1854. J. F. Ranke, \emph{Des Kindes erster Unterricht aus Gottes Wort angeschlossen an die 40 Bilder der Kaiserswerther Bilderbibel. Eine Vorstufe z. d. ersten Religionsunterricht}, ²1873. F. Wiedemann, \emph{Wie ich meinem Kleinen die bibl. Geschichten erzaehle}, about 1865. Graefin Poninska, \emph{Biblische Geschichten. Nach d. Schnorrschen Bilderbibel erzaehlt f. kleine Kinder}, 1879. L. Wangemann, \emph{Der erste biblische Anschauungsunterricht. Anweisung z. Gebrauch der 20 Anschauungsbilder}, ²1907. P. Staude, \emph{Praepar. f. d. ersten Religionsunterricht}, 1889 f. A. Gieseler, \emph{Der Religionsunterricht auf. d. Unterstufe}, 1900. G. Kaelker, \emph{Bibl. Anschauungsunterricht},
\quellseite{488}
²1905. E. Zeissig, \emph{Von d. bibl. Gesch. auf d. Unterstufe in Katech. Zeitschr.}, 1907. R. H. Evers, \emph{Die bibl. Geschichte u. deren Behandlg. auf d. Unterstufe}, ³1909. M. R. Unger, \emph{Beitrag zur erbaulichen Behandlg. des Religionsunterrichts auf der Unterstufe der Volksschule}, ab. 1912. *M. Reu, \emph{Fuer die Kleinen}, 1913. K. Schneeberg, \emph{Uns’ Herrgott un sin Lued}, 1914. Mrs. H. Gaskoin, \emph{Children’s Treasury of Bible Stories}, 1896. James Baldwin, \emph{Old Stories of the East}, 1896. R. Bird, \emph{Jesus the Carpenter and Joseph the Dreamer}, 1901. Laura E. Ragin, \emph{Kindergarten Bible Stories}, 1905. Geo. Hodges, \emph{The Garden of Eden, The Castle of Zion, and When the King Came}, 1904 ff. R. G. Moulton, \emph{The Art of Telling Bible Stories}, 1904. *Sara Coni Bryant, \emph{How to Tell Stories to Children}, 1905. Rich. T. Wyche, \emph{Some Great Stories and How to Tell Them}, 1910. *E. P. St. John, \emph{Stories and Story-Telling}, 1910. M. Reu, \emph{For Beginners. Fifty Stories for the Little Ones}, 1916. M. Affolter, \emph{Kommt, Kindlein, hoert}, ²1926. Leipziger Lehrerverein, \emph{Im Strome des Lebens, Sammlung von Erzaehlungen und Geschichten zur Belebung der religioesen Jugendunterweisung}. Sneath, Hodges, and Tweedy, \emph{The King’s Highway Series}, now 8 vols., containing much valuable material for telling the Bible stories and for the application upon the life of the child.

That the material assigned for any particular lesson in Biblical History should not be loose conglomeration of isolated events, but a catechetical unit, has already been shown (p. 424). If this statement is expressly repeated here, it is done for the reason that the truth involved is the postulate for all that is to be said in this paragraph concerning the method underlying instruction in Biblical History. While it is quite permissible in the upper grades or on the occasion of a review to place a number of narratives under a common point of view and thus to weave them into a unit, the catechetical unit for the first treatment is always the single narrative. How to treat the single narrative, then, will be the chief problem. Many are of the opinion that the main object of instruction in Biblical History is merely the memorizing of its principal contents, for the one purpose of having an opportunity to refer, in sermon and catechesis, to the material stored up in the memory. If that view were correct,
\quellseite{489}
the science of catechetics might well refrain from investigating the various methods as to their appropriateness. In that case we might be satisfied with reading the story several times and conducting a quiz for the purpose of ascertaining whether there has been an appropriation of the contents. But if instruction in Biblical History is to do its share toward reaching the aim of religious instruction and training as it is given on page 280; if it is to be of present value for the children, if it is to influence their whole inner life and to implant the desire to live according to God’s will, if it is to become a dynamic factor in their young lives, then the question as to the didactic treatment of the individual biblical narrative proves of uncommon importance, and requires specific treatment.

In order to attain this purpose, there will have to be in evidence what we described in §32 as the essential features of the whole process of instruction. In other words, the treatment of the individual narrative should cover the three steps of presentation, penetration, and application. The penetration will also comprise what the Herbartians call “association” (cf. p. 448) in case the latter seems pedagogically valuable; and invariably it will issue in the “generalization” of the specific truth educed from the individual narrative (fourth step of the disciples of Herbart). Again, presentation should be preceded by a statement of aim and frequently also by a brief preparation.

\section{1. Statement of aim}

A sermon that pursues only the general aim of edifying the congregation, will degenerate into an aimless and useless talk; every time the preacher ascends the pulpit he must have a clear and specific aim, and he will say nothing except what serves to attain this aim. His success is proportionate to the clearness with which he formulates his aim and to the “goalwardness” of the sermon. Whether he states his aim in the form of a theme or not—he must
\quellseite{490}
have one and constantly pursue it. The same is true of the catechist. He should never present a Bible story to his children without determining exactly what he means to accomplish. But should he specify and name his aim at the very beginning of the catechesis? Certainly not the inner aim, the ultimate aim to be realized with his pupils. In that case he would anticipate the best that is to come and himself thwart the children’s attention; he would hand over to them what they should have gained by earnest efforts, and he would deprive himself of the best opportunity to engage their mental co-operation. What he will do, however, is to state the aim formally. He will not say: “Today let us learn that nothing so pleased God in Abraham as faith”, but: “Let us learn today what pleased God most in Abraham.” The former is the announcement of the inner aim, the latter, of the outward one; the former is the material announcement of the aim, the latter, the formal. The latter is necessary in order to concentrate and focus the scattered thoughts of the children, and to direct them to the point on which the whole catechesis is to turn. The expediency of stating the aim must be recognized by anyone who realizes the variety of thoughts with which the children are still occupied when the Bible History period begins and who remembers that the “narrowness of consciousness” (cf. p. 207) prevents the soul from thinking of more than one subject at a time. As long as the soul clings in thought to the old, I cannot fill it with anything new. It will behoove me first to detach it from the old and lead it in the direction of the new: this is the very purpose of the statement of aim. In proportion as the catechist succeeds in making the statement attractive, brisk, interesting, he will accomplish his purpose. Statements of this kind would be: How sin came into the world; How God made man happy; Of Joseph, Jacob’s favorite son; How Joseph came to Egypt; How a shepherd lad
\quellseite{491}
was elected king; How little David fought with the tall giant; Of David’s dearest friend Jonathan; How the lad Jesus came into the temple for the second time; Of a man who fell among robbers; Jesus takes care even of the lost; Even believers may yet perish; How the first black man became a Christian, etc. Let the student compare the headings of the Wartburg Lesson Helps (Second Course; but note that the inner, or material aim has in several instances purposely been used there as heading). In connection with the statement of aim it should also be remembered that the child, and the common man in general, prefer the use of verbs to that of nouns. “The creation of the world” does not mean so much to them as “How God created the world.” Sometimes the statement of aim of the new story is readily suggested by the preceding one. For instance, the children’s sympathy has been aroused if the teacher has, with some degree of success, told the story how Joseph was sold into slavery; they are genuinely interested in his welfare; and on their lips lies the question: What became of Joseph in Egypt? Could the teacher desire a clearer statement of aim for the following story?

\section{2. Preparation}

It is often advisable that the presentation is preceded by a preparation. In the statement of aim a word has perhaps been used which points the child to conditions in life with which he is well acquainted. The discussion of these will supply the apperceptive faculty with the fulcrum which is to bear the weight of the new material in the narrative to be treated (p. 212). Should, e. g., the aim of the narrative have been stated as: “How the first black man became a Christian,” it is advisable to remind the child of the negro whom he has seen on the street, possibly also of the fact that even negroes can be true Christians. If, in addition, a short story should be told of some negro known as an earnest Christian, the child has not only been
\quellseite{492}
provided with the material which enables him to understand the statement of aim, but a torch has been lit which shines into the distance so that the erstwhile nebulous and obscure form of the Ethiopian eunuch has become something concrete and discernible. It is also possible that the new narrative deals with objects and affairs altogether remote from the present viewpoint and concept sphere of the children, but, nevertheless, of fundamental moment for an intelligent comprehension of that particular narrative. Such matters cannot be explained incidentally during the process of presentation; for this reason they constitute material for the preparation. As a case in point, in the story of the anointing of David for the royal office the concept “anointing” occupies a prominent place. The catechist will find far fewer obstacles to a successful disclosure of the thought underlying this narrative and the development of the inner aim (namely 1 Sam. 16: 7), if, in the process of presentation and penetration, he will not be compelled to dwell upon an explanation of the custom of anointing. In the preliminary discussion he will speak of the inauguration of our presidents, of the crowning of emperors, or of the practice of anointing priests and kings in Israel and the purpose and meaning back of it. Or, the catechist, very properly, may feel justified in connecting the new narrative with the one just disposed of, because only thus the children can be put into the right frame of mind to bring the needed amount of interest to bear upon the new material. It is the province of the preparation to bring about such a frame of mind. If, for instance, the sale of Joseph into captivity has been the subject of the previous lesson, the children should be reminded anew how Joseph must have felt on his way to Egypt—leaving his father, running beside the camel in the hot sand of the desert, sobbing and moaning in his lonesome grief, telling himself that he would probably never see his home again. Now the
\quellseite{493}
children are mentally prepared to follow Joseph’s further experiences in Egypt. Thus the preparation aims to enhance the interest of the child in the new story already aroused by the statement of aim; it will remove all obstacles in the way of an effectual appropriation of the new material, and attune the cords of the soul so that the new impressions may play upon it with power and without interference. It will often be well to let the preparation issue in a re-statement of the aim. — So far as the outward form is concerned, the preparation is not necessarily purely acroamatic; it is more likely to bear the character of an informal conversation. While the explanatory words of the catechist will usually form the bulk of the preparation, there may be instances when, to make it effectual, the whole preparation will consist of questions and answers (cf. the “method of \emph{construction}” discussed below). The “preparation” may be omitted entirely when the statement of aim and the subsequent presentation are immediately intelligible; at any rate, this is better than artificially to drag it in. — Compare the practical examples appended to the fifth part of this textbook.

\section{3. Presentation}

How the new narrative is presented, is of the greatest importance for the realization of the aim of instruction or training. If the presentation fails to arouse interest, to stir the emotions, to move the will, the other steps of the teaching process are not likely to have better effects, and the instruction as a whole remains without practical results. There are three modes of presenting the new material, each one having its advocates among the ranks of the educators. The new story may be offered through reading, narration, or “construction.”

\emph{Reading} is the most ineffectual form imaginable, especially when the children are expected to read the story themselves from the textbook or the Bible. Even though they be able to read with fluency and expression, it is impossible
\quellseite{494}
for them to receive so clear an intuition as when the story is read to them in an exemplary manner; and reading the Bible narrative over once, in most cases, will have no influence upon their inner life. But does it follow from this that the teacher is to read the story to the children? It must be conceded that, if the teacher reads well, the interest of the children can be aroused. But equally correct, and even more important, is what Schueren says upon this point (p. 15): “Whatever the teacher exacts from the children, he must first exact from himself. How repugnant it is when the teacher says to the children: “You should know this, even though I myself do not.” But that is exactly what he implies, when he reads the story to the children. He should not even use reference notes, since he does not allow the children to use these when the subject is reviewed. Even more injurious than setting this example of indolence is the resulting lack of discipline and the loss of animation in teaching. Children should know themselves under constant observation, if order is to be maintained among them. When the teacher reads to the children, he is unable to watch them, unable to nip disturbance in the bud;—and disturbance is likely to be particularly frequent in just this case, for the reason that a lifeless instruction possesses no fascination for the children. If there is to be life in instruction, the subject matter must live in the teacher instead of merely being in the book. How can the teacher arouse and maintain in the children the conviction that he himself believes the material that he presents, if he resorts to reading? Many other things besides the voice should speak during instruction, notably the eye. There is no telling what effect may be produced by a soulful eye turned to the hearer. The eye of the teacher should rest upon that of the child, and the eye of the child upon that of the teacher. Often the lips will strive in vain to utter what the eyes express without
\quellseite{495}
effort.” Furthermore, what text should the teacher select for reading? That of Scripture? This is the view of many, who think that veneration for Holy Scripture does not permit any choice whatever. But the Scriptures have not been written for the children in the sense that their language and mode of presentation is always familiar to them. On the contrary, the language of the Scriptures often requires to be adjusted to the needs of children, and to be translated into their own familiar speech. In order to bring home to the children the actual sense, it will occasionally be necessary to leave out parts of the Scripture text, and, again, to supplement it by explanatory comments. In due recognition of this fact, Bible Histories have been compiled with texts frequently quite different from that of the Scriptures, so far as the arrangement of words is concerned, texts that rise from comparative freedom from the Scriptural text in the lower grades to increasing conformity to it in the upper grades (compare the three grades of the Second Course of the Wartburg Lesson Helps). This practice is warranted without doubt; and it is even necessary in view of the needed reviews and drills; nevertheless, the objections to presentation by reading stated above apply also here; moreover, the mere reading of the story from the textbook does not suffice to unfold completely before the pupils the latent powers of the text of Scripture and to bring them to bear upon the threefold life of the soul.

The reason is this: In every story, the Bible story not excepted, a fourfold factor is to be discerned: (1) the historic act as such; (2) the framework of the story—for instance, certain features pertaining to the respective fields of geography, natural history, and cultural history within which the action moves; (3) the inner world of thoughts and motives of the story, i.e., the thoughts, emotions, reflections, motives, etc., actuating the persons described; and (4) the
\quellseite{496}
religious and moral nature of such thoughts and motives, which endows the action with its real value. These four factors must be comprehended by my pupils; until that has been achieved, I have not fully presented the story nor utilized all of its dormant powers. The ordinary school texts, however, permit only the first one to come into its own; the second falls short and has to be explained afterward whereby the effect of the story is marred; the third and the fourth are merely suggested; and so far as the fourth is concerned, it often is totally ignored. And yet, the educative force of the story is largely found in the last two factors. The pupil must be trained not to stop at externals in judging of an action, but to make his way to the motive of which it is an expression, and thus to form a judgment concerning its religious and moral value (pp. 223 f., 279 f.). To this must be added that the schooltext as well as that of the Bible often indicates important situations with a single stroke. This may be, and as a matter of fact often is, decidedly characteristic and significant; but the child does not understand it in its bearings, and thus loses something of great value. For instance, what anguish and despair, what remorse and futile hope, what horror of death and silence of the grave are expressed in that brief word of Gen. 7:21. 22! All this should be brought out; the class must vividly be made conscious of it. Or consider the situation sketched by Acts 9:22: “Paul confounded the Jews that dwelt at Damascus and proved that this is the Christ”; these words are only meaningless sounds for the class, probably not even understood in their literal sense, until the teacher unfolds them and presents to his pupils a vision of Paul’s evangelistic activity with its overwhelming arguments from the sphere of prophecy and fulfillment. Facts such as these call for a mode of presentation different from a mere reading of the respective text by the pupils or by the teacher. Even
\quellseite{497}
when subsequently the penetration supplies many things neglected by the presentation, this is a poor makeshift: until that time comes, the child has to worry along with much undigested material; no overwhelming effect upon his inner life, particularly his will and emotions, has been produced. Whether the soul will later be equally susceptible to such influence, is questionable: the first impression is usually also here decisive.

For the reasons given, many catechists today, instead of reading the text or having it read, prefer the method of \emph{construction}. They understand by this the method through which the teacher, in the form of a conversation, introduces the children to the situation postulated by the story, and then, by questions and answers, prompts them to find the several features of the story and thus in a sense to construct it for themselves. Where this method obtains, the above mentioned four characteristic features of the story, which require due attention if the story is to fulfill its purpose—the last three in particular—come into their own. Notwithstanding, this treatment of the story is still inadequate. While the mental activity of the children could hardly find a better field for independent exertion, the biblical story cannot, where this method is followed, produce the designed effect upon the emotional and volitional life of the children; and when climaxes are reached in the history of revelation and God’s redemptive acts are the heart of the story, this method fails completely. How should children discover God’s purpose in regard to man either in the past or present? Even though the thesis should be given to them: “God is love”, all attempts to discover His ways from that premise would result in aimless groping, but never in a knowledge of the ways which God has actually taken. If, however, this dangerous deviation is avoided and on all important points information about God’s redemptive work is imparted to
\quellseite{498}
the children, then it is thereby tacitly admitted that a story hardly requires any other form of presentation than that of narration. While the constructive method has originated in an exaggeration of the principle of training the child for mental self-activity, it is not without its element of truth, in that the catechist, who presents the subject matter in the form of narration, will find it expedient now and then to interrupt his narration by a question, which is intended not only to make sure of the pupils’ attention, but also to enable them to carry the action forward themselves or to give an account of the thoughts and motives underlying it.

But \emph{how should the story be told?} Here, too, a three-fold mode is advocated. (a) There are those who think that it ought to be told in precisely the same manner as the pupils shall afterward find it in the school-text, with only this exception that brief explanatory remarks may be interspersed. That this view is hardly correct, is readily inferred from what has been said above; for most of the arguments set forth against the reading of the Bible story by the teacher, apply to this method of narration as well. (b) Others have advocated a method diametrically opposite, in that they demand a narration that is not only altogether independent of the Bible text but also modernized to such a degree that the very coloring of the Bible is discarded, and the Bible characters appear like present-day country or city folks. While traces of this method are found already in Paul and Zurhellen, it is Gansberg and Scharrelmann who have consistently carried it out. Not only “full detailedness and exhaustive motivation,” but also “the most child-like modernization” is Scharrelmann’s motto. He says: “When I am unable to depict ancient Jewish and Egyptian conditions and scenes, I substitute modern ones for them,” and he follows this principle even more consistently than the great medieval painters by whose methods he justifies his own. We do not
\quellseite{499}
overlook the fact that these demands contain elements of truth which we ourselves recognize and adopt (v. i.); nevertheless, we must repudiate this mode of narration or presentation; for, to us, not only the soul of the child in whose behalf the demand is made is holy, but the material as well. Much as we emphasize the necessity of setting forth the truths of Christian faith and life contained in the several Bible narratives in such a way that the child’s interest is aroused on every hand, we are at the same time aware of the Bible story’s title to reality. It is to us an educational factor for the very reason that it is a link in the chain of that history which prepares or provides our salvation; if it were a mere vestment for garbing moral thoughts, we should not use it at all as an educational factor and prefer to have recourse to tales or stories from the present.

The proper mode of presentation is (c) \emph{free narration}. In far-reaching independence of the Biblical wording such narration gives a clear, vivid, and exhaustive picture of the event; it gives due attention to the progress of the action and to the whole “frame work of the story”; and tenderly it discloses the inner world of thoughts and motives and their moral and religious bearings. At the same time, however, in spite of all independence from the Biblical wording, in spite of all detailedness and motivation, it is in point of content, faithful to Scripture and in point of form noble and chaste; it intends to set in motion the intellectual, emotional, and volitional soul-life of the child by no other means than by the facts and thoughts lying open or hidden in the Scriptural narrative. Arriving at climaxes, when important direct utterances, especially of the Lord Himself occur, it conforms even to the exact wording of Scripture.

More specifically, the story to be properly told must meet the following conditions: it must be simple, intelligible, intuitional, vivid, and it must have movement. It must be simple in regard to language and tone. Not long but short sentences; simple sentences and not complicated
\quellseite{500}
ones with intertwining clauses. Nor should the catechist cultivate an unnatural and artificial tone, but a natural one, as if a mother tells a story in the family circle. The story must be intelligible; for though it be simple in language and tone, the children may still be unable to understand it, because there is not enough adaptation to the child’s vocabulary, or because words are used which may be familiar as to sound, without, however, producing clear concepts in their minds (p. 439). As a case in point, the teacher should not merely say: “Jehovah blessed the Egyptian’s house for Joseph’s sake”, but should show in detail wherein such blessings consisted: perhaps in this way: While Joseph was overseer, the horses were never hurt; the slaves did not quarrel or steal; and everything sown upon the fields bore fruit, sixty, seventy, and a hundredfold. He must not content himself with repeating the statement of Scripture that Saul was given letters to Damascus; for even pupils of the upper grades do not know what is meant by credentials, which are here meant, and thus they will miss an important feature of the situation. “Our instruction is often so utterly devoid of joy and success because we presuppose something that does not exist” (Schueren, p. 12). How is attention to be aroused; how the interest to be awakened; how the recitation to become acceptable if so much in our narration has remained unintelligible to the children? The narration, furthermore, must be intuitional. “To use the map of Palestine in Biblical History does not make the instruction intuitional. When the child, as the story is told, sees with the eyes of his mind what is heard with the bodily ears; when the external process becomes an internal one for the child, it is then that instruction is intuitional. Brush and pigment are not the only means of painting, nor chisel and hammer of moulding: these things can be done also with words: there is such a thing as a plastic art in language” (Schueren, p. 16); comp. pp. 204 f., 256, 260 f.). To tell a story intuitionally we must pay attention to the way children and mothers tell stories: the Bible, too, contains many sections that can serve as examples. For one thing, it pays attention to the lesser features (Abraham bowed down; let now a little water be fetched and wash your feet; here I am, my son!; Joseph shaved himself and changed his raiment when he left the dungeon; observe the plastic force in Gen. 22!); there are never any omissions (when Abraham entreated God repeatedly in behalf of Sodom, it does not say: He entreated God several times; but the six entreaties are given in accurate detail, and to this day each renewed entreaty causes the child’s sympathy to mount higher); definite numbers are given (five loaves, two fishes, twelve baskets); it does not deal in collective terms, but mentions
\quellseite{501}
concrete objects (e.g., in 1 Sam. 17:17, it does not simply say food, but enumerates the individual articles constituting the food); it never, or at least very seldom (as in Acts 1:1 ff.), deals in generalizations; it does not employ indirect speech with its difficult constructions, but direct speech (not: Jesus said that she should go in peace, but: Go in peace! not, God asked where Adam was, but, God said: Adam, where art thou?). All this enhances intuitional presentation. It is true that the narrative is thus made longer; but that is of no consequence: while breadth is not always an evidence of intuitional presentation, intuitional presentation does require a certain amount of diffuseness. It dare not omit miniature painting, whereby, at the same time, the phantasy of the child is given most acceptable food (pp. 210 f., 251, 256). The outward circumstances should be delineated in detail. Explanatory comments, for instance, such as are to shed light on the state of culture at the time of the event under consideration, should always, as much as possible, be joined to the persons of the story under discussion, unless they have been disposed of in the “preparation”; only thus can they really be made to enhance the intuitional presentation and to arouse the interest of the pupils (p. 257). Under no circumstances should they be accorded recognition as independent factors. While they may be allowed to retard the progress of the action, they must neither interrupt it nor ever crowd the essential features of the story into the background. The ivy twining around the pillar and all the details of decorative art to the right and the left are very fine, but the pillar itself is of greater importance. Special care is required for the illumination of inner processes, that is, the thoughts and sentiments of the actors in the story. The catechist must mentally transplant himself into their position; must envisage their pleasures and their burdens, must visualize their emotions, their fears and hopes now rising, now subsiding, and thus trace the act to its beginnings; he must try to catch a glimpse of the inner conflicts as they end in defeat or victory, and eventually take visible shape in actions; he must take into consideration the consequences of the good as of the evil deed; only thus can he prompt the soul of the child by word, voice, and eye to vibrate in sympathy with the scenes inwardly beheld.—Fankhauser: “When we tell of the campaign undertaken by Abraham for Lot’s rescue, we may fancy how the tempter whispered to him, ‘Serves Lot right! What tricks he used to play on you! Now he has gotten what is coming to him for his selfishness! Let him see how he gets out of the trap!’ But see how Abraham
\quellseite{502}
turned down all such advice and only thought of Lot’s misery—in chains, in fear and danger of death: and he his brother’s son, commended to his care by his dying father. ‘But what will you do, Abraham? With your 318 servants you would fight the four kings with soldiers probably a hundred times as many as yours!’ ‘God goes with me, if I do what pleases Him. If He goes by my side, we shall be the stronger of the two.’ ‘But, Abraham, you have never done anything but to herd flocks, while those men are veteran warriors!’ ‘I trust in God, He will show me what to do’. Zurhellen: ‘Cain’s fratricide, as pictured in the Bible-text, is, for the children, a ghastly and incomprehensible deed. Envy is the motive; and envy should be the subject of the subsequent penetration. But in that case the narration must have succeeded in making clear the nature of envy, how it probably originated, how it gradually corroded the soul, and what terrible results it will have unless opposed by a strong will. Therewith we come upon the second pedagogic motive of great value: ‘Unto thee shall be its desire, but do thou rule over it!’ Self-control is surely one of the most fruitful subjects of religious and moral instruction; and the story before us supplies a welcome opportunity to speak of it to the children. But in that case the presentation should also contain a reference to the conflict of the good will in Cain against sin and make its ultimate defeat conceivable.’ All this is necessary if the narrative is to be intuitional. The catechist will be successful in shedding light upon the outward as well as the inward situation, especially if he observes those inconspicuous hints contained in the Bible text which so often elude the attention of the cursory reader (cf. Gen. 7:21.22; Acts 9:22; above p. 474; or Gen. 15:1 in connection with Genesis. 14). With the instinct of a hunter the catechist must be after those hints; whatever in them appeals to the intuition he must track down; focus his interest upon them, and exploit them. But these processes—the delineation of the external conditions as well as the illumination of the promptings in the inner world of the actors, can enhance the intuitional character of the narrative only when the catechist associates and compares them with such external conditions and experiences of the inner world as the children are already conversant with. This is the germ of truth in these attempts at modernizing the stories of the Bible. With a loud voice they call to the catechist: “Do not forget the connecting threads between the home of the child and the home of your Bible stories; do not forget the fulcrums for the apperceptive faculty (p. 212) in the sphere of the children’s
\quellseite{503}
intellect and experience; look for them; turn them to account if your narrative is to take life in the soul of your child and quicken it to an inner experience of the same.”—If we name, in addition to intuitiveness, also \emph{vividness} as an essential attribute of an effective narration, we have in mind not only a noble, animated delivery, nor merely Fankhauser’s definition: “Vividly to narrate means so to narrate that the story lives in us and we in the story,” but also what Schueren emphasizes when he says that the catechist must have experienced the vital power of the story in his own heart, and that his narrative must become a vivid testimony of the divine life within him.—Finally, narration must have \emph{movement}, it must aim for the goal, it must be \emph{zielstrebend}. This means that it must not tarry too long here or there, least of all with external features; there must be no coloring for the sake of the color; but everything must be viewed in reference to the final goal to be attained, i.e., the precise religious or moral truth in the story which is to be set forth; and everything must be left out, no matter how beautiful, which does not serve the end in view—the realization of that truth.

If the story is told in the manner here set forth, a quickening of the emotions through the intellect will be the result; the story will appear precious to the soul; it will become a factor of its inner experience. The result will be an “ideal association” (p. 260) with the characters of the Bible, a holy \emph{συμπάθεια} (cf. pp. 226 f.). The children think, feel, grieve, pray with them; they plan, strive, suffer, believe, hope with them; they sing, rejoice, shout with them. What a power for culture, education, discipline that means, is known to every pedagogue (pp. 252, 254)!

Last of all, the question remains to be answered whether a story is to be told as a whole or in sections. The latter method still has the most advocates in the present, and the author has not hesitated to give an example of a story told in installments (see practical examples). In case a story is of undue length and makes two or three sharp turns it may be convenient to tell it in sections, and at the end of each section to reduce it to the words of the school text, to have the children reproduce it in their own words, and to let them find the proper heading for the section before proceeding to the next section. But when the story is a
\quellseite{504}
natural catechetical unit and not too long, the effect upon the emotions and the will is not subserved at all when this method is followed; we are convinced that the uninterrupted narrative is every time much surer of a telling effect; and this is decisive for the catechist. In the Biblical History by the author, an outline of every story is given at the conclusion; but this is done altogether in the interest of the pupil, to aid him in reviewing it and committing it to memory. When the whole story has been exhaustively told, the pupils should be directed to open their Biblical History and read the story or recite it extemporaneously, if the latter can be done without torture.—Practical examples are found at the conclusion of the section.

\section{4. Penetration}

With presentation, if effectual, the most important feature in the treatment of a Biblical History has been disposed of. A picture, clear and intuitional, of the story, has been painted for the children from both aspects, external and internal; their soul has been warmed and their will has been set in motion. But the process is as yet far from complete: unless other features are added its effect will be lost too soon and hardly ever cause more than a fleeting mood or sentiment. Instruction and training, however, dare not rest content with the production of religious and moral moods. The object to be attained is: Clear convictions, moral and religious value judgments, a consciousness of the Christian life as a real life—as the receiving of God's grace, as a struggle against sin, as a life of work according to God's will. That this effect, produced by the story upon the child's soul, be perpetuated, clarified, confirmed, deepened, and made fruitful unto life, the children must be directed once more, and that at leisure, to meet the characters of the story face to face; carefully to scrutinize their thoughts and motives—the whole complex of the mental and moral forces impelling them to action; to examine them from the
\quellseite{505}
standpoint of their religious and moral character; to trace carefully the consequences of wicked and of godly conduct; to make the relation of God to man, and of man to God a subject of meditation, so that faith, love, fidelity, unselfishness, etc., may in their majesty confront and captivate them; and that unbelief, unfaithfulness, hatred, envy, selfishness, lying, etc., may stand revealed in their disastrous and abhorrent hideousness. This is the purpose of penetration. It is to gather in leisure all the illuminating and cheering, but also all the repellent and dreadful, rays of the story; to let them exert their inherent force with power, and thus to turn to account the educative elements of the story for the training of the young to the full extent possible. Thus in penetration the last two constitutive elements of each story rather than the first two (p. 473) must be emphasized. For here the catechist comes to speak of the outward progress of the story and its external frame-work only in order that the catechumen may gaze upon the inner world of thoughts and motives, recognize moral and religious values, observe their vital importance, and gain respect and desire for them. When this is understood, it must be clear that penetration deals with persons, especially the leading persons of the story, in accordance with the fact that formative, educative force is mostly found in the personal, biographical element (p. 256). Of whom have we heard in this story? Or, whom do you like best in this story?—these are the questions most naturally asked when the penetration begins. When light is shed upon the thoughts and acts of those persons, the teacher should see to it that the children apply no other standard of judgment than that of God. While sinful acts are by no means excluded from illuminating discussion, care must be taken not to go unduly into detail in the delineation of sin, lest destructive germs are deposited in the children’s hearts. If such guidance to religious and moral value judgments is given by
\quellseite{506}
one who himself loves God and His will, it is by no means, as some maintain, the critical spirit, with its desire to judge others, which is awakened, but rather the children's conscience; that which is good in itself has been made to appear good and precious also to them; and what is evil in itself has been made to appear detestable. While shedding light upon men and their doings, the catechist should never forget shedding light upon God and His doings. His attributes and nature should be brought out in clearest light at the hand of his actions. This is the stage of instruction in which the catechist, in conformity to the example of Christ (John 17:6), is to reveal to his pupils the name of God; that is, he is to withdraw the veil, so that, behind His acts, they see God Himself in His divine nature and recognize Him as absolute personality and holy Love. (Comp. M. Reu, \emph{Unsere Erziehungsarbeit im Lichte von Joh. 17, 6}, Kirchliche Zeitschrift, 1911, pp. 453–466). It is often desirable that the most important points brought out in penetration, for the purpose of impressing them indelibly upon the children's minds, should be pronounced in unison by them and written upon the black-board in the order of their development. When this is done, an easier survey is attained, the final summing up is made easier, and the service of the eye is enlisted in order to secure a more lasting effect for the educative elements found in the story (comp. the complete practical example at the end of this fifth section).

If advisable from the didactic and pedagogic stand-point, there may be joined to penetration what has been called by the disciples of Herbart “\emph{association}” or “\emph{comparison}”; but only in that case. For just the constant—and hence sometimes unavoidably artificial—employment of association and comparison, has, perhaps not without reason, largely occasioned the discreditation of the formal step theory. What is to be done is that the children are made to recall
\quellseite{507}
for themselves the other stories in which the same features in the acts and nature of God, in the doings, experiences, reflections of men, in the relation of cause and effect, have aroused their attention. Whether such stories are taken from Biblical History, Church History, the Reader, or, on occasion, even from other branches; or whether they are taken from the children's own experience, is a matter of no consequence. Such stories being found, the truth developed from the new story is confirmed and enhanced in power; its universal character is recognized by the child, and thus it becomes in his eyes an incontestable certainty, and a factor ever to be reckoned with. It is hardly necessary to quote examples; may it suffice to call attention to the importance of trust in God as shown in the life of Abraham, of Moses, of David; to the effectual prayers of Jacob, of Moses, of Solomon, of Elijah, etc. While such comparison is not always required, there is another feature of discussion, called \emph{generalization}, which is indispensable. Invariably the main truth developed from the story should be summarized in the form of an easily remembered, classical sentence, which may be borrowed from Scripture, the Catechism, or the hymnodic treasure of the Church. Here everything is gathered as in a focus. Where it was possible previously successfully to employ the step of comparison, it is natural that the truth of this summarizing sentence is impressed even more deeply upon the minds of the children. It becomes the focus in which the truth not of one but of many stories is gathered. That such truth dare not be mechanically grafted upon the presented story, but should vitally grow forth from within, follows from what has been said. If presentation always should point “goalward”, penetration should do so even more; every constituent part of penetration should lead in a straight line toward the passage in which ultimately everything is gathered together, so that this summing up is not
\quellseite{508}
only the end but the consummation of penetration. While, in penetration, occasions arrive for the introduction of illuminating Bible verses, stanzas of hymns, or parts of the Catechism, these must yield precedence to the golden text, the ultimate goal. Since the work of memorizing must be done largely in school, especially in the lower grades, this golden text must finally be drilled by recitation in unison.

As to the \emph{outer form} assumed by penetration, not much remains to be said. It cannot be any other than a natural didactic conversation. While the didactic question has already been an occasional feature of the descriptive and detailed narration, the principle of independent activity requires that penetration should proceed entirely in the form of questions and answers, without, of course, excluding the occasional employment of the acroamatic form. The practical example based on the story of Zachaeus (at the conclusion of section V) should be converted into the form of a didactic conversation.

\section{5. Application}

The necessity for an application to follow penetration and, likewise, the aim of the same, has been demonstrated in ch. 32. But what character is the application to bear? At the very outset, it is necessary to caution the catechist against the fallacy that numerous applications are required. In presentation the catechist keeps in view one single aim; and during the progress of penetration he aims at only one summarizing sentence although he may employ a number of Scripture passages or hymn stanzas; in the same manner the catechist should not in the end commit the error of making many applications: the effect of that would be that none has made an impression deep enough to become a fruitful factor in the daily life of the children. What has been said of the preacher, applies to the catechist as well: “There are preachers that take too much for granted; they take a chest full of truths along with them upon the
\quellseite{509}
pulpit, just as someone might take a chest full of nails; and then, on the supposition that the hearers are posts, they take out a nail and expect it somehow to enter the post of itself. But that is not the right way. They must take the nail, hold it to the post, hammer it in, and then clinch it on the other side. Not until then have they the right to expect that the great Master will so fasten the nail that it will not fall out again.” This is true in an even higher degree of the catechist; for children are less able than adults to take home with them a mass of loosely applied truths and to make them factors in their daily life. And just as but one application is to be made, so this one application should not be a general one but one very definite, clearly bounded on every side; for any application necessarily loses force in the degree that it is general. E.g., when the story of Abraham’s departure for Canaan is to find its application, it should not be formulated like this: “Children, you, too, must believe in God,” but: “though God asks difficult things of us, we still should obey Him, confident that He will be with us and bless us.” Finally, the application should be something more than a mere doctrine. If the latter were the desideratum, the application would mark no progress; for the doctrine taught by the story has already been found by penetration and has been gathered into the golden text. What the catechist has to do, is to point out definite incidents in the life of the children and show, in contrast to their real conduct, what it ought to have been. To this end he will go with them into the past or future of their own life, in order to train the conscience for a recognition of sin, and to awaken in them wholesome purposes and resolutions. Moreover, it is their life, not that of others, to which application will have to be made; for in their present life, as children, they have to begin to adapt themselves in all seriousness to the religious and moral truths of Scripture. Instruction
\quellseite{510}
is to be more than a mere aggregate of directions to the children regarding their conduct in the future. Also as adults they will deem it a matter of course in all things to consult the will of God, if, as children, in their relation to parents and teachers, to brothers and sisters, to the neighbor and the neighbor’s children, in the house and in the garden, while at work and while at play, they were accustomed to conform to that will. The skill of the teacher will have its mettle tested in the attempt to restate the great deeds of the Bible characters in the terms of the children’s humble life. He who deems himself too great for that or even looks upon it as a triviality, may make many an application quite good in itself; but the souls of the children he will not touch: and it is his fault when the children fail to locate the channels into which the vital forces quickened in them are to issue; when the impulses awakened in their hearts through presentation and the penetration evaporate without fruitage, and thus already in youth that cleavage is formed between knowledge and life, worship and every-day life, which afterward, in adult life, proves so difficult of removal. Not seldom a short, well-chosen story germane in every way, is best suited to open the eyes of the children for the duty suggested by the biblical narrative with which they have dealt. Of course, the story intended to clinch the application should not be taken from the Bible; it should be one from life, from the present, from Missions; and it should be a juvenile story though this is not an invariable rule; but it must be only one, not a number of stories.

The catechist will find much material useful for presentation in H. Witt, \emph{Die Biblischen Geschichten Alten und Neuen Testaments mit Bibelwort u. freier Zwischenrede anschaulich dargestellt}, 3 vols., ²1883; G. Staebler, \emph{Jesus, der Schoenste unter den Menschenkindern}, Stuttgart, 1905 (however, everything requires re-statement in terms familiar to children); M. R. Unger, \emph{Beitraege zur erbaulichen Behandlung}
\quellseite{511}
\emph{des Religionsunterrichts auf der Unterstufe}, 1912; C. Stuckert, \emph{Jesusgeschichten}, 1910–14; especially in G. Fankhauser, \emph{Der Weg zum Kinde}, 1915 ff., and M. Reu, \emph{How I Tell the Bible Stories to My Children}, ²1926. M. Affolter, \emph{Kommt, Kinder, hoert}, ²1926. \emph{For penetration and application} in J. Nissen, \emph{\emph{Unterredungen ueber die biblischen Geschichten}}, 1888; E. Thraendorf and H. Meltzer, \emph{\emph{Der Religionsunterricht}}, 5 vols., ²1905 ff.; A. Reukauf and E. Heyn, \emph{\emph{Ev. Religionsunterricht}}, vol. iii to ix, ²1906 ff. H. Spanuth, \emph{\emph{Praeparationen fuer den ev. Religionsunterricht}}, 4 parts, 1909 ff.; but preeminently in G. Friedrich (Rienecker), \emph{\emph{Der Herr bleibt ein Koenig in Ewigkeit. Unterredungen ueber die bibl. Gesch. d. N. Testaments}}, 1911, and G. Fr. Rienecker, \emph{\emph{Gottes Wort bleibt in Ewigkeit. Unterredungen ueber Gesch. des A. Test.}}, 1912. \emph{For application} in G. Staebler, \emph{\emph{Sammlung von Erzaehlungen und Grundgedanken f. d. Unterricht in der bibl. Gesch.}}, 1902 f.; J. Heininger, \emph{\emph{Anwendungen bibl. Geschichten}}, 1900; E. Zeller, \emph{\emph{Andachten fuer Kinder}}, ²1914; some material also in G. v. Viebahn, \emph{\emph{Von der Landstrasse des Lebens}}, 1908; and in the \emph{King’s Highway Series}, edited by H. Sneath, G. Hodges, and H. H. Tweedy, 1917.

As far as the \emph{several grades} are concerned in which the Bible stories are to be treated, but little remains to be said. In the lower grades, of course, the whole process of instruction will be simpler and more child-like. But even here the steps named above can be retained; penetration and generalization may be successfully employed in the lower grade, as has been clearly shown by Fankhauser and Unger. It is quite true, however, that it will at times be sufficient in the lower grades to narrate the story in an intuitional manner, and to make the narration issue in the golden text which sums up its main truth. In the upper grade penetration and application are to be cultivated with particular care. If the same stories recur in all three Courses, as is often the case according to the curriculum for parish schools, given in Chapter 31, it need not invariably be the same truth which is placed in the foreground. The individual stories are at times so rich in content that they suggest a variety of inner aims; and often the mental development of the children does
\quellseite{512}
not permit the specification of the ultimate intent of the story until the Third or Upper Course has been reached. The story of the stilling of the tempest may be treated in such a manner that only this truth becomes obvious: “The Saviour saves His own when danger overcomes them”; but also in such a manner that it teaches: “If God be for us, who can be against us?” The pupils in the First or Lower Course understand the former; the second is required for the pupils of the Third or Upper Course. The story of Zachaeus can be treated from the point of view: “Even the publican is not too great a sinner for the Saviour; He enters his house, for the Son of Man is come to seek and to save that which is lost”; but also from this: “If anyone is in Christ, he is a new creature”; and even from the third viewpoint: “To receive forgiveness of sins from Jesus and to right the old wrong, is the way to true happiness.” The story of the marriage at Cana may serve to impress the children that the Saviour is almighty, for He can turn water into wine; but also that “Jesus glories in being a powerful helper even in the small troubles of daily life.” The consideration of the passion of Christ may stop in the First Course at the bodily anguish; in the Intermediate Course stress may be laid upon the inner anguish occasioned by the disloyalty of the disciples and the people's unbelief; while the most difficult feature may be reserved for the Upper Course: the enduring of the divine wrath. The story of Abraham offering Isaac may be made to show that Abraham loved God more than his son; but also that Abraham's faith was staunch and equal to the very greatest trial; and finally, Abraham and Isaac may be held up as types of Christ and God the Father.

“In Quarta the pupils are shown, in the light of Saul's conversion, “how God defeats and hinders every evil counsel and purpose which would not let His kingdom come.” That class, according to the curriculum, studies this story after having learned of the conflict of the Sanhedrin with the Christians, the martyr death of
\quellseite{513}
Stephen, and the dispersion of the Christians. The assurance that Christ is risen from the dead, that He “lives and reigns” is nothing new for them. This assurance, based on what had been taught in Quinta is now merely confirmed. In Tertia where the Life and Teaching of Paul is treated, it is a different truth that should be brought out. There the theme of the story (outer aim) is not: ‘How is the most furious persecutor of the Church conquered by God?’ but: ‘How did Paul come to faith in Jesus Christ?’ And the answer is: ‘God revealed in me His Son’ (Gal. 1:16). ‘He was laid hold on by Christ’ (Phil. 3:12). ‘By the grace of God I am what I am’ (1 Cor. 15:10). ‘By grace are ye saved’ (Eph. 2:8): this is now the saving truth drawn from the very same story. In Quarta the time is too short to bring out this phase, and the pupils would very likely fail to fathom it. And still different will be the statement of aim and the generalization of the saving truth in Upper Secunda. There the subject under consideration is the apostolate of Paul, the divine origin of his Gospel, the history of the conflict between a free Christian faith and the Christian faith as perverted by Pharisaism, or of the transition from Judaistic Christianity to Gentile Christianity; and Bible verses like 2 Cor. 4:6, 1 Cor. 2:9 f. (in connection with 2 Pet. 1:16–21; John 1:14) now summarize in classic words the saving truth of the self-same story.” While these words of Zange (p. 186) presuppose other educational conditions than those prevailing among us, they constitute a notable testimony to the cause to which we stand committed. Let this be noted by our American sticklers for method, who deem it unpsychological and unmethodical to permit the same story to recur in a “Graded Series”, and who take recourse to all kinds of material to avoid this. Zange, a man of pre-eminent training, both as psychologist and as educator, goes so far as to add: “Another proof how necessary it is to return to the same material in the successive grades.” Cf. p. 425.

In the Sunday School, the narration of the story can naturally move toward one aim only because it is intended for all the children in the Second Course (cf. p. 420); but the higher aims can at least be suggested to the pupil by the “Questions for Review” in the Lesson Helps for the Intermediate and Junior Departments when the teacher of the class is not mature enough to lead up to them himself.
\quellseite{514}
The method of teaching here recommended requires much time: in most instances one hour will not suffice for all the successive steps. Nor is it necessary to mount them all in one hour. In the \emph{parish school} the work of the day may close with the presentation and the subsequent reading of the story from the textbook on Biblical History. The pupils are expected to read the story several times at home; and in the next lesson penetration and application take the place of the mechanical questioning which is still widely in vogue. In the \emph{Sunday School} the story is told to the whole Second Course or Department in a detailed manner, as explained above. Because such narration is designed to arouse the feelings of the children and to move their will, it should \emph{never be omitted}; the preparation may be left out (but not the statement of aim) for lack of time. After the story has been told, the several grades may read it under the supervision of the teachers from the graded Lesson Helps, and the Bible verses will be drilled by recitation in unison. On the following Sunday, after the Superintendent has presented and the class has read the new story, the teacher of the individual grade reviews the story of the preceding Sunday with the aid of the review questions appended to every lesson of the \emph{Wartburg Lesson Helps}. Penetration and application of the story taught in the Sunday School may occur in the Saturday School (cf. p. 429). The same method recommended above for a regular parish school can without difficulty be observed in the partial parish school, comprising the children of 12 and 13 years. It is true, not every story can here be thoroughly studied according to our method although a lesson in Bible History is assigned for every forenoon; but that is no great loss when one has done his duty in the Sunday School and Saturday School. Some of the stories are taken up merely by way of review, leaving for thorough treatment only the more important lessons upon
\quellseite{515}
the Sunday School calendar, and those which in Sunday School had been omitted entirely.

When the end of a group of stories has been reached, a review may be undertaken: a detailed sketch of the outstanding character may be drawn, or the principal characteristics of God and the important persons may be assembled, always with a view to the training of souls (p. 483). A sketch germane to the points here made is found among the practical examples at the end of this section.

At times \emph{Bible reading} may be connected with the treatment of the Bible story. In connection with important phases of the discussion, especially when the steps of comparison and generalization have been reached, the upper grade may read—directly from the Bible instead of the textbook—a Psalm or another section that comports well with the truth developed (cf. p. 371; for instance, in connection with the story of creation, Ps. 104 and 8; the passage through the Red Sea, the song of Moses and Miriam; David’s shepherd life, Ps. 23; the story of his fall and subsequent repentance, Ps. 51 and 32; the destruction of Jerusalem, Lamentations; at the turning points in the history of the kings, passages from the prophets; the exile, sections from Isaiah 40–66; the return of Israel from Babylon, Ps. 126; the episode of Cæsarea-Philippi (Matt. 16), Isa. 49; at the foot of the cross, Isa. 53; in connection with the ascension, Ps. 110; on Pentecost, the pertinent section from Joel, etc.).—That the most important facts of \emph{\emph{Biblical Geography, Natural History}}, etc., should be associated in part with the preparation, in part with presentation and penetration, follows readily from the foregoing. There are places where such facts should be brought out with telling effect. When, e.g., the journey through the wilderness is the subject, let the needed geographical information about the peninsula of Sinai be given; when the class enters Canaan with Israel, an accurate survey of the
\quellseite{516}
country should be afforded; when it is taught that God gave Israel every blessing promised, an exhaustive picture of Israel’s life and happiness should be painted (cp. Practical Examples); later the class should follow Israel to Babylon and back; the Christmas story with its mention of Nazareth, Bethlehem, and the land of Judæa furnishes the occasion for a discussion of the division of the Holy Land in that period of its history, which should be illuminated by constant reference to the map; the mention of Emperor Augustus suggests the same treatment of the Roman Empire; when the pupils travel with the twelve-year-old Jesus lad, the journey should be traced upon the map; Paul’s missionary journeys can not possibly be taught without that adjunct. Places like Jerusalem, Bethlehem, Nazareth, and Capernaum lend themselves to geographical illumination (cp. Practical Examples).

Finally, concerning the use of \emph{Biblical pictures} in connection with the several stories, the following remarks are in place. That the use of biblical pictures deserves encouragement as a stimulus for the intuitive faculty, goes without saying; however erroneous the opinion is that the mere use of pictures constitutes intuitive instruction. There are catechists who use pictures diligently, and yet fail to impart intuitional instruction; for the inward intuition, as above described (pp. 478 ff.), is more important than the employment of the finest pictures. This much is true, however, that the internal intuition is enhanced and prolonged when the external intuition is facilitated by a really good picture. But when is the picture to be exhibited and put to use? Before the presentation, during the presentation, or in connection with the application? Each of these three views has found its advocates. The exhibition of the picture before the presentation presupposes that the picture shall be made the starting-point for the whole process of instruction. For the so-called intuitional method of instruction in
\quellseite{517}
the home or lowest grade of the Sunday school that is the proper method; for there the inward intuition is still secondary (compare \emph{Wartburg Lesson Helps}, first course; Wangemann in his \emph{Zwanzig Anschauungsbilder fuer den ersten Unterricht in der Biblischen Geschichte} and in his \emph{Anweisung zum Gebrauch der 20 Anschauungsbilder} goes too much into detail). But in the same degree in which the mental development of the children progresses, facilitating as well as requiring the inward intuition, this method loses its usefulness. While the detailed narration is in progress, the children are still so occupied with the picture that they are not open to the inward intuition intended by the narration; and even the teacher may tend to confine himself to the cultivation of a purely outward intuition.—He who puts the picture to use in connection with the stage of application, is actuated by notions far removed from the purpose of biblical instruction and based upon the presumption that the pupils are possessed of a degree of mental maturity that is ordinarily rare even in the Upper Course of our parish schools. The children are here expected to express an opinion upon the picture and its most salient features, upon the ideas conveyed by the artist upon canvas; this tends to place the picture in the service of esthetic culture.—When we consider that the picture, if it is at all valuable from the catechetical standpoint, represents the climax of the action, it is not likely that any other use of it can be justified than its exhibition during the illuminating narration; that is, just when the climax has been reached and the inward intuition has already taken place. In that case the picture does not divert the attention of the pupils but rather gives strength and depth to the inward intuition, which is now aided by the eye. Having reached the climax, the narrator can properly make a pause and give a little time to a discussion of the principal features of the picture. If the latter crystallizes as it should, the thoughts of the
\quellseite{518}
characters of the story, there will be occasion to recur to it during the penetration.

Among the series of pictures which can be put in the service of the schools, those by Julius Schnorr von Carolsfeld still deserve first place, although this pertains rather to those on the Old Testament. Heilmann, as late as 1911, says: “No one else could bring so much artistic power and religious understanding to bear upon the sacred material as he. In regard to magnificence of composition and passionate interpretation these Bible pictures have hitherto never been surpassed. . . . To illustrate the story of creation, of the patriarchs and prophets, no grander illustrations are to be found. A valuable feature of pedagogic and religious moment is their immaculate moral character and sublime religiousness.” Another merit is that they usually represent the climax of the action. For the New Testament the series by Hoffmann and Schramm excels. There are also good features in the collection of Morgan and Copping; the series by Karl Schmauk (equally realistic, but superior as regards execution of detail), and in the collection made by the Gesellschaft für christliche Kunst in Munich. Recently the Catholic Fugel has produced a good series (twelve Old Testament and as many New Testament pictures have been published to date). Among the pictures illustrative of the parables, those by E. Burnand have justly won renown. Compare the instructive monograph of the Catholic Alfonse Heilmann, \emph{Bibel-Bilder. Gedanken zur religionspädagogischen Wertung biblischer Kunst}, Kempten, 1911.

\quellseite{519}
\chapter[§ 35. THE METHOD IN THE CATECHISM]{§ 35. THE METHOD OF INSTRUCTION IN THE CATECHISM}

Compare the literature in ch. 14, 26, and 32. *G. v. Zezschwitz, ii 2: \emph{Die Katechese od. die kirchliche Unterrichtsmethode}, 1859. *F. Lange, pp. 167–241. *A. Eckert, pp. 115–163. J. Gottschick, pp. 179–186. J. Berndt, pp. 99–105. R. Kabisch, pp. 246–247. J. Steinbeck, pp. 201–218. M. Reu, \emph{Quellen z. Gesch. d. kirchl. Unterrichts}, 1904 ff. O. Schoenhuth, \emph{Methodenlehre}, 1903. *B. B. Hahn, \emph{Luther's Katechismus erklärt aus bibl. Theologie}, 1891. A. Scheller, \emph{Die pädagogische Bedeutung d. luth. I. Hauptstueckes}, 1905. F. Niebergall, \emph{Die paulinische Erlösungslehre i. Konfirmandenunterricht}, 1908. K. Eger, \emph{Taufe und Abendmahl im kirchl. Unterricht der Gegenwart}, 1911. H. Matthes, \emph{Der 2. Artikel auf bibl. geschichtl. Grundlage}, 1913. D. Vorwerk, \emph{Gebet und Gebetserziehung}, 1913. PRACTICAL HELPS: J. Nissen, \emph{Unterredungen ueber den kl. Kat. Luthers}, 1885. F. W. Schuetze, \emph{Entwürfe und Katechesen zu Luthers kl. Katechismus}, 3 parts, 1888. K. Ender, \emph{Handbuch zum kl. Katechismus L's fuer Lehrer und Pfarrer}, 1874. L. Wangemann, \emph{Einführung i. d. Verständnis d. Luth. Katech.}, 3 parts, 1880 ff. H. Cremer, \emph{Unterweisung i. Christentum nach L's kl. Katech.}, 1883. O. Zuck, \emph{Katechesen ueber L's Katech.}, 1883. G. v. Zezschwitz, \emph{Christenlehre}, ²1883 ff. J. Chr. G. Schumann, \emph{Handbuch d. Katechismusunterrichts}, 3 parts, 1894 ff. F. H. Kahle, \emph{Der kl. Kat. L's, anschaulich, kurz u. einfach erklärt}, 1896. K. Boehruecker, \emph{Der Katechismusunterricht}, 1895. A. Nebe, \emph{Der kl. Kat. ausgelegt aus Luthers katech. Werken}, 1891. Fr. Oehme, \emph{Die 5 Hauptstuecke d. luth. Katech. katechetisch bearb.}, 1891. Ziethe, Rohden, Heydt, \emph{Die unterrichtl. Behandlg. d. 6. Gebots in der Schule}, ²1894. A. Gruelich, \emph{Skizzen z. unterrichtl. Behandlung d. kl. Katech. L's}, 1899. Th. Hardeland, \emph{Die katech. Behandlung d. kl. Katech. L's in Unterredungen}, 1899. *J. H. A. Fricke, Handbuch d. Katechismusunterrichts, 3 parts, ²1899. *S. Bang, \emph{Katechetische Bausteine zu christozentrischer Behandlung des 1. Hauptstueckes}, ²1901. G. Heimerdinger, \emph{Praeparationen f. d. Katechismusunterricht}, 1901. W. Beyschlag, \emph{Christenlehre auf Grund d. kl. Kat.}, ²1903. *R. Stange, Der Katechismusunterricht, 3 parts, ²1903–08. M. v. Nathusius, \emph{Die christliche Lehre nach L's kl. Kat.}, 1904. G. Kaelker, \emph{Der Katechismusstoff f. d. 5–8. Schuljahr}, 1904. A. Eckert, \emph{Der Katechismusstoff d. Konfirmandenunterrichts in synthetischem Gang}
\quellseite{520}
\emph{u. psychologischer Stoffordnung}, 1905. A. u. F. Falcke, \emph{Praeparationen z. Unterricht i. L’s kl. Katech.}, ⁵1906. *R. Steinmetz, \emph{Katechismusgedanken; Beitrag z. katechet. Behandlung d. 5 Hauptstuecke i. Kirche u. Schule}, 3 parts, 1906–13. Hahne, \emph{Praeparationen f. d. Katechismusunterricht}, ²1909. *A. Eckert, \emph{Kinderkatechismus f. Schule u. Kirche, I. u. 2. Hauptstueck}, 1909. *J. Kolbe, \emph{Der kl. Katech. L’s in ausgefuehrten Katechesen}, ⁷1911. *K. Eger, \emph{Katechismustafeln}, 1911. *O. Hardeland, ⁵² \emph{Konfirmandenstunden}, ⁵1914. F. Heyn, \emph{Katechismusunterricht in Reukauf u. Heyn, Ev. Religionsunterricht}, ²1913. *K. Eger, \emph{Evangelische \emph{Jugendlehre}; Hilfsbuch z. religioesen Jugendunterweisung nach Luthers Katech.}, ²1912. *Th. Kaftan, \emph{Auslegung d. luth. Kat.}, ⁶1913. *H. J. Schuh, \emph{Catechizations on Luther’s Small Catechism}, 1914. Geo. W. Lose, \emph{Catechism Bible Narratives. A Series of Bible Narratives on the Five Chief Parts of Luther’s Small Catechism}, 1915. ON APPLICATION: *K. Caspari, \emph{Geistliches und Weltliches}, ²³1915. P. v. Zychlinski, \emph{Illustrirende Aussprueche, Sentenzen und Geschichten z. Gottes Wort}, 1900. L. Pestalozzi, \emph{Die christl. Lehre i. Beispielen}, ³1901. Baum, \emph{Erzaehlungen und Beispiele z. kl. Katech. L’s}, ²1907. *J. Besch, \emph{Aus der Lernstube des Lebens}, ²1913. Leipziger Lehrerververein, \emph{Im Strom des Lebens. Narratives on the Catechism} (Columbus, Ohio). *G. Warneck, \emph{Die Mission i. d. Schule}, ¹⁴1912. Th. Schaefer, \emph{Die Innere Mission i. d. Schule}, ⁷1912. O. Koenig, \emph{Die Mission i. Katechismusunterricht}, 1913. F. W. Foerster, Jugendlehre, ⁶⁵1915. Compare also \emph{The King’s Highway Series} by Sneath, Hodges, and Tweedy, and \emph{The Golden Rule Series} by Sneath, Hodges, and Stevens.

True as it is that instruction in Biblical History, if it be of the right sort, is at the same time an introduction into the truths of the Catechism, separate instruction in the latter remains necessary notwithstanding. There exists no other way of giving the young a concise and comprehensive introduction to the faith as it lives and manifests itself in the mature congregation (pp. 310 ff.). The method of such separate instruction in the Catechism is the next subject claiming our attention.

If the aim of religious instruction as a whole has been rightly described on page 280, and if instruction in the Catechism is a part of religious instruction in general, it must contribute its part to the attainment of that aim. Every
\quellseite{521}
method, therefore, is erroneous by which those vital sacred truths fail to find that firm and deep anchorage which is needed for the stimulation of the emotions and of the will. Now, the method in vogue in \emph{Reformation times} does not measure up to the standards which we have set. It is true that all the elements necessary to an efficient method were on hand. Luther had inserted biblical pictures in his Catechism, and most catechists followed his example in their editions of Luther’s Catechism. In the preface Luther had advised teachers to adduce many illustrations from Scripture (p. 114); also in this regard there is no lack of followers, among whom Loener, 1544, and the authors of the Joachimsthal Catechism, 1574, are prominent. The practise prevailed of connecting the Catechism with the Bible and of adducing proof passages from the Bible; for instance, Loener, 1544, Huberinus, 1544, Jos. Opitz, 1583. The hymn was linked to instruction in the Catechism, not only in so far as, in the catechetical period on Sunday, hymns were sung that comported with whatever part of the Catechism happened to be under discussion (a very common practice), but also by adducing illuminating hymn stanzas in connection with the several parts of the Catechism (e.g., G. Walther, 1581). Victorius, in 1591, even published a catechism in which the individual sections were accompanied by a biblical picture, the text of a pertinent Bible story, printed in full, and a number of Bible passages.

This material, however, was not turned to account for the development of an efficient method. It was deemed sufficient that the father should recite the several parts of the Catechism before the family, and the teacher before his class, whereupon they were repeated until retained by the memory. When catechisms with explanations appeared, the explanations were taught in the same way. Camerarius and Lossius improved the method slightly; but they found few
\quellseite{522}
followers. The former would explain the Catechism to the upper classes in coherent discourses in which illustrative stories (taken chiefly from ancient history or invented by himself) were interwoven with the catechetical material, and then address a few summarizing questions to the pupils (cf. Reu, Quellen, i 1, 2 p. 98); the latter often introduced Objectiones, i.e., misgivings to be cleared away by instruction. We must indeed admire the exemplary zeal with which the Catechism was “pounded” into the young and common people, and the unlimited faith in the reconstructive and regenerative power of the Word of God; we must also recognize that some of the teachers succeeded in keeping their catechetical explanations free from specifically dogmatical elements (although the number of such teachers was steadily diminishing). The manner, however, in which the Catechism was offered to children was anything but exemplary. If, notwithstanding, the German people experienced a new birth and if a generation grew up which could endure the terrors of the Thirty Years’ War, such results, instead of demonstrating the correctness of the catechetical method in vogue, should rather be considered as evidence that God may bless the efforts of men who are faithful even though they may use imperfect methods. Moreover, the pronouncing and repeating method was much more expedient in an age when church membership was a matter of course and the consciousness of church authority lived as a power in the hearts of the people to a degree impossible today.

The Magdeburg pastor and educator Andrew Cramer, living toward the end of the sixteenth and beginning of the seventeenth century, is possibly no unbiased witness, but he is probably right in regard to the main point when, in his \emph{Anleitung, wie die Jugend in Gottesfurcht, Kuensten und Sprachen zu unterweisen}, he says concerning religious instruction as practiced in his time: “The Catechism is indeed in use in the schools, but it is taught in a very inefficient manner. A lesson is assigned to the boys, which they have to memorize and cram into their heads. While pronouncing the Catechism to the boys and their repeating it is good enough as far as it goes, it does not go far
\quellseite{523}
enough. No explanation is offered in the form of sermons adapted to children or by means of Bible passages and illustrations; and they are not given an adequate spiritual understanding of the Catechism. Nor is any application to life and Christian practise in evidence. When the boys are about to enter the adult stage, the Catechism is treated in the same manner by means of a foreign tongue: with glosses in a foreign tongue the explanation is made; but all to no purpose, ‘\emph{memoria coacta, sine intellectu, sine praxi}’, Christianity, therefore, has not been planted in the right fashion!”—An element of progress is found in the method of \emph{John Valentine Andreae} (pp. 125, 127). He had his pupils first memorize the text of the Catechism, or of his \emph{Evangelische Kinderlehre aus heiliger, goettlicher Schrift} (1621). That done, however, he would enable them, by a freer treatment, frequent reviews, clear generalizations, sound explanations and skillful application, to reproduce intelligently the words memorized and “to give a ready and reasonable account of their faith”. Cp. his \emph{Theophilus} (edited by V. Fr. Oehler, 1878, pp. 22–25 and pp. 28–49). In the same work we read (p. 128): “A good teacher leads, while a bad one drags; the one sheds light, the other darkness; the one teaches, the other confuses; the one guides, the other pushes; the one stimulates, the other depresses; the one scatters cheer, the other fear; the one builds up, the other tears down. In short, unless the teacher himself be a book, yea, a library, an itinerant museum; unless he himself be a personification of the right method of handling and superintending the task; unless he himself be the embodiment of the genius and method of all language and science, and, in addition, an ornament and a flower of State and Church, he, once for all, is not according to my taste. For to take up books and finish them one after another, to drive and to goad to exertion, to enforce injunctions, rules, and fiats and to hammer them in—that is something anybody can do; but to lay down the main contents of a lesson, to prepare the way to an understanding of it, to make the application, to teach the correct use, to go ahead with a good example, and finally to bring everything into harmony with Christ:—that is the teacher’s true function, that is a task worthy of a Christian, for which all treasures on earth are not sufficient remuneration.” Also the Danzig school superintendent and pastor John Maukisch made an effort to improve the method. In Gemeine Nachricht, wie man die Jugend zu Hause und in den Schulen den Katechismus Lutheri abfragen koenne (1653) he advocated the analyzing method, which attained to dominating influence since the days of Spener.

\quellseite{524}
Neither is the analyzing method (pp. 135, 137, 140) recommended and adopted by Spener and Francke, and later influenced by the philosophy of Wolff, a practicable way to the true goal of education and training. Whether it was employed merely to bring into the light the structure of Luther’s explanation (as was done occasionally even in the sixteenth century), or whether by the use of the question and answer method, the several sentences were made easier of interpretation through grammatical analysis into subject, predicate, object, etc., a remedy may indeed have been found against inattention, but the method remained inadequate. It certainly contains an element of truth, but it does not even succeed in thoroughly stirring the intellectual life of the soul, let alone the feelings and the will, in spite of the fact that, in other respects, Spener and Francke were very much intent upon reaching the emotions and upon the demonstration of religious knowledge through deeds.

\emph{Spener} says in his \emph{Bedenken}: “Examinations should consist not merely of questions and the repetition of the answers found in the Catechism, but in the analysis of each question into many other questions, which, in turn, should be so formulated as to possess explanatory force, so that the young are trained to reflect upon the subject and to answer from their own mind” (vol. i, 631 f.). In the preface to his catechism (1677) he says: “As to the use of this little work, the idea is not to expect a few people merely to commit the questions and answers to memory—a torture of the memory which I would rather warn against than advise. It is my opinion that writings such as this should be used as an aid to the intellect rather than as a burden to the memory. For this reason I require that, by the use of these questions and answers, the subject should be so thoroughly imparted that people can answer from their own mind in their own words; this is better than they should commit to memory the very finest formulas. . . It may be laid down as a general principle among us that the youth entrusted to us should not only understand the truths of the faith with which they have become familiar, but also appreciate them as a divinely established rule of life and as an incentive to true godliness. Occasional exhortations between examinations for the purpose
\quellseite{525}
of adapting the truths learned to life will be found expedient. It would also be exceedingly profitable for the preacher who is occupied with this humble work to train the young to go to Scripture itself for proof by turning to the Bible, or at least the New Testament—a book which they should always have with them. Instruction ought to be given how to understand a passage according to text and context; how to analyze it; how to determine its meaning, and how to use it in proving the truth. This is a fine intellectual exercise and discipline; it results not only in such knowledge of the Bible that any passage can be readily found, but also in attention when the Bible is read and in the earnest desire to miss nothing. The ultimate gain will be that the people’s faith, which, after all, is based not on the Catechism, but on Scripture, becomes well grounded, and that the conviction is wrought in them that their Catechism is really founded upon the Word of God.” — \emph{A. H. Francke}, in his \emph{School Statutes} of 1702, says: “The method of catechization consists 1. in the recitation; 2. in the explanation; 3. in the application. 1. The teacher lets the children recite the part which he intends to treat; 2. he shows them the true meaning of every word of the Catechism, lest they rattle off the catechetical text without making the progress to be expected; 3. he points out to them how the things taught and explained should become for them (a) a solid foundation for their faith and (b) a means of testing and amending their life. All this is not to be imparted by long-winded talks but by the simple method of questions and answers, and with great love, meekness, and kindness. However, the catechetical method of some is beset by faults, which should be noted: 1. there are those who formulate the question in such a way as to require an unvarying yes or no; this has the effect of discouraging the children or so habituating them to shouting yes and no that they lose attention; 2. others ask very few questions when they catechize, but they lecture incessantly and preach regular sermons, which has a tendency to make the child restless; 3. then, it is found that some catechists do not abide by the subject; they pass from one topic to the other and go roving everywhere; but the subject of the catechization is but lightly touched or altogether ignored, whereby it is obscured rather than explained: in these circumstances the children learn nothing whatever with thoroughness; 4. the attempt is also made to arouse the flagging attention of the children by the infliction of bodily punishment; this does more harm than good. . . . Furthermore, a Bible passage should be made clear through questions; for instance, Christ gave Himself for us. Who gave Himself
\quellseite{526}
for us?—For whom did He give Himself?—What did He do for us?—What did He give for us?—When the children have learned to understand a passage in this way, the catechist should induce them to apply it; for instance, Who is it that gave Himself for us? Christ.—For whom did He give Himself? For us. (For me.)—Should we not love a Savior who gave Himself for us? Oh, yes, etc.—This method of imparting knowledge to the children through questions and answers is earnestly recommended to the teachers, not only because this is the best method of explaining a subject, but also because it tends to quiet the unruly minds of the children at a time when there is danger of feeling bored”.

\emph{Loeseke}, in his \emph{Zergliederte Katechismus} (²1758), treats the whole Catechism according to the following plan: “Which article is this? The First Article. What does the First Article treat of? Of creation, Gen. 1:1. What does it say? I believe . . . earth. Who believes? I. What is that you do? I believe, Rom. 1:17. In whom do you believe? In God, John 1:9. What do you call this God? The Father, Eph. 4:6. What kind of maker do you call this God? Almighty. As what is He almighty? As maker. Of what is God the Maker? Of heaven, Ps. 33:6. Of what else? Of earth, Jer. 31:15. How does Luther’s explanation begin? What is meant by this? What do you do when you hear all that is said in the First Article? I believe, 2 Tim. 1:12. What do you believe in regard to your own creation particularly? That God has made me, Eph. 4:24. (Who has made you? Whom has God made? What has God done for you?) Together with what has God made you? Together with all creatures. Together with how many creatures? With all. To whom has God given many blessings and goods through creation? To me. What, for instance, has God given you at your creation? My body. What else? The soul. Gen. 2:7, etc.—Rambach’s manner, indeed, is not quite so dull; a number of valuable hints as to method are given in his booklet, \emph{Der wohlunterrichtete Katechet} (1722); but he does not pass beyond the restrictions of the analytical method. Only in his posthumous work, \emph{Wohlunterwiesener Informator} (1737), noteworthy beginnings of the developing method are found, especially on pp. 140 ff. Here he clearly expresses it as his opinion that Luther’s Catechism cannot be successfully treated with the children until there has been imparted to them 1. some simple preliminary instruction concerning God and divine things, 2. a brief outline of the plan of salvation, and 3. an outline of Biblical History (p. 140–147). Moreover, he clearly recognized that the will of the children, no less than their intellect,
\quellseite{527}
needed correction. “Most informatores”, he says, “who themselves lack information, make a mistake in this matter, and fancy that they have done their duty, and more than their duty when they have imparted to the child the Catechism and a few passages, and when they have taught him to read and to write; but the state of the child’s will appears to concern them precious little. It is all the same to them when they grow up savage and uncouth, their life a field for Satan’s work. In that manner the children are made angels intellectually, but devils according to the state of their will, the aim being not godliness, but knowledge.” What he thereupon says as to the method employed in correcting the will (pp. 162–235), is worthy of attention, but we cannot enter upon it here since he has failed to show how his views concerning catechetical instruction might be instrumental in training the will.

A great step of progress is represented by the Socratic method of the period of rationalism. For the first time now, there appeared a real didactic conversation between the teacher and the pupil, a real development of the unknown from the known (pp. 142 ff.),—a method which really put the intellectual life in motion, at least when employed by its more eminent representatives, especially Dinter. However, the catechesis was too exclusively an intellectual operation, the feelings remaining cold and the will unstirred, quite aside from the substance of instruction, which rarely transcended natural religion. H. Chr. Schwarz: “The insistence upon knowledge was quite in place; but the one-sided way of stressing this principle led to the error that ignorance was worse than wickedness. The heart remaining empty, such catechizing was in the long run no better than the mechanical memorizing previously in vogue.” Cl. Harms: “Question follows question lickety-split; there is discussing, developing, explaining; example and illustration are brought in; punctuation is drawn upon for light, and assumption must do its share, and whatever else is required by catechetics is there, with all of which, if all goes well, logical efficiency is attained and the joy: Now I have caught it. But the adult who already knows
\quellseite{528}
those things and craves a little religious inspiration, cannot stand it. To ask question after question and to call for an answer is not instructing the children. In addition to questions that address themselves to the intellect, such are to be put which, like nails and spears, make for the feelings—the heart. . . Nor should we be content with questions as means to make our way to the heart. Talks, admonitions, entreaties, stories, reading by the teacher, reading by the pupils, prayer and such things are the instruments through which the Kingdom is taken by force.”

Here is an example from the \emph{Unterredungen ueber die fuenf Hauptstuecke des Katechismus} (1806 ff.), by Dinter: Teacher: Has the garden back of your house come into being of itself? Child: No, my grandfather planted it.—Who, then, was there first, the garden or the grandfather who planted it? Grandfather was there first.—Why did that have to be the case? Otherwise he would have been unable to plant it.—What do we call that which is produced, or brought about by someone? The effect.—And what do we call that through the power of which an effect is produced? The cause.—When you think of your grandfather and the garden in this connection, can you use these words? Yes, my grandfather was the cause, and the garden was the effect.—Do you think that the same is the case elsewhere? Yes.—Think of the watch and the watchmaker. The watchmaker was there first.—Think of a similar example! The cabinetmaker was there before the table.—Now think of a case where the effect was there before the cause! That is impossible: if the cause is not there first, it cannot make it at all.—It would be better to say: it cannot bring about an effect: now can you state this in general terms? The cause always comes before the effect.—Now let us apply this to God and the world! God is the cause, and the world is the effect.—What follows from this? God must have been there first, or he could not have made the world.—As soon as I think of a maker of the world, it becomes clear that there must have been a time when the world began to be. Before it was there, there was nothing but God: God made the world: by whom, now, was God made? By nobody: nobody was there but God.—Well, if He was made by nobody, He must have come into being of Himself! That is out of the question; for nothing can come forth by itself.—But is that quite sure? What is meant by saying, You have reason? That I can see the connection between
\quellseite{529}
cause and effect.—Suppose, now, someone would fancy an effect without a cause; what would you call him? Unreasonable.—But coming into being is surely an effect. Therefore there must be a cause for it.—If someone, accordingly, should say: “Something comes into being of itself”, what would you say of him? He speaks unreasonably; where there is an effect, there must be a cause; nothing, therefore, can come into being of itself.—And God? Neither can God have come into being of Himself.—Now let us put all this together: God was there before the world; nobody can have made Him; nor can He have come into being of Himself; but one possibility remains; which might that be? He cannot have come into being at all; the world came into being; that is put also in this way: it had a beginning; as to God, He never came into being.—How can that be best expressed? He has no beginning; he has no end.—Good that you remember this from the lower grade; but there you had to learn many things without knowing the reasons why we should believe that it is so; now your attention is to be called to the reasons, too.

Do you see yonder tree? What do you believe concerning it? Will it be there forever? No, some time it shall be cut down.—Why do you believe that? Because the people need wood, if they are not to freeze.—Now, when someone cuts it down, does the cause of its destruction come from the outside or from itself? It comes from the outside.—Suppose that nobody should cut it down, will it always stand there, and will it never be destroyed? Sometime it would decay and cease to be.—In that case, would the cause of destruction come from without? No, it would come from within.—This is the case with the tree; this is the case with every other thing; when it is destroyed, only one of two things is possible; where may the cause of destruction be found? The ground of destruction is either without or within the thing destroyed.—No third possibility can be imagined; let us apply this to God, and ask the question whether God can ever take an end. If God were to have an end, one of two events would have to take place. Which? The ground of God taking an end would have to be within Him or without Him.—We will speak in connection with God not of destruction, but of extermination. Why? That shall be explained in one of the following conversations; if, now, of two things the one destroys the other, which of the two is bound to be the stronger? The one that does the destroying is bound to be stronger than that which is destroyed.—Now, suppose that God, the living God, God the Creator, should be exterminated by someone, what would have to be the case? The other would have to be stronger than God.—But is
\quellseite{530}
that a possibility? No, nobody is stronger than God; for God is almighty.—What, accordingly, is no living being, no power of nature, able to do? To exterminate God.—The cause of an extermination of God, therefore, cannot come from without; but we assumed a second possibility: could the cause of extermination lie within God as it does within the oak? No, that, too, is impossible.—Why? Suppose that yonder oak tree had a duration of life of five hundred years, could it, in that case, have been there eight hundred years ago? No, it would have decayed a long time ago. How many years ago? Three hundred years ago.—Or could it be possible that it had no beginning at all; that it was there from the beginning? In that case it would have decayed even sooner.—This much, then, must be clear to you: a thing that bears in itself the cause of destruction, cannot have been there from the beginning; why? It would have been destroyed long ago.—Hence, if anything has no beginning, what are you sure of in regard to it? What can it not bear within itself? The ground of destruction.—Apply this to God! God has no beginning; for that reason He cannot bear the ground of extermination within Himself.—There is no ground of extermination within Him or without Him. What do you conclude from that? That no ground whatever of His extermination exists.—Therefore, just as surely as God has no beginning, He has . . . . can you continue? He has no end.—What it is now that you have said of God? He has neither beginning nor end.—Do you know what He is called for that reason? He is called eternal.—Right; sometimes we call Him something else; whatever comes to an end, we call what? Finite.—Him who never takes an end, we call what? Infinite.

Let us once more survey the path which our conversation has taken! What did we first say of cause and effect? The cause is there always first, before the effect; God is the cause of the world; hence God must have been there before the world.—Very well; but that does not say that He was there from eternity; someone may have produced Him. No; for no one was there but Himself.—How did we argue further? He cannot have come into being of Himself; for nothing comes into being of itself.—Now, then, if He has not been produced by anyone else nor come into being of Himself, what must we conclude? He did not come into being at all; He had no beginning.—But whence do we know for certain that He has no end? If anything is to come to an end, the cause must be found within or without it.—Which of the two, now, may be the case with God? Another cannot exterminate God because God is almighty; nor can a ground of extermination be found in God; for He has no beginning.—From all this
\quellseite{531}
follows the divine attribute about which we mean to converse today; which? God is eternal.—That means? He has neither a beginning nor an end.

In the period of a more general \emph{return to the faith of the Reformation} many catechists made the regrettable mistake of resuscitating the old method. In contrast to the Socratic method of the rationalistic period, they would simply impart the catechetical statements to the children by the force of authority, draw them out again by means of questions, demand the proof for them in the form of an occasional Scripture passage, and, possibly tell a story from life in order to shed light upon it or insure its retention by the memory. When they would really employ examples from Biblical History, they did not utilize them fully; least of all did they use them as sources of deduction for the truths of the Catechism. For instance, when a definition of faith had been given and proved as correct by several passages from Scripture, it was added: Thus Abraham, David, the centurion of Capernaum, the king’s officer, the jailor of Philippi, and others, believed. That such proceeding was worthless, is evident. Others recognized the indefeasible element of truth in the Socratic method; they saw that it does not suffice merely to impart and draw forth again by questions; to analyze, and, perhaps afterward, to demonstrate and illustrate—a process during the whole of which the attitude of the children is purely receptive, but that, through the didactic conversation, there must take place a real developing, and thereby a training of the child for self-activity (so Harnisch especially). But which was to be the starting-point for such a developing process? Entertaining the perfectly correct view that instruction in the Catechism has accomplished a great deal when by means of it an understanding of the words of the Catechism has been imparted to the children, some would proceed from the linguistic usage in which the concepts of the Catechism had found expression in the language of the Bible and the vernacular in general (Brieger
\quellseite{532}
particularly). Others would make the Bible passage their starting-point, endeavoring to develop from that the particular Catechism truth to be taught (they were of the opinion that only in that way the “formal principle” of the Reformation could come into its own); Materne, 1853, and, even more, \emph{Crueger}, 1860, represented this method. The latter says plainly: “To adduce Scripture passages as proof material, after the sentence of the Catechism has been discussed, presupposes doubt in the correctness of the doctrine, and for the reason that the passages, for lack of time, are afterward not carefully explained, it confirms such doubt.” As a matter of fact, by that kind of development (especially the type recommended by Crueger, see example below) little was accomplished for the intellectual life of the pupil, and nothing at all for the feelings and the will, let alone for practical life. Everything was abstract rather than concrete and intuitional; the explanation was not gained from the concept sphere of the child. Such a method did not only leave the heart of the child cold; it might even bring about the notion that the Bible is merely a code of doctrines requiring absolute submission. Thereby, however, joy in Scripture is rather choked than fed; and yet, to arouse and nourish it, is one of the very objects of instruction.

Something decidedly better was achieved by G. v. Zezschwitz and all those who, more or less independently, followed the path blazed by him both in theory and practise. At heart in agreement with the truth of Scripture, he traced his steps back to Dinter in regard to form. Repudiating every catechetical method but the developing one, he laid all possible stress upon the elaboration of clear conceptions and precise definitions. He became the father of the positive scientific catechization (Kunstkatechese), which, by a series of severely logical questions, forced the pupil to reach a series of conclusions along a gradual ascent of logical development, and,
\quellseite{533}
by way of summing up these conclusions, the final theme—the exact and precise determination of the particular catechetical truth to be ascertained (cf. Zezschwitz’s Model Catechization on the Ninth and the Tenth Commandments, Katechetik ii, 2, p. 562 ff. and his \emph{Christenlehre}, three vols, ²1883—1888). It cannot be denied that the co-operation of the children was secured in this way to a rare degree, and that their mental powers were subjected to an extraordinary discipline; likewise, that they attained accurate religious knowledge. At the same time it is true that, in point of contents, there was a deplorable overloading by the introduction of material not rooted in the Catechism but dragged in from biblical theology and dogmatics; that an undue amount of attention was given to the formation of conceptions; that through failure to accord all-sided recognition to the intuitional principle only little influence was brought to bear upon the emotional and volitional life; that too much was expected from the average child, and that insufficient scope was given to pupils to give unrestrained expression of their own thoughts.

In his \emph{Versuch einer biblisch-sachlichen und sprachlichen Erklaerung des Kleinen Katechismus} (³1853), \emph{C. F. Brieger} suggests a proceeding something like this: Suppose the term to be explained is “redemption”. The root-word is emo = to purchase. Inasmuch as the purchase of a human being postulates a condition of slavery, redemption means that one who had been a captive or a slave has been purchased for the purpose of restoring him to liberty, the prefix re designating a return to the condition preceding the state of captivity. That which requires redemption is, as it were, in a state of slavery, precluding freedom of movement. With this linguistic conception Holy Scripture agrees, for whenever it speaks of anything from which redemption is required, it is something which cramps and imprisons, such as enemies, distress, sin (Ps. 18:3; 102:20; 130:8). “By the use of such a method,” Brieger thinks, “the conception is built up, as it were, before the very eyes of the pupil, and impressed so much more thoroughly upon his soul because of the part taken by himself in the erection of the structure.”—Crueger, in his Entwurf einer entwickelnden Katechismuslehre (1860; ¹¹1889), explains “with
\quellseite{534}
all cheerfulness and confidence” in the Introduction of the Lord’s Prayer as follows: “According to the words of Luther’s explanation, the dear children ask their dear Father with all cheerfulness and confidence. We are reminded of the passage Rom. 8:15: “Ye received not the spirit of bondage again to fear . . . Abba, father” (8th Trin.). We Christians have received a child-like mind; by this we are prompted to pray: Dear Father! We have not received the mind of a servant, a slave; a slave fears that some evil will be inflicted upon him by his master. When a slave asks anything of his master, he is afraid that some harm may be inflicted upon him on account of his petition. He who fears that harm will befall him on account of his petition, does not pray cheerfully. We pray cheerfully when we are not fearful of harm as a result of our petition. The dear children pray not only with cheerfulness, but also with all confidence (1 John 5:14: “This is the confidence that we have in Him, that, if we ask anything according to His will, He heareth us”). We have the cheerful confidence, the joyful faith, that God will hear our prayer. We pray to God with the cheerful confidence that our prayers will be heard. This cheerful faith is the right kind of confidence when we pray. We pray, accordingly, with all confidence when we pray with the joyful faith that our prayers are heard. Faith in the love of the heavenly Father moves us to pray.”

When \emph{L. Kraussold} proceeded to revise Dinter’s book: \emph{Die vorzüglichsten Regeln der Katechetik} (1801) but instead published his own book on Catechetics from the orthodox standpoint (1843), he became the forerunner of the positive scientific catechization. E.g., the thesis: “The Grace of God in Jesus Christ, is the Christian’s Solace in the Hour of Death” he develops in the following manner: “What does sin merit? Punishment. What punishment has God decreed for sin? Damnation. What is the sensation of the person who looks for punishment? Anxiety. What, therefore, is the sinner bound to feel when he thinks of his sin and punishment? Great anxiety. Whereby could such anxiety be taken from him? (Whence anxiety? When may a man cease from anxiety? When, accordingly, will anxiety have passed away? When will anxiety disappear? Whereby, accordingly, can he be relieved of anxiety?) By being relieved of sin. What does sin merit? The penalty of damnation. What passes away at once when sin is taken away? Punishment. Who punishes? God. Who alone, accordingly, is able to remit the punishment? God. When God remits the penalty of sin, what is that called? To forgive sin. Whereby, accordingly, is anxiety removed?
\quellseite{535}
By the forgiveness of sin. When one is relieved of anxiety, we call that giving comfort. What, therefore, is it that the sinner seeks when he seeks the forgiveness of sin? Comfort. What, in consequence, is the sinner’s comfort? The forgiveness of sin. Has the sinner deserved it? No. What is it called for that reason? Grace. Who has secured for him such grace? Jesus Christ. In whom, therefore, does the sinner find the grace of God? In Christ Jesus. What, therefore, is the Christian’s comfort in his sin? The grace of God in Jesus Christ. When does the Christian stand in need of such grace? Always. Why? Because he is always a sinner. What must he do to secure grace from God? He must ask God. But how long may he ask God for it? As long as he lives. No longer after death? No. What is it that causes him torment if he has not secured grace before death? Damnation. When, therefore, is he nearest to damnation on earth? At death. When is anxiety bound to be the greater, when punishment is near or distant? When it is near. When, therefore, is the sinner’s anxiety bound to be greatest? At death. When, accordingly, does he stand in need of comfort most? In the hour of death. What, therefore, is the Christian’s comfort in the hour of death? The grace of God in Jesus Christ.”

With this model of a scientific catechization before us, we appreciate the views expressed by \emph{Schmarje}. Though himself not averse to catechizing in this manner, he is more moderate than Zezschwitz. While emphasizing the developing question, he repudiates the „\emph{zwingende}“ question of the typical scientific catechization—those severely logical questions whereby the pupil is forced by the teacher to move along the predetermined path of the latter’s reasoning, nor is he a zealous advocate of the “final theme” as the ultimate goal of every catechization. In \emph{Das katechetische-Lehrverfahren auf psychologischer Grundlage}, ²1892, p. 18, he says in regard to the scientific catechization: “The immediate impression of such a conversation is that the pupil can be conceived only as being helplessly drawn toward the goal with invisible tongs. He trots behind the teacher; but he does not know why he follows him, nor whither he follows. The logical consistency of the questions presents no gap for a detour to the right or to the left which might lead astray; but neither does it present an opening through which a glimpse might be obtained of the inner life of the child.” Upon the thesis: “We should fear God”, he himself offers the following model of a catechization conducted by him in Altona with boys ranging from eleven to twelve years and recorded stenographically: “Teacher: You know that Luther began the explanation
\quellseite{536}
of every Commandment, with the sole exception of the First, with the words: We should fear and love God, etc. . . Why does he do that? Pupil: We shall keep the commandments of God when we fear and love Him.—T.: Today we want to explain what it means: We should fear God. Which of the two expressions will it be necessary for us to explain first if we want to understand what is meant by this demand? Pupil: We should know what it means to fear.—T.: Name persons from Biblical History who showed fear. Pupil A.: Moses had fear of God when He appeared to him in the fiery bush. Pupil B.: Paul had fear, when, on the way to Damascus, he suddenly heard the words: Saul, why persecutest thou Me? Pupil C.: Adam and Eve had fear when they had eaten of the forbidden fruit. Pupil D.: The shepherds in the field, when the angels appeared to them. Pupil E.: Mary had fear when she wanted to see the grave of the Lord and saw the angel sitting before it. Pupil F.: The Israelites had fear when Pharaoh pursued them.—Name also some examples from secular history. Pupil A.: Desiderius had fear when he heard that Charlemagne moved upon him with a large army. Pupil B.: The Romans had fear when Hannibal had defeated them.—Can you name also some examples from poetry? . . . Think, e.g., of “The Ring of Polycrates” Amasis feared the envy of the gods.—There we have quite a number of examples of fear. Now tell me: Why did the Israelites have fear when they stood on the Red Sea? They had fear because Pharaoh was behind them; he wanted to take them back.—Why did Adam have fear? He was afraid of punishment.—Why did Jacob have fear when he returned from Laban and was about to meet his brother Esau? He thought that Esau had evil intentions and would take vengeance.—What is it we want to know when we ask “why”? We want to know the ground.—In how many cases may we inquire about the ground of fear? In every case.—Why? There is always some ground for fear.—You named a short time ago a good many examples of fear. Compare the causes of fear with each other! They agree with each other.—Why? When anyone has fear, he is afraid of somebody’s anger.—Is that your opinion? Have you ever had fear? Yes.—Of what, for instance? Of a thunderstorm.—What did you think when the dark clouds massed themselves upon the sky and the lightnings began to leap? I thought that the lightning might strike.—Do you mean to speak of the wrath of a thunderstorm? No.—You see, therefore, that your answer was not correct. Now compare
\quellseite{537}
all the grounds of fear with each other. No answer.—Well, what was Adam afraid of? Of punishment.—Jacob? Of the vengeance of Esau.—Of what is the child afraid during the thunderstorm? Of the lightning-stroke.—Now, if Adam is afraid of punishment, Jacob of the wrath of his brother, and the child of the lightning, you surely must be able to express an opinion when you compare these several grounds of fear. They are different.—That is correct. And yet you contended at first that all the causes of fear were alike. Can you not give the reason for your assertion? Pupil A.: Yes, when one has fear, he thinks of God.—Really? Pupil B.: No, he always thinks of that which he fears.—What does he expect of the persons or things that he is afraid of? He expects that they will do him harm.—What, therefore, is it that one thinks of when he is fearful? The person who has fear always thinks of harm coming to him.—That is very true. But tell me in other words that the grounds for fear, though outwardly different, are alike inwardly. Pupil J.: When anyone has fear, he always thinks of coming harm. Pupil K.: He always thinks of an approaching danger.—What, accordingly, is the reason for fear? The reason for fear is the thought of some unavoidable danger or a threatening disaster.—Good. Express the same thought also in this form: “We have fear when we, etc.”. We have fear when we think of an unavoidable danger or threatening disaster.—Now we know the reason for fear, and we want to remember it (It is written on the blackboard).

“Let us go a step farther. You said some time ago that, on a certain occasion, you had fear of a thunderstorm. How did you know that you had fear. I had that feeling.—You (turning to another pupil), have you, too, ever had fear? How did you know it? I felt it.—What, then, is fear? Fear is a feeling.—What kind of feeling is fear? Fear is a disagreeable feeling.—Who can express the feeling of fear in a different way? Pupil B.: Fear is a feeling of dread. Pupil F.: A feeling that excites you very much.—That is still better. It may happen that, when fear is felt, we become speechless; yes, that we have difficulty in drawing breath. Fear is a depressing, crushing feeling. The answers are all correct. But we want to keep in mind for the present only the last definition. Answer, accordingly, the question fully: What is fear? at the same time telling us how it arises. Fear is a crushing feeling arising from the thought of some unavoidable danger or threatening disaster.
\quellseite{538}
—This sentence, too, we want to remember. (It is repeated a few times and written on the blackboard.)

“Name feelings similar to fear! Pupil: worry and care.—What do you call the feeling that creeps over you when suddenly danger comes upon you? Fright.—What do we feel when the danger coming upon us is not only sudden but at the same time great and overwhelming? Terror.—Summarize what feelings similar to fear we have found? Worry, care, fright, terror, and expectation.—You have named expectation as a feeling similar to fear. When, on Christmas eve, the pine-tree has been decorated by father and mother, while you stand in an adjoining room full of expectation, do you really think that you have fear? No.—Your answer, therefore, was not quite correct. We now know what is meant by fear or what we mean when we say: We fear. What, however, did we designate as the subject for today? We are to learn what is meant when we say: We should fear God.—Give another expression for, We should fear God. We should have fear of God.—Have fear of God! What is God’s relation to our fear? No answer.—You know what I mean when I say: Jacob had fear of Esau; the Israelites had fear of Pharao. What do we mean to say when we speak of fear of God? Pupil A.: God is the one that excites our fear. Pupil B.: God is the ground of our fear.—What had we just recognized as the ground of our fear? The thought of an unavoidable danger or a threatening disaster. And now you contend that God is the ground of your fear? How do you reconcile the two statements? Pupil M.: They cannot be reconciled; the fear of God is something different.—Who is of a different opinion? Pupil C.: They can be reconciled. God, too, can send us disaster.—Explain that more fully! God has the power to do so, and it may be intended as a punishment.—What can be intended as a punishment? The disaster.—What disaster did Adam think of when he hid from God in the garden? He thought of punishment for sin.—In how far can we apply the statement that fear comes from the thought of a threatening danger or disaster, to the fear of God? The fear of God comes from the thought of a threatening danger, because the punishment for our sin is a threatening danger.—What kind of disaster does the godly man consider punishment for sin when he compares it with every other evil? As the greatest evil of all.—Why? Pupil K.: Every other disaster passes; this one lasts forever. Pupil F.: Every other disaster comes upon us in this life, while that one
\quellseite{539}
comes upon us in the other life as well.—Which parable comes to your minds in this connection? The parable of the Rich Man and Poor Lazarus.—Name just one feature of that parable that illustrates the terribleness of the punishment yonder! The rich man said to Abraham: Send Lazarus, that he may dip the tip of his finger in water and cool my tongue; for I am in anguish in this flame.—If the godly man considers punishment for sin the very greatest evil, what will his fear of God be like? Compare it with every other fear! It is greater than any other fear.—Make that clearer! Pupil J.: The godly man has more fear of God than of men. Pupil K.: The godly man fears God above all things.—Give me some examples! When the Christians were persecuted the demand was made upon them to cast aside their faith; but they refused to do it.—But why? They feared God more than any man.—Other examples! Abraham was willing to sacrifice his own son.—What did he prove by that? That he feared God.—Go on! Luther at the diet of Worms, when he was told to recant his writings.—What did he say? I cannot recant. Here I stand; I cannot do otherwise; God help me. Amen!—What do you seek to prove by that example? That Luther feared God more than he did the emperor.—We have become acquainted with one ground of the fear of God. But we must look at this ground somewhat more closely. In what light does man consider God when he fears His punishment? As the judge.—What do you think of a fear of God in which man thinks of God merely as a judge about to inflict punishment? That is a wrong kind of fear of God.—Can you give me another name of such wrong kind of fear of God? Servile fear.—Why? It is the same kind of fear which a servant has of his master.—What is your opinion of the fear of God felt by the Jews? The fear which the Jews had of God was the servile kind.—Why was that? They thought of God merely as a severe judge.—What do you know of the fear which the heathen have of God? The heathen have no fear of God at all.—How is that? The heathen do not believe in God at all, but only in idols.—That's right; in the same sense as in the case of the Jews we cannot speak of the heathen having fear of God. But they fear their gods and idols. What is your opinion in regard to this fear? It is a servile fear.—Prove that. Pupil B.: The heathen performs sacrifices in order to appease the wrath of their gods. Pupil M.: They even sacrifice their own children.—Quite true. Think of the poem: “The Ring of Polycrates.”
\quellseite{540}
There, too, we behold the servile fear of the heathen.—How? The king is to sacrifice his last jewel, lest the gods should do him harm.—Which words of Amasis can we quote as proof of the truth of our statement? The envy of the gods I dread; From birth until he joins the dead, No mortal always lives in glee.—What kind of fear, then, is it which we Christians should feel? A child-like fear.—Assign a reason for it. Pupil C.: God is our father, and we are His children.—What kind of fear this child-like fear is, I want to make plain to you through an example from life. (Here follows an incident taken from life.) The boy had fear. Of what? Of punishment.—That, too. But therewith you have not sufficiently described the ground of his fear. He feared to cause grief to his father.—And why was the thought so painful? He knew that his father loved him.—Now we have brought out two thoughts that we can give as grounds of his fear. Which are they? The boy thinks first of the punishment he has to expect, and, second, of the possibility of grieving his father, who loves him so dearly.—Apply this to the fear of God. But give a complete answer, in this form: The servile fear arises from the thought, etc. The servile fear arises from the thought of the punishment which a person has to expect for his sin; the child-like fear arises from the thought . . .—Well? Remember that we have given a twofold cause of the child-like fear. The child-like fear arises from the thought of the punishment for our sin and of the grief that we caused a loving God by our sin.—Let us remember this statement (it is put down on the blackboard).

“Let us now repeat the main statements that have been made, by finding answers to the following three questions: 1. How does fear arise? 2. What is fear? and 3. How does the true fear of God arise? Fear arises when we think of some unavoidable danger or threatening disaster. Fear is a crushing feeling arising when we think of some unavoidable danger or threatening disaster. The true fear of God arises when we think, on the one hand, of the punishment that we have merited by reason of our sin, and, on the other hand, of God, the loving father, whom we have grieved by our sins.—We have so far considered only the ground of the true fear of God. Let us view next the effect of the true fear of God. What will the servant do who fears the wrath of his master? Pupil G.: He will avoid him. Pupil H.: He will flee from him.—Apply this to the servile fear of God! When a person has a servile fear of God, he will flee from God.—Examples!”
\quellseite{541}
When Cain had slain his brother Abel, he fled to another country. Pupil D.: When Adam had sinned against God he hid from Him.—What effect, now, is the child-like fear likely to have? Think of the boy of whom I told you a short time ago. He went to his father and confessed his shortcoming; he thought of it again and again and avoided sin after that.—What effect, accordingly, has the child-like fear of God upon the disposition and conduct of people? It makes them better.—Quite right. The servile fear of God causes us to flee from God; the child-like fear of God, to flee from sin: we may accordingly say that it sanctifies us. What kind of feeling, accordingly, is the true fear of God? It is a feeling that sanctifies.—Give, therefore, a complete answer to the question: What is the true fear of God?, the ground from which it comes into being to be included in the answer. The true fear of God is a sanctifying feeling arising from the thought that we have sinned, and that we grieve the loving God through our sin.”

This example of Schmarje’s reminds us of von Rohden’s comment upon the scientific catechesis: “As, in the game of croquet, the ball is driven by many blows, both direct and auxiliary, to a small goal through narrow arches, immovably fixed at the outset, so the master of the scientific catechization, by means of many chief questions, and, at every sideward leap of the childish fancy, of even more numerous auxiliary questions, drives the thought of the pupil with unwavering determination through the previously fixed iron gates of his own pre-arranged course of thought to the goal before him—the conception formulated by the book or the teacher, a definition or something like it.” While this criticism is too pungent, it corroborates most of what has been said on the subject.—Zezschwitz, in his model catechesis on the Ninth and the Tenth Commandments, first achieves the following two results by a gradual development: 1. The Ninth and the Tenth Commandments forbid selfishness, the root of every sin against neighborly love; 2. the Ninth and the Tenth Commandment forbid evil lust, the root of every sin. These results achieved, he uses them to establish the final theme: The Ninth and the Tenth Commandments constitute the summit of the whole Law. If, at the same time, he succeeded in influencing the emotions and the will, the effect was due, as in the case with many other catechists, to his personality rather than his method. He himself appears to have felt the insufficiency of his method; for he supplemented his “dialectic-didactic” method by another, which he
\quellseite{542}
called the “teleological-paraenetic” method. The latter was to be employed chiefly in the instruction of confirmands and in \emph{Christenlehre} when the teacher discusses Baptism, Confession, and Communion.

Just when the scientific catechization (\emph{Kunstkatechese}) had attained to dominating influence in catechetical instruction, a new idea began to gain ground. Sponsored by the psychology and pedagogy of Herbart, as developed especially by Ziller (p. 151), the thought gradually gained ascendancy that the way from the “concept” to the “conception” (cf. p. 205) is through the intuition, and that therefore a fundamental principle of a correct catechetical method must be first to feed the intuitive faculty of the child, and thus to enable him to make progress by independent thinking and judgment. Another factor making for a change in the method of instruction was the increasingly advocated fundamental idea of religious psychology that specifically religious conceptions can exert vitalizing power upon mind and soul only as there is a beholding—an intuition—of the religious and moral life, of which such conceptions are the expression. Thus the view at once was cast aside that linguistic usage or the Bible verse is to constitute the starting-point in the development of the truths of the Catechism. At the same time the scientific catechization (\emph{Kunstkatechese}) could no longer be advocated; for though intuitive material taken from history or the life of the present was employed in that method, the law governing it in all its parts was not the paramountcy of the intuitive faculty. The scientific catechization did not proceed from an appeal to the intuitive faculty; that is to say, it did not, first of all, strive to let the children behold the nature and meaning of the religious life upon the background of the actual life of religious-moral personalities, thus enabling them to recognize it in its worth, significance, and beauty, and thereby in their hearts to arouse interest in it, desire for it, and the determination to lead such a life themselves. And just because this was neglected
\quellseite{543}
by the advocates of the scientific catechization, their efforts were bound to be abstract, doctrinaire, anything but pulsating with life, and, in consequence, anything but vitalizing, unless the defects of the method were atoned for in large measure by the personality of the catechist. Whatever fault one may find with the school of Herbart and Ziller, and the modern religious philosophy, it must be clear from the whole second part of this book (especially pp. 223–226, but also 253, 260, etc.) that the principle: “from the intuition of the religious-moral life to the religious-moral conception” has to be a dominant and indispensable rule of the catechetical method. The catechist has to bring the children face to face with the characters of history, in order to help them to see what truth and Christian life are. This was the method of God the Father when He sent His Son into the world, so that men might behold in Him the Father, the Father’s nature and attributes, His holiness and love. While this was not the sole purpose of sending Christ into the world, it was nevertheless closely related to redemption as it was, one of the reasons of His coming (John 14:9; Matt. 11:27). This was the method of Christ, who said to John and Andrew: “Come and ye shall see!” (John 1:39, compare 46); who told the delegates of the Baptist: “Tell John the things which ye hear and see” (Matt. 11:4); who, in order to lead men to faith, did not formulate definitions or syllogisms concerning God, His own person, and the Kingdom of Heaven, but presented Himself in person, so that, beholding His divine-human life as a Redeemer and perceiving the powerful appeal of His personality, men might obtain a true vision, gain vital experiences, and form wholesome judgments (cf. John 4:42). This was the method of Paul also, who “painted” Christ before the eyes of the Galatians as the crucified (Gal. 3:1).

The catechetical rule here established has indeed been exaggerated, with a twofold error as result. There are
\quellseite{544}
those who think that all that is necessary now in religious instruction is to picture to the children the life experiences of grace-endowed personalities in fascinating concreteness and with “contagious” warmth, so that, coming under the spell of the life so beheld and mysteriously laid hold of by God in the depth of their own souls, they might “meet God”. It is evident that, where that view prevails, the emphasis is laid upon feeling, upon patterning one’s own feelings and experiences after the life of the personalities beheld; far less interest is evinced in the thoughts and truths uttered by them or shown forth in their lives; still less is it believed to be the object of religious instruction, in co-operation with the children and with the life of such grace-endowed personalities as object-lessons, to develop definite, precise, religious ideas and truths. In that case, the intuitive principle, by reason of gross exaggeration, has led to the complete abandonment of all instruction in the Catechism: the chief parts of the Christian faith would hardly be rated as anything higher than the expressions of the religious life of a Moses, an Athanasius, a Luther, and at best would be incidentally treated in connection with Biblical or Church History. That we cannot join anyone on such a path, is clear from ch. 26.—Others thought that the observance of this fundamental rule led of necessity to the abolition of instruction in the Catechism as an independent branch of religious instruction. The intuitive material for the religious-moral life, it was held, is found in the Bible stories. When, in the process of presentation, by a detailed description of the characters of the story an attractive, alluring picture has been painted of true religion and morality, which has been made even more precise, more clear, and more appealing in the process of penetration (and if need be, of comparison), nothing more remains to be done, it is concluded, than to gather the several features of the religious life that have been set forth, into
\quellseite{545}
a terse classical expression. This is supposed to be most readily accomplished by using for the purpose a hymn stanza, a passage of the Bible, or an appropriate sentence from the Catechism. Thus the conception, it was thought, grows out of the intuitive material. The catechetical thesis would not impress the children as something discovered and dragged in from the outside, to which they must submit at any price; but as something ascertained through their own effort in co-operation with the teacher, to which they cannot refuse to yield without self-disapproval. But, by the side of this “catechetical step” in instruction in Biblical History, it was held, there is neither room nor need nor justification for anything like specific instruction in the Catechism. In contrast to the view previously discussed there is here a recognition of the value to be accorded, also in the religious life, to clear cognitions and sharply defined truths and conceptions; and there is a lively interest in working these out. While this touches a sympathetic cord in one who is bound, with Paul and Luther, to lay stress upon “sound doctrine” (2 Tim. 4: 3), we know ourselves to be opposed also to this view (compare ch. 26). We believe that it is possible to hold to that fundamental rule in all its bearings and apply it in practise without the alternative of discarding specific instruction in the Catechism.

Two possibilities have to be considered in this connection: \emph{The first} is that the catechist treats Luther’s Catechism as something not yet existing; in every catechesis he proceeds from an intuitive basis; from this he gleans some religious truth or directs the children to glean it, garbs it in the words of the Catechism, and thus builds up the Catechism by a method purely analytical (this word to be understood in the sense of p. 442), until it stands as a unit before the eyes of the children. This process is likely to comprise the following steps in every catechetical lesson: 1. The
\quellseite{546}
Catechism truth is developed from some incident in the life of the child or in history; 2. When the children have received a clear idea of the subject matter, it is summarized and, with the aid of the teacher, first expressed in the child’s own words; 3. the teacher garbs such explanation in the form of a statement of Luther’s Catechism; 4. in this form it is written on the blackboard, drilled with due attention to proper expression; then—not before—it is looked up in the Catechism, drilled, and memorized. In this way, when, e. g., the Fifth Commandment is treated, with the story of Cain’s fratricide to serve as starting-point, the following truths can be established in an ascending line through the didactic conversation: 1. We are not permitted to hurt our neighbor in his body; 2. we are not allowed to harm our neighbor in his body; 3. we should help our neighbor in every bodily need; 4. we should befriend our neighbor in every bodily need; 5. all this we should do or leave undone out of fear and love of God. These sentences written on the blackboard, are then, at the close of the whole conversation, gathered into a whole in the words of the Catechism.—\emph{The second possibility} is that the Catechism is put in the children’s hands at the outset; the part to be considered is looked up; where necessary, its structure is set forth by a careful analytical process, in which the children take part; the analysis completed, its clauses are written on the blackboard, whereupon the children are made acquainted with the several sentences or clauses by means of the analytical development (as set forth on p. 442), which always rises from the intuition of religious-moral life to the religious-moral conception. If, for instance, the Fifth Commandment is the subject under consideration, let it be analyzed into its two component parts: 1. We should fear and love God and not hurt or harm our neighbor in his body; 2. we should fear and love God and help and befriend our neighbor in every bodily need. In
\quellseite{547}
connection with each one of the four truths contained in these two sentences, the catechist will make a beginning with the intuitive material offered by a certain Biblical story or the experience of the children, in order thus to lead them by the developing process from that point to a comprehensive understanding of the specific truth to be taught. A detailed picture of the slave labor of the Israelites in Egypt, for instance, may be painted, in order to let the children vividly see and feel the impossibility for anyone who fears and loves God to hurt or harm his neighbor in the body. Or a word-picture of the Good Samaritan, with due attention to detail, is painted, in order to let his noble traits arouse the sympathetic, moral, and religious interest of the pupils and to move them to the definite resolution to follow his example. Or if the Second Article is to be explained, Luther’s explanation is first analyzed, whereupon the several clauses may be developed in the way of induction (analytical), with an intuitive basis as starting-point in each case. In order to enable the child really to apprehend the truth “without Christ, man is a lost and condemned sinner”, a word-picture of Adam and Eve at the moment when God drove them from Paradise should be painted, or of some heathen from the history of Missions, whose conscience has been awakened, so that the children may read in the outlines of that picture, and in a measure reproduce in their own feelings, what the sentence from the Catechism expresses.

—There the Catechism, in its several clauses, and finally as a whole, is the result of the discourse; and here it is the starting-point; but in both cases the fundamental rule: from the intuition of the religious-moral life to the religious conception, dominates the process. In the first mode it lies at the root of the process as a whole; in the second it is of moment at least in the teaching of some specific truth of the Catechism. In both cases the important element of progress
\quellseite{548}
introduced by the method of rationalism as well as the scientific catechization is given due consideration; we have a real didactic conversation, in which the pupil is led from the unknown to the known, i. e., results are gained through diligent development.

Which of the two modes described deserves preference? The superiority of the first, at first glance, appears so great that there is no room for the second. All learning is based upon apperception (p. 212); for this reason all teaching of new truth, if it is to be correct methodically, is bound to take place in connection with such truth as is already found in the soul. This principle is fully observed in the first mode; for the starting-point is not some foreign quantity, as “Catechism”, but a personal experience or a well known story: therewith the cell-wall is supplied, as it were, against which the new cell is built. As a matter of fact, however, the case of the second mode is by no means so bad as might seem at first sight. It should be borne in mind that, for pupils who have been instructed in Biblical History according to the rules laid down above, the Catechism has ceased to be a “foreign quantity”. While, so far, it has not been presented to them as a unit or treated as such in special lessons, it is nevertheless true that every Biblical History period had taken up a certain part of the Catechism and garbed the truth gleaned from Biblical History in its language. It is not likely that there are many sentences of the Catechism which have not already been elaborated by the instruction in Biblical History and have not yet become familiar to the pupils. The need hardly exists at that juncture to construct the Catechism independently. On the contrary, the child is prepared for it, and desires to become acquainted with it as a whole. The child has by this time attained to such intellectual maturity that this life manual—intended and qualified to serve as fingerpost and compass
\quellseite{549}
for his whole religious and moral life, by the help of which he is to become more and more familiar with Holy Scripture—can safely be put into his hand as a summary of the faith confessed by the Church, full membership in which is the goal of his present tutelage. The second method is the very expedient calculated to impress upon the child the difference between the instruction in the Catechism now in progress and the instruction in Biblical History which preceded it and still continues as a parallel course. There is no doubt that the pupil will become more firmly grounded in the Catechism; for it is the Catechism which he now has always before him, and which now forms the starting-point for every catechesis and, properly understood, also its aim. Nor is it true that Christian children lose interest when the aim of the catechesis has been announced to them at the outset in the words of the Catechism; as, for instance, when a new general subject is announced in words such as these: “The next lessons shall teach us what Luther’s explanation of the First Commandment requires of us Christians”. The same holds good when some subordinate point is treated, as, “Today we want to learn how, as Christians, we should trust in God above all things.” There is the additional advantage here that the children are enabled at the outset to look upon the intuitive material to which they are introduced from the definite point of view indicated by the statement of aim, and that a close inner relation is at once established in their souls between the truth expressed in the Catechism and the intuitive material laid before them, so that the one cannot awake in the memory without the other. For this reason the second mode should receive preference as a rule, especially in the circumstances prevailing in our country. It adopts all the elements of truth contained in the catechetical methods heretofore employed and presents the correct method in a fusion of all of them. The
\quellseite{550}
catechist here presents the whole Catechism and holds it up to the young as the sum of truth, experienced and proved as such by himself and the whole mature congregation:—this is the element of strength in the method obtaining at the time of the Reformation. At the same time the catechist endeavors to analyze the several parts through questions, and thereby to bring about, or, in any event, facilitate, a rudimentary understanding of them:—this is the truth in the method of pietism. While constantly stimulating the pupil to independent thinking, he makes it his object to effect a real comprehension of the subject by means of the didactic conversation, and thus to put in operation the element of truth found in the scientific catechization as well as the Socratic method. But all this is done only by leading from the intuition of the religious-moral life to the religious-moral conception, whereby the compatibility of this method with the fundamental demands of present-day psychology stands demonstrated.

It remains to add one, final, requisite of correct method: The whole didactic conversation must pursue aims altogether practical. The needs to be satisfied should not merely be of an intellectual nature. While these are by no means to be overlooked, the molding of the life should invariably be the final object. Not this is of moment that the children, on the basis of the intuitive material supplied to them, are able to give a correct definition of “trusting”, but that, at the sight of David risking unequal combat with Goliath on the strength of his trust in God, they receive a vivid intuition of living faith; that their hearts are warmed by what they see; that they recognize the great value of trust in God; that they begin to perceive the utter inadequacy, the sinfulness of any other trust; that they recognize the fact that there are situations in their own life where everything depends upon just such trust; and that, if it please God
\quellseite{551}
(p. 238, 278), they leave the room with the resolution, even now as children, ever to trust in no one as they do in God. It is only thus that the personal character of Luther’s explanation with its “I”, “we”, “us”, comes into its own (p. 105); it is a direct challenge to the catechist not to content himself with the rôle of a mere preceptor, but by all means to be personal and practical in his instruction, and to keep in mind the individual and the molding of his life. Also here the threefold life of the soul is to be set in motion, which means that the intellect with the emotional and the volitional life must receive harmonious attention. Quite true, clear conceptions are to be achieved; for hazy ideas are seldom effective; but not these conceptions in themselves should be the aim, but only as a power for the practical Christian life. Luther may well serve as our pattern in this respect, who, no matter what important truths he states, e.g., in the First or Second Article, in conclusion always establishes the connection between them and the practical life.

Accordingly the process of a catechization (especially on the First, Second, and Third Chief Parts) will be as follows: \emph{(1) Statement of aim}, in order that the children collect their thoughts and concentrate them upon one point. The stage of preparation may usually be safely passed over, especially since the new catechesis is ushered in by review and drill questions concerning the substance of the preceding catechesis, whereby the children are led up to the point where a new start is made. The statement of aim should be couched in a single sentence of the Catechism, as, “Today let us learn that Christians should use the name of God for the purpose of praise and thanksgiving”; or, “Today let us learn that no one has ever become a believer in the Lord Jesus by his own reason and strength”. \emph{(2) The developing of the specific truth from the intuitive material}. If this truth
\quellseite{552}
is to be of value, result in clear cognitions, and move feelings and will, the intuitive material should not be merely touched upon: all its points should be utilized as far as they may contribute to the feature indicated in the statement of aim. Only thus can an adequate picture of God’s or man’s conduct be produced. There may be occasions when it will be necessary to bring in a second or a third story, in order to obtain all the features of the catechetical truth under consideration. When, for instance, trust in God, as required by the explanation of the First Commandment, is to be set forth, it will be advisable to adduce not only David’s combat with Goliath, but also his conduct under Saul’s persecution; for while the former story illustrates trust in action, the latter illustrates trust in suffering. The closer the relations between the story drawn upon, whether from the Bible or life, and the truth of the Catechism, the better: another reason, why nothing has been gained when a story is merely touched upon. \emph{(3) A comparison with similar or even contrasting stories} will often help to bring out more forcefully the truth under consideration. Thus David’s trust in God is rendered more conspicuous when compared with Goliath’s self-trust. Likewise, the impression made by David’s trust is deepened when paralleled with that of Moses which prompted him to pass through the Red Sea. God’s faithfulness appears so much greater when viewed against the background of man’s unfaithfulness. If the specific truth has been developed from the intuitive material and its true nature been more fully recognized by comparison, there follows, \emph{(4) a valuation} of the action of God or of such persons as have been used for intuitive material, or drawn upon for the purpose of comparison. The children are trained to rate the value and significance of an action in itself and in the light of its results; human acts are to be examined as to whether they are good or bad, worthy of emulation
\quellseite{553}
or of detestation. Thus has already been introduced what need not be emphasized as a special step in every case and must occasionally be omitted altogether, namely, \emph{(5) the application of the truth to the children’s own life}. Such application is effectual and in place only when it does not form a mere appendage but forces itself upon the mind as the natural conclusion of what has gone before. If, e.g., the high value of trust in God has been recognized, questions like the following will naturally suggest themselves: What shall you do when you are to solve a difficult problem; when you are to undertake a dangerous trip; when you are alone at night; when a severe thunderstorm is coming up; when you should lose your father; when poverty oppresses you; when you might become rich by leaving the right path, etc.? If the children, in connection with the Second Article, have been made conscious of the fact that Jesus Christ is true God, the question will naturally suggest itself: Who, then, will be your comfort in trouble, even in the greatest trouble; but what, in view of that fact, will you always have to guard against? If the truth has dawned upon the children that their Lord and Redeemer is true man also today—the same man who has gone through every temptation here on earth, they are shown how to apply to themselves the comfort given in the Epistle to the Hebrews, that the Saviour will surely have pity upon our weakness. If a description has been given of the price of our deliverance from sin, death, and devil, namely, Jesus’ blood and life, the question will at once occur how God evidently judges of our sin, and what care we should take to guard against it, etc. In connection with the fourth or the fifth step the Bible verse or the hymn may usually be drawn upon with advantage. In connection with the fourth, the Bible passage will usually, though not exclusively, come in as \emph{dictum probans}—in connection with the fifth usually as an admonition for life (p. 379).
\quellseite{554}
It is still indispensable as \emph{dictum probans}, modern objections to the contrary notwithstanding. For instance, when the redemptive value of the work of Christ is under discussion, who is authorized to render a final valuation in the premises save (in addition to the prophets) Jesus and the apostles? Who can be relied upon for a thorough and trustworthy disclosure of human sin and helplessness and of saving grace as well, man himself, ever ready to deceive himself and to trust in his own power for help, or God; the patient who believes himself much sounder than he is, or the physician, who understands the disease and is aware of the healing power of his remedies? To be sure, the Bible verse can be applied in a manner so mechanical that the children are led to feel that the Bible is merely a code of doctrine, to whose injunctions and statements they must blindly submit. But when, again and again, the truth is emphasized that, in a history of over two thousand years, the Bible has shown itself as truth and life; that it is the grace of God which has given it; that true association with it has never yet brought forth slaves but always free children of God—happy and cheerful in every trouble, joyous and bold even when face to face with death; when the children are made to feel that it is holy and precious to the catechist himself; that it is the true book of life, such thoughts will not arise. What the children need here, as in connection with instruction of children in general, is a little confidence in the greater experience of the teacher. Often it may be advisable already in connection with the fourth step, and, still more, with the fifth, to introduce passages from the Bible as testimonies expressing the experience of those holy men, provided only that there is no encroachment upon its absolute correctness and universal application. Cases in point are Ps. 23 in connection with the First Commandment and the First Article; Rom. 8:31 ff. when the blessed state of God’s children is under
\quellseite{555}
discussion; Ps. 145: 18. 19 in connection with the Second Commandment. Whatever Bible passages are called for under point five, should bear a practical character throughout, with Ps. 119: 9 and Ps. 19: 8—12 governing the perspective. When the time has come for the class to form a valuation of the act under discussion or to apply the truths derived, it may be appropriate to read lengthier sections from Scripture, as, Isa. 53 in connection with “purchased, redeemed, and won . . . with His holy precious blood, etc.”; Heb. 11 in connection with the treatment of the nature and the fundamental character of faith; 1 Cor. 13 in connection with the subject of Christian love of the brethren, etc. (p. 371). The hymn takes its place by the side of the Bible passage as the expression of the experience of the later generations and, at the same time, as evidence of the fact that also today men may have the same experiences which those ancient men of God experienced and expressed, under singular divine guidance, in Holy Scripture. There are occasions when whole hymns may be aptly quoted.

—The final requirement is \emph{(6) the drill}. If a catechism with explanation be the textbook, the section dealing with the subject is looked up when application is finished. It goes without saying that it will be of service only when it is the result or thesis of the catechization that had been conducted in the class-room. This section should then be read by the pupils at the close of the catechization. At home it should be studied. The next day, before the new subject is taken up, a number of questions are asked for the purpose or reviewing and drilling. The questions should not require answers that literally reproduce memorized material. It is better that merely the essential facts be recapitulated, as has been emphasized on p. 447, nor should they be restricted to intellectual features alone; see Practical Examples at the end of this section. If the catechism in use contains no
\quellseite{556}
explanation, a notebook should be procured by the pupils. In that the drill questions asked by the teacher should be entered and studied until the next lesson, when they are to be answered in writing or orally. For light upon this point, as well as for information how to give the catechization a practical trend, we refer to Practical Examples. At the same place is to be found a catechetical “excursus”, such as may be undertaken at the end of a Chief Part or after completing the Catechism in all its parts.

It has been emphasized above (p. 320) that the subject for catechization is not the Catechism text together with Luther’s explanation, but only the latter. In connection with the Ten Commandments it may possibly be necessary to enter briefly, and in a preliminary manner merely, upon the Old Testament significance in order subsequently to ascend to the lessons they contain for us Christians according to Luther. In connection with the Second Chief Part a verbal explanation of the text together with an outline will suffice; at most a few historic remarks are needed, so that afterward all available strength may be devoted to the setting forth of what we Lutheran Christians mean to confess in the Apostolic Creed, which is so often upon our lips. When, in the Third Chief Part, which requires but brief treatment, and in the Fourth and Fifth, which require a lengthier and somewhat different treatment from that of the first three, the biblical text is made the starting-point, the whole exposition will issue in Luther’s explanation.—When, finally, the requirement is made that the explanation of the Catechism should be \emph{Christo-centric}, all that is needed is to follow the trend of Luther’s explanation and the \emph{Christo-centric} character of the explanation will take care of itself. Cf. especially pp. 312 f., 324 ff., 351, but also 346.

\quellseite{557}
\chapter[§ 36. TEACHING THE REMAINING MATERIAL]{§ 36. METHOD OF TEACHING THE REMAINING MATERIAL}

Compare the literature in chapters 26, 29, and 30.

The \emph{hymn} cannot generally be accorded independent treatment while our school system remains subject to present conditions. We shall have to be content with illuminating it, as occasion may require, now as a whole now in part, in connection with instruction in Biblical History, the Catechism, and Church History. When it is there brought in at the right moment, read by the teacher in an impressive manner, and the pupils are then drilled in it in order to read it well and with expression, and, moreover, are given a short explanation of those words that need an explanation, all is virtually done that is necessary for its comprehension (p. 379). When a festival approaches, selection is made of some striking festive hymn; festive feelings are aroused in the children's minds; they are shown how those festive feelings are given expression in the hymn, everything unintelligible meanwhile being explained as briefly as possible; it is then read in an impressive manner, at first as a whole, then the first stanza; the class is thereupon taught to read the first stanza with due attention to expression; last of all, the teacher sings or plays the melody and then practices it with the class. By this time the children almost remember the words, so that it will soon be possible, perhaps in the next lesson already, to have a few stanzas sung from memory. Frequent reviews, especially through singing, will make the children's knowledge of it secure. That no hymns should be assigned for memorizing without having been treated in the manner here set forth, should appear self-evident. Should there be time (cf. also p. 431) to devote more attention to
\quellseite{558}
a hymn, let care be taken to avoid the superabundance of other material under which teachers are wont to bury it. Even where more time is available, it will be sufficient to treat the hymn according to Schueren’s suggestions (pp. 119 ff.), whom we quote in the Practical Examples given at the end of this section.

In what manner the reading of the Bible is to be connected with the instruction in Biblical History and in Catechism, has already been shown (pp. 371 f., 416, 532 f.). Where Bible reading is treated as an independent branch, the following order will be found to answer the purpose best: 1. The teacher, using the acroamatic or erotematic method, establishes the context to the extent that is desirable and practicable, with frequent recapitulations covering the whole Gospel or Epistle; 2. when such review does not suffice, he prepares the ground for the new section by pointing out the fundamental thoughts and reading the whole with due expression; 3. he lets the children read the section (those rather intelligent are called upon first, then the backward ones; the more important parts are read in unison); he intersperses the reading with the necessary explanations (only such), by paraphrasing difficult expressions, dissolving involved sentences, explaining geographical, archeological features, etc.; 4. he has the class trace the fundamental thought, connects it with something already known, etc., and concludes with the question: “What does this section teach you, children?” Especially when the Epistles are read, one must guard against losing oneself in details, but try to discover and pursue the thread of the discourse: when the Gospels are read, it is advisable occasionally to read the parallel sections, or to have the brief account in the section read supplemented by additional features gleaned elsewhere in Biblical History, and thus to explain Scripture by itself (cf. Schueren, p. 47 ff.).

\quellseite{559}
In regard to \emph{Church History} we generally have to confine ourselves to a rousing narration of epochal events, with special emphasis upon educative features, and to drill-questions. Where the requisite amount of time is available, one may readily accommodate himself to the “formal steps”, among which presentation and association require the greater amount of care.—In the \emph{Young People’s Society} more freedom is naturally vouchsafed. While the acroamatic form of instruction is bound to preponderate here, the need for self-activity should not be overlooked. Where a City Library or a Young People’s Library renders the requisite literature available, historic sections, delineations of characters, etc., may be read, followed by brief essays on the subject, which, in turn, can be made the starting-point of discussion. This applies especially to such questions regarding church and religion which are at the time being of general interest (p. 432). Also questions relating to church usage and discipline (p. 433), difficult passages of Scripture, etc., can be illuminated in the manner suggested. It falls to the leader to remain in control of the discussion and, at the end, to give a lucid, concise, and impressive generalization. But the converse rule, too, will be found available: the leader delivers a lecture, and the discussion follows. When, for instance, in connection with the biography of Athanasius, the Nicene Creed becomes the topic for discussion, the young people may be invited to show at the next meeting in how far inalienable elements of the Christian faith were at stake; in how far the testimony of Scripture favored Athanasius; in what churches views akin to those of Arius are found today. Or, the origin of the Papacy is pictured along purely historic lines, without anticipating a verdict, whereupon attention is directed to such Scripture passages as Matt. 16 and others and the proof furnished why these can in no wise support the papal idea. Or, the origin of monasticism has been made the subject of
\quellseite{560}
a lecture; this is followed by a discussion in which the fundamental traits of a truly Christian ideal of life, with its stress upon the duties of the secular calling, are substituted for the monastic ideal. The fact that, under Constantine, the Church becomes a State Church, supplies an opportunity to exchange opinions upon the relation between State and Church. A description of the ruin of the Medieval Church may issue in a series of questions regarding the scripturalness of the distinctive doctrines of the Catholic Church or of the legitimacy of the charge made against the Reformation that moral life did not decline until after its advent. A study of the life of Luther may give rise to a discussion of the fundamental difference between the Lutheran and the Roman Church, of the meaning and the bearings of the doctrine of justification by faith alone, of the relation of Luther to Paul, of the difference in the conception of faith in the Lutheran and the Roman camp respectively (cp. especially \emph{Reu, Luther’s Life, Sketched for Young People’s Societies, and the Necessary Directions for General Discussion Appended}, 1917). A lecture on the Thirty-years’ War will furnish an opportunity to discuss the ultimate aims of Rome in regard to our country, etc.

\quellseite{561}
\chapter{§ 37. METHOD AND PERIOD}

F. Zange, \emph{Schulagende}, 1893. O. Zuck, \emph{Schulandachten fuer das ganze Jahr}, 1861. Knaudt, \emph{Schulandachten}, 1910.

There must be careful preparation for every catechetical period. That the catechist is constantly occupied with the study of catechetical material and will not ignore new facts as to method (cf. ch. 38), is a self-evident postulate; but each individual period, too, requires preparation. The catechist must have clearness regarding the ground to be covered, not only in general but down to the last detail; regarding the aim to be fixed, both the outer and the inner aim; and the way to reach it. While for the beginner, it is advisable, under certain circumstances necessary, to write down his catechesis in minutest detail, it will later be sufficient to draft a sketch; eventually a careful mental survey of the subject matter will prove adequate. The catechist will do well by providing a notebook both for his instruction in Biblical History and the Catechism, in which should be entered whatever he has found valuable for the several steps of preparation, statement of aim, presentation, penetration, application (Biblical History) and, likewise, for the intuitional basis, comparison, valuation, application (Catechism). Entering the class-room thus equipped, he will also recall what the next paragraph has to say about his personal relation to the children and the Word he is to offer them; then he will invoke the blessing of God and address himself to his task.

Every religious period is to be opened with prayer. Whether the class speaks in unison with the catechist or the catechist speaks the prayer alone, it should always be a real prayer, brought before God with a concentration of all the energies of the soul. The example of the teacher is of
\quellseite{562}
great importance here, the more so as many children are under no constraint at home in this respect, often even without opportunity of hearing prayer. Prayer, therefore, is not the least part of the education to be imparted. It is, furthermore, often advisable to connect a hymn with the opening prayer, which should be sung with vim and vigor; while it may bear in point of content upon the character of the lesson, this is by no means essential. Often one may simply sing a few stanzas of a morning hymn and read the morning prayer contained in the Catechism. Our fathers knew why they instituted this practise; often, when the lesson took place in the afternoon, it was brought to a close with the evening prayer.

In the partial parish or catechumen school, the opening prayer will at once be followed by Biblical History, when the new story should be presented and afterward read by the class several times from the textbook; an opportunity will thus be presented for a brief explanation of an occasional obscure word. After a Penmanship period or the like, arranged for the purpose of resting the children’s minds, there follows the penetration and application of the story which was presented on the day before, and meanwhile has been studied by the children at home. A period as provided by the curriculum of the Grammar school, and the singing of a hymn, brings the forenoon to a close. The afternoon is entirely devoted to branches of study prescribed by the curriculum of the public school. On the next day instruction in the Catechism takes the place of Biblical History. The opening with song and prayer is followed by a recitation in unison of the whole Chief Part of which a part is to be treated in detail on that day. In this the teacher will insist upon scrupulously exact, precisely articulated, but at the same time euphonious, speaking. Then comes the drill in the lesson of the previous Catechism period, to be followed
\quellseite{563}
at once by the new catechization. The rehearsal of the material to be memorized is connected with penetration and application (Biblical History), or with the drill (Catechism). The preparation for the new material to be memorized is connected with presentation (Biblical History), or with the new catechization (Catechism). Should one say, This arrangement does not allow a successful covering of the secular studies prescribed by the regular curriculum of the public school, we answer from our own experience and that of many others that this objection is not valid. In rural districts parents often are satisfied if their children finish Grammar school a year or so later. In this case instruction in Biblical History and Catechism can be given on each school-day—the one in the forenoon, the other in the afternoon, and the remaining hours can successfully be devoted to written work based upon the religious lesson, to Bible reading, Bible Geography, Church History, hymn drill and explanation.

The curriculum for the \emph{Saturday and summer school} has already been determined (p. 429 f.). In the \emph{regular parish school} instruction will follow the same general lines as in the partial parish or catechumen school. For the \emph{Sunday school} the following order is recommended: 1. Liturgical opening; 2. the leader tells the whole Sunday school or, at least, the whole Second Course (children from 7–12 years of age), the new story as clearly as possible, brings in at the proper place the material to be memorized, and lets the presentation distinctly issue in the golden text or the principal hymn stanza. 3. Thereupon the pupils retire with their teachers into their respective class-rooms, where the story that had been presented by the leader is read from the Biblical History and the material to be committed to memory repeatedly rehearsed (speaking in unison!); thereupon the teachers review the lesson of the previous Sunday by means of the review questions appended to the lesson. 4. Finally,
\quellseite{564}
the whole school with all its teachers assembles again, and the service is concluded according to the liturgy.—The \emph{Young People’s Society} has an order of services as fixed by its constitution: a hymn, prayer, the Bible lesson with informal discussion, hymn, lecture with discussion, closing hymn—these are the fundamental features generally observed.

\quellseite{565}
\chapter{§ 38. METHOD AND MAN}

Augustine, \emph{De catechizandis rudibus}, comp. ch. 7. G. v. Zezschwitz ii 2, 1, § 30 and 31, 1869. Th. Harnack, \emph{Katechetik}, pp. 114–118. F. Zange, pp. 46–57. A. Eckert, pp. 175–195. J. Berndt, pp. 73–78. E. Chr. Achelis, pp. 345–356.

But what kind of man should the catechist be, if he shall succeed in imbedding and anchoring the catechetical material in the intellectual life of the catechumen, in arousing his interest, and in guiding him in the paths in which the Spirit, in due time, shall lead him to personal faith? While much of what has been said in the second part of this book might be resumed here, some features of superlative importance remain to be pointed out. What has been said in regard to method finds its supplement in what is here said about the person of the catechist.

The first requirement of the catechist is a good \emph{education}. Theological education first of all. That he knows better and more accurately than the young what things are to be taught, is not sufficient; there must be mastery of the \emph{material}; that is, it should be known in its organic connection and be a real mental possession, so that all details may be given in their central bearings. Only one who has a systematic comprehension of Christian doctrine is able to recognize the relative unimportance or importance of any doctrine; only he is conversant with the threads that connect the individual doctrine with the sum of truth; nothing is viewed by him singly, but always as a part of the whole. Where systematic training is absent the catechist will lose himself in the mass of detail; will not know how to ask questions properly; nor shall he be able to rate the answers according to their content. He stumbles from one embarrassment into the other, and impresses one as a man who has
\quellseite{566}
no sure ground under his feet. He who would be an efficient catechist should accordingly study dogmatics. While it is perfectly true that catechetical instruction should never bear a dogmatical complexion, it is just as true that the safe foundation required for instruction is an outgrowth of a diligent study of dogmatics: what the spine is to the human body, dogmatics is to catechetical instruction. But inasmuch as all his teaching material should be based upon Scripture, and catechetical material has been taken from Scripture, the continued study of dogmatics should be accompanied by a pursuit of exegesis. Only as he goes down, again and again, into the ocean of redemptive truth gathered in the Bible, does his instruction possess the required vigor and, with all freedom, the stamp of noble dependence. Alongside of theological education dialectic versatility is required. The development and the drill particularly are dialectic exercises. But the dialectic process has to be learned and practiced. The laws that control thinking have to be understood if one wishes to think logically himself and teach others to do likewise; the mind has to acquire the necessary agility and alertness. Attentive hearing of model catecheses, and, wherever that is possible, diligent study of good printed catecheses is, in addition to the study of logic, the best aid in the acquisition of dialectic skill. The beginner had best put down his catecheses in writing; even the mature catechist should not shrink from mentally fixing the outline of the catechization—according to its major and minor points. While this is desirable from the viewpoint of dialectic clearness and skill, it is even more so from that of the importance and responsibility of such instruction. Finally, the main results of the psychological science should not be foreign to the catechist; by constantly taking into consideration the psychic life of the children he will make his task easier for himself and more successful.

\quellseite{567}
The second requirement to be met by the catechist is \emph{love to the children}. The heart of the child, and that of the adult catechumen likewise, is hungry for love; and love is what the catechist owes his pupil as a Christian, still more as a trainer. Without love, the training effort is without a soul. To be sure, the love here spoken of cannot be that earthly, effeminate, spineless love which rates the children according to their several degrees of endowment and rank, becomes extremely partial in consequence, and, for this reason, arouses in the pupils envy, jealousy, and a feeling of being slighted. The love of which we speak is rather that which is brought to bear upon all alike and which has its ground in the thought that they are children of God that have been entrusted to him. Such love is not blind to the sins and evil propensities of a child; it rather clears one’s vision to see them and generates the power to fight them with irrepressible determination. Such love can inflict punishment—painful punishment; but the punishment is inflicted for the purpose of correction, and the child feels that love has inspired it. Though the severest reproof be administered, and even corporal punishment be resorted to—a penalty never to be inflicted except in extreme cases, the child searches for love in the eye of its monitor nevertheless; and when it is beheld, together with sadness over the necessity of administering punishment, he will not only willingly submit to it but permanently benefit by it. Such love finds the tone that reaches the heart; from it flows that \emph{hilaritas} already demanded by Augustine, the counterpart of all acerbity and melancholy—a ray of vernal sunshine in the school-room. Such love will never weary of descending to the intellectual status of the child and becoming a child among children. From such love flows that faithfulness which does not cease improving one’s methods and raising the instruction to a higher level of efficiency; which does not despise
\quellseite{568}
the task of teaching children; which works with gentleness and humility in the field of catechization, though the opportunity to shine before men is utterly absent and one’s efforts are seldom controlled by any one, knowing that the Church is blessed by one’s labors; which does not permit itself to be embittered or to be turned aside from its care for the young by thoughtless men; which always bears in mind the future of the Church and the welfare of the children, and conscientiously “redeems” the time of youth, when the spirit can be molded like wax (Eph. 5:16). Such love will not forget prayer for the catechumens, least of all for the frivolous and obstinate.

The third requirement is a \emph{believing personality, sanctified by the Holy Spirit}. If the souls of the children are to yield a real anchorage to those sacred truths; if the will also is to be moved by them, the influence brought to bear upon the intellect through a clear and intelligible presentation and a lucid and logically correct development must be supplemented by a stir of the feelings; for the way to the will is found only through the intellect and the feelings. Not only interest in the truth, but also the sympathetic, ethical, and religious interests, for which the feelings serve as bearer, require cultivation. This, to be sure, is already in a measure accomplished when the catechist, during the process of instruction, brings out how one ought to share in the joy and the grief of his fellow-men, how goodness invariably finds its reward, and wickedness its punishment; when he makes the souls of the catechumens acquainted with the characters of the Bible and of Church History—men and women morally holy, meek in the presence of God, upheld by faith; but most powerfully such feelings and interests are aroused, and most enduringly the will is prompted to similar conduct when all this is surmounted by the example of the teacher as of a man intent upon the truth alone, \emph{full}
\quellseite{569}
\emph{of sympathy for whatever of grief or glee befalls his fellow-men, a champion of the right and quickened by sincere godliness}. There is a movement with which the world of pedagogy is now astir, which advocates the axiom: Religion is taught, not by word, but by example. That is a dangerously one-sided truism, an undervaluation of the Word and of knowledge, a slighting of the historic side of the Christian religion: but the element of truth in this movement, that, in the achievement of enduring effects, felt also by the will, the living example of the teacher whose life is an embodiment of Christianity, is necessary—this is a fact hardly capable of exaggeration. All strictness of discipline, all didactic accuracy and punctiliousness, all dialectic skill and faithfulness in preparation, will prove utterly inadequate where such personal influence is lacking; where the children are not made to feel that what the teacher offers to them appears precious, holy, important also to him—is, in truth, the greatest treasure and the highest norm of his life. The instruction imparted by such a catechist is never dry: it is ever fresh and vivid; for everything said by him is constantly born afresh, as it were, from his life of faith; a grasping of truth constantly renewed; a subjection to the Word willing and incessant. They are not empty formulae which he handles like lifeless tools; facts and verities they are to which the soul draws near for sustenance! Many a thing once upon a time ushered into the soul of the child through careful development may be forgotten when the graciously earnest and genuinely godly form of the old teacher and shepherd still abides in the memory as a vital force. Such a teacher cannot fail to secure and maintain that degree of authority over the children which is his due. Little need will he have for pains and penalties to enforce authority: his very presence and attitude will fascinate the obedient children and impel the obedience of the others. Without exception they will all feel: Here
\quellseite{570}
stands one who has the right to teach us; for he does not teach what he himself fails to believe and to practice: a Johannine personality is he; to lead us to Christ is his aim.

Finally, the catechist as well as the homilist should have a \emph{living faith in the creative power of the Word of God}, which is the merchandise and tool of his trade (Matt. 13: 33; Rom. 1: 16; 2 Cor. 2: 15–17). Likewise he should possess that patience and perseverance which can wait until the Word of God imbedded in the heart begins to stir and reveals its vital power, saving or judging (Mark 4: 26–29). It is this thought of the divine power in the Gospel; of the receptivity of the chatechumens in the time of their youth; of Baptism, through which those children objectively became children of God; and of the transcendent importance of his labors for the future of the Church, which, again and again, fills the catechist with strength and cheer for his calling. And the same thought tides him over the failure which, upon more occasions than one, confronts him.

This is the ideal. A zealous suppliant at the heavenly Father’s feet shall reach it more and more. He shall be enabled rightly to anchor in the whole inner life of the young those holy truths which sustain the life of the mature congregation, to arouse an interest therein and desire for participation in the divine services, so that, when the Holy Spirit, in His own time and hour, arouses them to faith by means of those sacred truths or of a renewed presentation of the Word, they shall be able to exercise their rights and duties as members of the mature congregation.

\quellseite{571}
\chapter[PRACTICAL EXAMPLES]{PRACTICAL EXAMPLES\footnote{Where the name of the writer is not given, the examples are written by the author of this book.—We purposely give not only models after which to pattern, but also some which are designed to evoke criticism.}}

\section{1. Instruction in Biblical History}

\subsection{1. PREPARATION:} \emph{(a) Aim: We would learn today how God made the world and all that is in it.}

Children, how beautiful is this world in which we live! Which season of the year do you like best? Yes, springtime is the most beautiful. Why do we like this season best? Certainly, because everything again becomes fresh and green. The snow melts, the grass grows in the meadows, the trees shoot out their leaves, the flowers blossom, and the birds return again with their songs. How did nature appear in winter? Yes, and the snow was like the white cover which is placed over a dead person. Then what happened to all nature that seemed to be dead? It was as though born again, made new all over. What do we hate to do when everything outside looks so fresh and inviting? The house seems too small, and we eagerly hasten out of doors, breathe deeply the fresh air, and try to outdo the birds in their joyful song. Who makes everything so new and fresh and beautiful in springtime? Correct, Frank; we could never have caused snow and ice to disappear. But the almighty Lord has done this. He let the sun shine more warmly, which caused the ice to break, the snow to melt, the grass to grow, the flowers to bloom, and all things to become like new. God does this every year, but there was once a time when He did much more. He created this great, beautiful world out of nothing. Think of it! What a wonderful deed! We could not cause the smallest blade of grass to grow or even make the tiniest grain of sand, but God made heaven and earth—this entire world, out of nothing. We would learn about this today—how God made heaven and earth and everything that is in it.

\emph{(b) Aim: How the Saviour once helped His disciples on Lake Gennesaret.}

Who of you has ever seen a big lake? Well, you have certainly seen a small lake or a pond. The Lake of Gennesaret is eleven or
\quellseite{572}
twelve miles long (that is about as far as from here to . . . .), and about six miles wide (about as far as . . .? That’s a lot of water, is it not? And it is very deep; and it has a wonderful deep blue color. If one wished to go from one shore to the other, one had to take a boat. Who of you has ever had a boat-ride? Tell the others about it. At the time of which our lesson speaks there were no steamships or gasoline boats, but only sailing vessels. A large cloth was raised on a mast (i. e., a large pole), and the wind blowing against the cloth moved the boat. Picture of a boat on the Lake of Gennesaret. Today we shall hear how Jesus was once on the Lake of Gennesaret, and helped His disciples out of a very great trouble.—Fankhauser, page 55.

\emph{(c) Aim: How the Lord Jesus was brought to a child sick unto death.}

A child “sick unto death”—what does that mean? You have seen a sick child. But have you ever seen a child “sick unto death”? It is a child so sick that he must die. Is that not terrible? I have known children who were so sick that they could not get well, but had to die. Have you ever seen or heard of such? What was the matter? “One had scarlet fever.” What did the poor parents do? (Sent for the doctor.) What did the doctor do? (Watched the child, examined him, and wrote a prescription.) Why did he do all that? Would that make the child better? (Sometimes—but not always.) The poor parents! Why do I call the parents poor? (They are in great trouble and very sad.) What does the mother do? (She sits at the child’s bedside.) And when the father comes home, he, too, is worried, and softly asks mother whether the child feels better. Then father and mother look sadly at each other; and as the child grows worse, they think more of the only one who can help when even the doctor’s skill is helpless. Who is that? And what will father and mother then do? (They will fold their hands and pray to the dear Lord.) “Let our dear child live, make him well”—so will they pray. They want to bring God, Jesus Christ the Saviour to their child that is sick unto death. Today we would hear how a father brought the Lord Jesus to the child which was at death’s door.—H. Spanauth, \emph{Praeparationen fuer den Religionsunterricht. Unterstufe}. ²1910 (according to the method of “construction”).

\emph{(d) Aim: How Jesus raised His friend Lazarus from the dead.}

Two weeks ago a member of our congregation died, whom all of you knew. Who was it? Body and soul were separated from each other, and that is what we call death. Where do the souls of pious
\quellseite{573}
persons go? Where is the body placed? Yes, it is placed in a coffin and then taken to the church where the person when alive so often worshipped; after that it is taken to the cemetery. What is then done? Oh, that is a small, narrow room! It makes no difference how rich or how poor, how high or how low, all—all must go to the grave. Death is like a heartless reaper who with his keen sickle mows down everything in his way. He lays all into the grave. How do we make graves? In eastern lands caves were used as burial places; sometimes these were natural caves, formed by nature, but, more often, they were cut out of the solid rock. The dead were laid into these underground rooms. Grandfather, father and son, besides many others were buried there side by side. Who once upon a time bought such a family tomb? Who was the first one whom Abraham buried there? Who was the next? And the next, etc.? We learned a short while ago of another man who bought such a tomb for himself and his family. Whom do I mean? Who was the first to be laid in this cave? Among us the dead are laid in the earth, and covered with earth; that is why we use a casket with a cover or lid. What is done with the lid before the casket is lowered into the grave? Of course, we have another protection for the casket, which likewise has a lid. Why? We do not wish to bury a human being like a beast, which one simply throws into a hole and covers with soil. We do everything possible to keep the body as long as we can, especially if the departed was near and dear to us. In the East no coffins were used. The corpse was wrapped in a large white cloth with sweet smelling spices, in order that any odor from the corpse might be killed; the body was placed upon a bier, or a stretcher, and so carried to the tomb. What was done in order that wild beasts might not find their way to the body? Where is it plainly stated that a great stone was rolled in front of the cave? We also do something that no animals get into the place where our dead lie. What is it? We put a fence around our cemeteries that our departed ones may rest in peace, undisturbed. What should you therefore not do when you go into the cemetery? It is a place of rest and peace, which should not be disturbed by loud talking or noise. When you go to the cemetery, then think of your grandfather or grandmother, of father or mother, or of others who lie there. Their bodies sleep beneath the sod and wait for the great resurrection day. What do we believe as Christians? Do we know when that shall happen? But this we know, that Christ will some time come again and all the dead shall arise. Then our graveyards will be as full of life as they now are full of the stillness of death. Can the Lord Jesus really call the dead to life? He can do so because He is almighty.

\quellseite{574}
We would hear today how Jesus raised His friend Lazarus from the dead.

\emph{(e) Aim: We learn how the most furious persecutor of the Christians is overcome.} (In abbreviated form.)

The Sanhedrin resolved to persecute the Christians. Stephen is accused by Saul and stoned to death. The Christians flee from Jerusalem, for Saul rages against them. Will their life or, at least, their freedom be in danger outside the capital city? Will their persecutors’ follow them into Samaria, Galilee, Syria, Damascus? Saul, the most violent of all their enemies, will seek them there. He thinks especially of Damascus; for that is the place where many of those have gone who believe that Jesus Christ, crucified and risen, is the only way to salvation. But how shall Saul find entrance into Damascus, and who will give him the right to work against the Christians? He has provided for that; for he has letters to the authorities of Damascus which tell who he is, whence he comes, what he wants, that he comes in the name of Sanhedrin, and that the authorities should help him against the Christians. These letters will open all the doors of Damascus to him. Thus equipped, he goes there. Will he be successful? If so, Christianity shall soon be no more. Or will God hinder and overcome the evil intentions of His enemies? He will; for His name shall be made holy and His kingdom shall come. How shall this be done? We hear how He overcomes Saul, the most furious persecutor of the Christians, on his way to Damascus.

\subsection{2. PRESENTATION}

\subsubsection{(a) The Stilling of the Storm} (narrated in sections for small children).

\emph{1. The Pleasant Voyage.}

The Lord Jesus often sailed with His disciples upon the Sea of Gennesaret, not for pleasure, but so that He could preach to the people on the other side of the sea. (Show map.) Sailing over the waters is far more pleasant than travelling by automobile or railroad. There is no jolting or bumping. Gently and smoothly one glides along. Sometimes the boat rocks gently as a cradle, and one is easily lulled to sleep. The Lord Jesus had worked hard all day. What had He done? Walked much, preached, healed the sick. When you have walked to school and have studied hard there, you are tired when evening comes and easily fall asleep. Jesus was tired, and as the boat went gently gliding through the water and the cool breeze fanned His cheeks, He soon fell asleep. He lay in the back part of the ship with His head upon a pillow. What do you do when you come into the house and see your little brother or your grandmother peacefully sleeping? You walk lightly and close the door
\quellseite{575}
gently, because you love the sleeping one and do not wish to disturb his slumbers. The disciples acted in the same way. They loved the Lord Jesus and wanted Him to have His rest. They were careful and did not talk loudly or laugh, lest they should awaken Him. Thus they sailed on through the golden sunshine. It was a delightful voyage. Now I shall briefly tell you the story just as it stands in the Bible, and you may repeat it to me. Narrate according to Matt. 8: 23. 24; Luke 8: 22. 23.

\emph{2. The Storm and the Danger.}

Now I want to tell you how the Lord Jesus helped His disciples. Do you need help when you are playing? No, but if you have fallen into the river then you need help. When things go wrong help is necessary. But how could matters have gone wrong with the disciples? This morning before Sunday school I showed you the watering trough. The water was smooth as a mirror, but when you blew upon it, it became rough with waves. On the Lake of Gennesaret the wind came suddenly and blew strongly, much worse than if all of you would blow together. Besides, the Lake of Gennesaret is many thousand times larger than a watering trough. The waves rolled high—as high as a house. What happened to the little ship? It was tossed to and fro like the nut-shell in the watering trough. It was no longer pleasant in the boat. Perhaps you have been sitting in a wagon when the horse was frightened and ran away. So the disciples became fearful. It seemed as if their boat must be upset and they must drown. They let the sail down and dipped the water out. They had to hold fast to keep from being thrown into the angry sea. Dark clouds, thunder and lightning, the raging sea,—all threatened death every minute. That was an hour of great need. Narrate according to Matt. 8: 24; Mark 4: 37; Luke 8: 23. Repetition by the children as in first part.

\emph{3. The mighty Helper.}

What do you think? I wonder that the Saviour did not wake up. But He slept peacefully. The disciples saw only the storm and the waves, and in their fear they forgot that Jesus was with them, else they would have thought: He can help us, for He has done many other wonderful things, healed the leper, the man with the dropsy and the like. But they did not think of that. They did not believe that Jesus would certainly see to it that they would not perish. Their faith was weak; they had only a little faith in the Saviour. Finally they awakened Him. What should they have said? They cried: Master, Master, we perish! The Saviour arose and rebuked them, saying: “Why are ye fearful, O ye of little faith!” If you had only believed that so long as I
\quellseite{576}
am with you the ship could not sink, you would not have been so fearful. And now the Saviour helps. How? By dipping out water? By steering to the shore? Oh, no! Listen: “Then He arose and rebuked the winds and the sea, and there was a great calm.” Could any one else have done that? No one, no teacher, no king, no emperor. Had we been there, we should have looked at him in wonder and surprise. “But the men marvelled, saying, What manner of man is this that even the winds and the sea obey Him? Narrate according to Matt. 8:25–27; Mark 4:38–41; Luke 8:24, 25, and repetition by class.—Recital and repetition by class of the entire story in the words of Scripture.—Fankhauser, page 55 ff.

\subsubsection{(b) Zaccheus the Publican}

It was a week before Easter, and the streets of Jericho were filled all the day long with people coming from the North to go to Jerusalem to celebrate the festival of Easter. Naturally, the people of the town did not stay in their houses, but went out on the street; and if they had time, they walked as far as the city gates in order to see the pilgrims and talk to them. In his booth by the city gate sat Zaccheus, the chief tax-collector of the town. During these busy days he had to be at the gate constantly and have his eyes everywhere. He was a small man, well-fed and well-dressed. The tax-collectors under him worked in silence; they did not like him at all. He was a rich man; and he was hated by the people because he made the taxes too high and loaned out money at too high rates of interest. But for some time he had acted as though he cared but little for his fine business and had to force himself to attend to it. He hung his head as if he were really upon the publicans. It was not pleasant when he was about.

Zaccheus also worked in silence, but he was not so much taken up ashamed before the Pharisees and other people who looked down with his accounts as in former times. Again and again he stopped to listen to the talk of the passers-by; and he seemed to be frightened. Why did they all mention one name that caused his heart to stand still? He heard it again: “The Prophet of Nazareth”, “Jesus the Nazarene”. Was He coming to the Easter festival, and did they expect Him at the gates of Jericho? Zaccheus was worried, for he had suffered for quite a while because of that name. A wanderer from Galilee had once come while he was dealing with a poor debtor, and had spoken to him so earnestly and scornfully: Have you never heard what the Prophet of Nazareth teaches: “What shall it profit a man if he gain the whole world and lose his own soul?” That
\quellseite{577}
word had struck deep into his soul, and since that time he had asked more than one traveller about the Prophet, so that by this time he had a good idea of Him and knew much about His teachings and miracles. So it came that, more and more, he felt himself to be hateful and filthy and mean in comparison with Jesus; he had begun to hate his calling, which made him rich through its dishonest dealings. But he could not give up his fine clothing, expensive furniture, soft bed, rich wines and luxurious meals. And besides, what would his wife and children say if they suddenly found themselves without all their luxuries and compelled to live like ordinary working people. It would simply not do:—such had been his difficulties for quite a while. The pencil trembled in his hand:—a little while ago he had ceased writing; his thoughts were far from his work. Then from out the jabbering multitude he heard once more the name of Jesus. He listened intently. Oh, if he might only see Him once—the good, the pure One! Perhaps it would help him! Then he heard a newcomer greeted with the words: “Have you heard, Matthew, that the famous Nazarene is coming?” “What—the Prophet?” Then they whispered, but Zaccheus caught a word that made him tremble. Did they not speak of the Messiah? The mocking words and laughter of another man came to him: “Messiah! Well, I tell you, I expect a different Messiah; my Messiah will wear the finest clothes and command a hundred thousand soldiers! But this Galilean prophet with His few disciples—fishermen and women, yes, even publicans and sinners!” Zaccheus listened more intently. They had said of the Messiah that His followers were publicans and sinners—publicans and sinners like he! Then he must see Him! Breathing deeply, he closed his book and left the gate, while his clerks sneered and scoffed as they watched him depart.

He passed through the crowd as quickly as possible, stopping to listen only when he heard the name of Jesus mentioned. He learned that Jesus was expected toward evening. Reaching his home, he walked restless to and fro. He felt that the greatest hour of his life had come. His wife became worried as she watched him, and finally laid her hand upon his shoulder and asked: “What troubles you?” Then he told her as well as he could what he knew of Jesus, and his desire to see Him. She trembled and folded her hands, and perhaps she imagined that this day might bring a great change into her life. After a little while, Zaccheus went out. He had clothed himself in his finest garments, not for show, but in honor of the Prophet. The
\quellseite{578}
streets were crowded, and with difficulty he made his way to the gate. There the multitude was greater. All Jericho was there, and besides, many of the Jerusalem pilgrims had remained to see the Prophet. So closely were they crowded that one might think a king was coming. Doubtfully Zaccheus looked about. He was hemmed in by the crowd and jostled by many elbows. He reached only to their shoulders. His costly clothing was already torn, and it was impossible for him to force his way through and to see anything. Should it all be in vain? Then he looked up and saw a nut-tree which stood on the side of the street and stretched its branches wide over the crowd. He was only a few steps away from it, but it caused him much effort to reach it. He looked at the slender trunk; he was not used to climbing; but, perhaps, he could make it. Quickly he gathered his garments about him, grasped the trunk of the tree, and began to climb.

The onlookers laughed loudly at the sight of the fat little man with the gold bracelets tearing his shimmering yellow and red clothes on the twigs of the nut-tree. One of them recognized him. “Look at the fat publican Zaccheus,” he cried, and the crowd joined with him in his jeering laughter. With flushed and sweating face Zaccheus reached a place from where he could see. His long hair clung to his face, his clothing was dirty and torn; his white hands scratched and bleeding. The people mocked him, but he cared not, for he had reached his goal and could see the One whom he so much desired to see. From the distance came a murmur which grew to a subdued roar: “He comes! He comes” Zaccheus quivered with excitement. He pushed back his long hair from his forehead, and felt how wet and dirty he was. But what mattered that! There they came—a company of Galileans, easily recognized by their dress. Three men went in advance, and the one in the middle—it must be He! Such eyes full of spirit and power, so much goodness and love in the lines about the mouth—there could not be another such in the whole wide world. He walked so modestly and humbly; He did not look about Him; and yet He moved like a king; every one had to think: If He should command, all would hasten to serve Him. They had drawn near. With a smile one of the men with Jesus touched His arm and pointed out the strange little man who looked down upon them with burning gaze. Those about Him laughed, and a fresh little boy cried: “Yes, yes, that is the rich publican Zaccheus.” Zaccheus flushed hot with shame, and his eyes grew wet with tears. Jesus looked up at him and stopped. He stood there and looked straight into the eyes of the
\quellseite{579}
trembling man. Then He said in friendly tones: “Come down quickly, Zaccheus, for I would sup with you.” Then He turned and went on. All regarded Zaccheus, but none of them laughed. They stared with wonder, envy, astonishment and anger, and many murmured: “So, He will eat with that sinner!” Zaccheus reached his home, but he did not know what had happened to him or how he had climbed down the tree or how he had made his way through the throng. Breathless he entered his house; his eyes gleamed like stars in his flushed and dirty face. His wife hardly recognized him. She regarded him with mixed fear and wonder as he embraced her and stammered: “He—is—coming.—He—is—coming!” “Who?—Who?” “The Prophet—Jesus of Nazareth is coming to us!” “Then let me go. I must get things in readiness,” cried his wife as she broke from him. Should she not prepare the finest of food and decorate her home with the most costly of her possessions? But—no! Such display suddenly seemed shameful. Should she not rather hide all those things, gained by so much wrong? She hurried about the house, while Zaccheus looked about him in bewilderment. Then some one rapped, and Jesus entered. He stepped quickly to Zaccheus and took him by the hand. Forgotten were the beauty of the room, the torn clothes, the soiled hands. Everything—everything was forgotten; for beneath the gaze of those eyes he could feel but one thing—his tremendous debt,—the great penalty which he had loaded upon himself during all these years. And now Jesus had come to him—He did not despise him: No—he saw it all clearly—Jesus had forgiven him. He felt small and poor and rich and free at the same time—so free—so free! Jesus had forgiven him! Jesus had come to him! The burden of his life fell from him like heavy chains torn in two. No, he could no longer live as in the past. That was impossible. Then he began with stammering words: “Jesus—Master—the half of all my goods I give to the poor; and if I have taken any thing from any man by false accusation, I restore him fourfold.” Then Jesus pressed his hand and said: “This day is salvation come to this house.” “Salvation,” yes, truly. That which had tormented,—had seemed so impossible to him, was now all so easy, so beautiful. He was freed from his burden, because Jesus had come to him.—From E. and O. Zurhellen, \emph{Wie erzaehlen wir den Kindern die biblischen Geschichten}, pp. 322 ff.—N. B.! We are by no means of the opinion that Zurhellen’s example should be followed without modification; He takes so much for granted at which one may only guess; and he is not always simple and chaste enough; but, on the
\quellseite{580}
other hand, one can plainly see from the example what we (p. 474) said above about the motivation, the revealing of the acting person’s thoughts, and portrayal of the origin of the deed.

\subsection{3. PENETRATION. Zaccheus the Publican} (Recapitulation, comparison, and application are here combined with penetration.)

Of whom have we spoken in this story? Of Jesus and Zaccheus. We wish to pay the chief publican two imaginary visits. First we would visit Zaccheus before he knew Christ. We go to Jericho and ask for the house of Zaccheus. All know him, and any one can show the way. His house is larger and finer than the majority of houses in the town. The floors are covered with costly rugs. Along the walls are low, soft couches. Many servants come and go. All this costs money—much money. But money is plentiful with the chief publican. In a private room stands a large chest full of gold. Where did he get all this? We shall see that presently. Those under him, the publicans, come to him and he gives them his orders. “You must get money for me,” he commands, “much money. For every hundred-weight of grain that goes to Jerusalem I want a dollar, for every bushel of raisins, twenty cents, and for every yard of purple half a dollar.” He knew well that it was from two to four times as much as the emperor demanded; but he wanted to get ever more money, and therefore he did not care how much he demanded. And when the people complained and asked: “Is the tax really so much?”, his tax gatherers were instructed to answer, “Yes, so much.” What do you think of that? It was not right. That was lying and cheating. He took the people’s money in an unrighteous manner. Deception is as wrong as theft. The publicans, however, knew that they must obey; for if they brought their master too little money, they lost their job. So they oppressed the people, for in addition to what their master demanded, they took a little for themselves. How much scolding and cursing must have been done at the counter of the publicans. The people knew that they were being robbed. But what did Zaccheus care for their scolding if he got their money! He wanted to be rich.

Matters went favorably for the chief publican, and he became richer from day to day. But how was it with his heart? Certainly, not well. When he saw the poor little traders creep by with their thin, heavily laden asses, his conscience must surely have said: “Listen! It is not right that you oppress the poor in order to enrich yourself from the little that they have.” And when the pious Pharisees went by with looks of contempt, as if to say: “Miserable publican and sinner”,
\quellseite{581}
—that also hurt. So he had no true pleasure in his riches. Perhaps he could not even sleep in peace, because those whom he had robbed walked like ghosts before him. Thus you see that riches do not make one happy. What would be far better? A good conscience. But who could help him to get that? The Scribes and Pharisees despised and scorned him. If he had gone to their schools on the Sabbath to hear God’s word, he would have been thrown out. He would have been told: “You publican and sinner, cheat and rogue, enemy of your people and friend of the Romans,—get away from here! You have no business among good people, and our God has rejected you.” The Scribes and Pharisees were proud and hard and had no mercy. Poor rich Zaccheus! This is what we see on our first visit.

After a while we pay Zaccheus another visit. How different are things now! He is not the same man. His body has not grown; for he is the same little, fat fellow. But things have changed in his heart. One can see that in his face. The lines of care are gone from his brow and the sad look has disappeared. He looks calm and peaceful. He is friendly toward his servants. He sends for all who have ever been cheated by him or by his clerks. Ere long a brown camel driver comes to him. “Have my men ever taken too much money from you?” “They certainly have,—many a time.” “How much would the entire amount be?” “Hm, I think about fifty dollars.” Then the chief publican opens his desk and counts two hundred dollars before the wondering man. Why so much? That is four times as much as has been taken. “Is that satisfactory?” The man hardly knows what has happened to him. Thus one after the other comes and goes. The chest of gold has grown considerably lighter. Better look out, Zaccheus, else all your plans to become rich will fail. But Zaccheus acts as if he did not hear.

In the evening we behold a strange gathering before the house of Zaccheus. It is composed of unfortunates of all kinds, poor, lame, blind, cripples, etc. The chief publican who a while ago would scarcely have noticed them, now goes to each with kind words and gives them liberal gifts of money. In this way he distributes half of all his wealth. But will he not again force the money into his coffers? O no! He has strongly impressed upon his assistants that they are to exact not a penny more than the law demands. And with it all he is happier than he has ever been, and his happiness is written upon his face. What has happened to cause this wonderful change? is the question his fellow townsmen are unable to answer.

\quellseite{582}
We know the answer. The Lord Jesus alone has wrought the change. We have already told how kindly the Saviour treated him when He came to his house. He, the good, holy man in his, the sinner’s house; the love that shone forth from His eyes and never a word of censure or blame! It softened the publican’s heart. He knew that he had not deserved so much kindness and love. From all that the Saviour said and did, it became clear to him that He was holy, sinless. In the presence of Jesus Zaccheus saw himself for the first time as a great sinner, and all his deception and seeking after riches loomed before his eyes as evil. “What shall I do that I may become clean?” he asks. “I will do what I can; perhaps the Lord will forgive.” And then he said what we have already learned, “The half of my goods”, etc. You have seen that he was in earnest, and it happened as he desired. Certainly the Saviour forgave him: “This day is salvation come to this house.” Now the page has been turned. To be rich is no longer the greatest thing in his eyes; but his desire is to please the Saviour, to do what He wishes, and to hate what He hates. This the Saviour wrought by His love. Tell me the Scripture passage you have learned about His love. “Oh, taste and see that the Lord is good.” Zaccheus tasted that. “The Son of Man is come to seek and to save that which was lost.” Who else of his time experienced that? The publican Matthew.—And now we hardly recognize Zaccheus. He is a different man, a new creature. Old things have passed away. What, for example is past and gone? All things are become new. What do we see in him that is new?

\begin{center}
\begin{tabular}{ll}
\multicolumn{1}{c}{\emph{The Old:}} & \multicolumn{1}{c}{\emph{The New:}}\\
Seeking after riches, & Seeking to please the Saviour,\\
Love of gold, & Love to the Saviour,\\
Cheating, & Making good,\\
Hoarding, & Distributing to the poor,\\
Unfriendliness, & Friendliness,\\
Bad conscience & Forgiveness, a good conscience,\\
Sadness, discontent & Happiness and peace\\
\emph{is past.} & \emph{have taken its place.}
\end{tabular}
\end{center}

In Zaccheus has this word of Holy Writ been fulfilled: Therefore if any man be in Christ, he is a new creature (2 Cor. 5: 17). Let us learn this passage.

You can doubtless recall another man with whom all things became new. It was Saul, or Paul. What things became old with him? Hatred against Christ, persecution of the Christians. And what became
\quellseite{583}
new? Love to the Saviour and all believers, obedience to Him, consecration and work for Him. Who caused this change? Tell me what more you know about Saul’s conversion. Thus Paul learned the truth by his own experience—the truth which he so beautifully expressed in his letter to the Corinthians.—Even now Jesus makes new creatures of all who receive Him into their heart and home. There is much in our hearts also which does not please the Saviour. Our heart is by nature like a dark, dreary room. One does not know how ugly the room is until a light is brought in. So we do not understand the condition of our hearts until Jesus reveals it to us. Then we see all kinds of evil thoughts: Envy and hatred like Cain and Joseph’s brethren, and covetousness; for we constantly desire the best for ourselves. There are hateful creatures in the room of our hearts—spiders, insects, worms. The Lord shows us this by means of His word and the Holy Ghost. A pious man once said to two bad boys: “Pray for the Holy Spirit as often as the hour strikes.” They obeyed, and ere long saw what bad boys they had been. They continued to pray, and received hearts that loved the Saviour. We can not see the Saviour and can not, like Zaccheus, give Him our hand and say: “Welcome to our house, Lord Jesus!” But He is nevertheless here and sees and hears us if we mean it. But we must be in earnest. How did Zaccheus show that he was in earnest? He not only promised Jesus that he would be different, but he also kept his word. He did not do like many children who say: “I am sorry and will not do it again,” and within an hour forget their promise. They do not mean what they say.—If there is anyone here who has secretly taken something from someone else, or at home, he should do as Zaccheus did. He should at least confess, and if possible, give back that which he took, and not steal again. Otherwise he is not in earnest. How different things become in the home of a sinner, a drunkard, for instance, when Jesus enters there! Father, mother, and children are different in spirit and in outward appearance. There is even a change in the furniture, the walls and the windows. The heathen leopard becomes a lamb when Jesus works in the heart through the preaching of the Gospel. There is a wonderful change in a country when the people become Christians. A shipwrecked sailor was cast upon an island where cannibals used to live. In fear and trembling he climbed a hill and saw—a chapel. Then his fear fled, for he knew that where the cross is planted, he need not fear. Behold, I make all things new!—Fankhauser, page 141 ff. Naturally, the conversational form should here be made prominent.

\quellseite{584}
\subsection{4. APPLICATION. The anointing of David} Because God looks upon the heart, we can do nothing wiser than to strive for a clean, beautiful heart that pleases Him. It pleases Him best when we, like David, think much of Him and love Him. If you do not know how to do this, pray to the dear Lord. He will gladly give you that which He gave David—the Holy Spirit. He will teach you to love God and the Saviour aright. He will make you pleasing to God.—Certainly now and then evil thoughts will slink into your heart,—a hateful word, hateful thoughts against your companions, a foolish trick. But the Spirit of God tells your conscience kindly and gently that you did wrong. Would you remain pleasing to God? Then follow David’s example. How? Sing a hymn, repeat the passages of Scripture you have learned, or think of David the shepherd boy, or of Joseph. What did Joseph do when he was tempted to wrong? How did Joshua admonish his people? Therefore keep watch over your hearts, that you may please the Lord, your God.—You liked David’s courage, boys, did you not? I will tell you a story of a boy who had the right kind of courage. Story—The Proper Courage. Little Edward lacked courage to do wrong, to lie, to steal; but he was brave enough to save a child’s life at the risk of his own.—Above all, we want to equal David in one point more. Do not feel proud if you have been good and obedient, and do not boast how much you read your Bible or pray: only think, it is enough that the good Lord knows it. Many small and large children are like hens. If they have given a poor child some little thing, or if, for a change, they have obeyed their parents, or perhaps have not talked in school, they cackle loudly about it like a hen when she has laid an egg; or they look around to see if some one will not praise them. Do not do so. Rather think of the dear Lord. Rejoice if you can please Him. He sees it, and at the proper time it will be revealed.

The Lord no longer needs kings—yet He will make of us something finer than kings. He watches to see whom He can make His child and take into heaven, where each one will have far more glory than David ever possessed. But He will first conduct an examination. Each one must stand before Him to be tried, as the children of Jesse were tried by Samuel. But God will not only look upon us. He will look us through and through. On that great examination day all things will come to light that were hidden and secret, the good as well as the bad. Then perhaps a child will stand before Him and say: “Behold, dear Lord, here am I, Mary, the daughter of the
\quellseite{585}
rich miller. My father has much, much money, and all people say, You dear child, you pretty Mary! Dear Lord, do you not know me?” But God will sadly say: “I know you well, but I cannot use you. You have constantly thought of your pretty clothes and have forgotten Me entirely. You have kept yourself clean and nice, but you have not kept your heart and mind clear of evil thoughts. I look at the heart. Depart from Me”.—Or big John will say: “Dear Lord, I have always been the strongest in our school, and all have praised my strength and agility”. But God will perhaps say: “Yes, you have overcome others, but you have not conquered your pride and anger and the evil in your heart. Your heart is not good and you have never prayed to Me to cleanse it. Depart from Me!” Then will come clever Lydia, and say: “Dear Lord, do you not know that I have always been the first in my class? You will certainly give me first place in heaven, will you not?” But God will say: “Lydia, I not only take into account the Scripture passages and Bible stories they all can tell, but also this, that they practice what they have learned. But you have not obeyed your parents, and much less have you obeyed Me.”—Perhaps a small, shy boy will come and stand silently before the Lord. He was not strong, and the others often teased him. He was poor and had no shining shoes or warm cap, and was often cold in winter. But the good Lord will regard him kindly and say: “Come to Me, dear child. I have seen how you have thought of Me and loved Me. I know well how often you have shared your bread with other poor children and how you have overcome the temptation to lie or steal because you did not wish to make Me sad. I know you. You are My dear child. Enter into the eternal joy which I have prepared for you”—See how differently God looks at things! How easily could rich Mary, big John and clever Lydia have pleased God, if only they had kept watch over their hearts! Despise not the poor child in torn clothes. Perhaps God will be pleased with his heart rather than with yours. The Lord looks upon the heart. “Search me, O God, and know my heart: try me and know my thoughts: and see if there be any wicked way in me, and lead me in the way everlasting.”—Fankhauser, page 130 ff. The catechist would do better to confine himself to what Fankhauser says of the examination.

\subsection[5. A complete example]{5. A complete example: How God made the first people happy}

\subsubsection{Preparation} From whom does everything that we see on earth come?—Correct. It all comes from the dear Lord who made it in six days and still preserves it. Therefore we call God the Creator of all
\quellseite{586}
things and say on Sunday in church: “I believe in God the Father, Maker of heaven and earth.” But for whom has God made this world with all its beauty? What do you think, Anna, for whom has God made this earth with its mountains and valleys, rivers and seas, fields and meadows, birds and animals? Indeed, God made it all for us; we should live here; it should be our home and give us all that is necessary for our life. Therefore we should be glad and happy. How rich is our God, who made everything so beautiful; and how good and kind He is for calling all this into being for our sakes. And yet we have not heard the best that God made in the beginning, in order that we might be perfectly happy. What could that have been? We shall hear about it today, for we shall hear how God made the first people happy.

\subsubsection{Presentation} When God created the world, He made a most beautiful garden. You have perhaps seen pretty gardens, but none so beautiful as that. It was called Paradise, that is, Garden of Pleasure, because it was a pleasure to be in it and to live in it. There were beautiful trees of all kinds—forest trees which gave cool shade and fruit trees which yielded the finest of fruits. And Oh, the many, many flowers that grew there! They bloomed and gave out perfume as flowers never since have done. There were also all kinds of animals, large and small, from the little mouse to the lion, the king of beasts. Happy birds fluttered from twig to twig or flew high into the air and sang until the heavens resounded. There was plenty of water, for there were many springs, clear as crystal, flowing from the rocks. These flowed down the mountain-side and became brooks, and the brooks grew into rivers and streams. Four such streams flowed through Paradise and watered it, keeping the grass and trees and flowers green and fresh. Rich people, kings and emperors have made themselves fine gardens, but believe me, children, none of them, however large and beautiful, could compare with the pleasure garden which God made when He created the world. And think of it!—Adam was permitted to live in this beautiful garden of Paradise. Oh, how happy did God make our first parent!

He could go through the garden, rejoice over the flowers, eat the fruit and lie in the cool shadow of the trees. He could see all the animals and play with them. But with all this, something was yet lacking. What was it? He had no one to speak with. The animals followed him when he called them by the names he gave them and let him stroke them. The birds came and perched upon his shoulders
\quellseite{587}
or rested on his hand. But when Adam spoke to them they could not reply and they understood nothing of his thoughts. Then Adam wished for some one who would be his equal, with whom he could speak, and who could be a companion in his joy; for his heart was full to overflowing when he beheld all the beauty that surrounded him. The good Lord, understanding Adam’s wishes, let him fall into a deep sleep. Then He took one of Adam’s ribs and out of it made a woman, Eve. When Adam awoke, he saw her stand before him. He beheld her with wonder and shouted with joy: “This is now bone of my bones and flesh of my flesh: she is a fitting companion and just what I have wished for. I can speak to her and she to me. I can tell her all my thoughts.” How happy was Adam then! He went about the garden with Eve and showed her a pretty spot here and another there. They sat beside the brook and listened to the babbling water, or enjoyed the sweet songs of the birds. And when they became hungry, they plucked berries or apples and pears from the trees. But many a time they must have stood still and said to each other: “Oh, how much our God must love us since He has made everything so beautiful and placed us in this lovely garden.” They worked, too, for God had told them to take care of the garden. They watered the tiny plants, making little ditches, so that the water could flow to them; they lifted up a branch of the grape-vine here one and there one and fastened them so that they could climb up the tree. But their work was not hard, it did not make them tired. It was a pleasure and a joy and only increased their happiness.

God had made Adam and Eve masters not only over the garden, but also over all animals that were in it, the small as well as the great. “Be fruitful and multiply,” He had said, “and replenish the earth, and subdue it; and have dominion over the fish of the sea, and over the fowl of the air, and over every living thing that moveth upon the earth”. In our picture the artist has shown us the Lord God pointing to the animals with one hand, while He raises the other before Adam and Eve and speaks these words to them. See, here is a mighty elephant with his big ears and long trunk. Above him sit two mighty birds. Farther down is a proud horse with high head and flowing mane. And near by—what is that just showing its little head?—Sure, it is a lamb. And that other is a deer, which can run so swiftly and flees when it sees human beings. What comes next? Correct—that is a lion. The lion is a powerful animal, and therefore he is called the king of beasts. Beside him a cow is peacefully lying,
\quellseite{588}
while between her legs a serpent crawls. Could we today put all the animals together as the painter has done? Why not? Certainly. They would kill and eat each other and the people, too. We would have to put some of them behind iron bars. But then it was different, for when God created them the animals were not enemies to each other. It was easy, too, for Adam and Eve to be their masters, and it was a part of their happiness when God said to them: “You shall rule over all these creatures and over the whole earth and subdue them”. So God made the first people kings over the earth.

But, my dear children, I have not yet mentioned the best and greatest happiness of Adam and Eve. What could that be? It was when the Lord God came to them in the garden and spoke with them as a father with his children. Oh, how they rejoiced when He entered the garden, and how they hastened to meet Him, as children do when their parents return from a journey. How bright their eyes!—They had so much to show Him and so many questions to ask. Today we can not see God, but then He came to them in the form of a human being, so that they could see Him and feel Him. What delightful conversations those must have been between God and them! He would explain to them what they did not understand, how He had made the world, that it was all for them for their enjoyment and happiness and that they should rule over it. “Yes,” said God to them one day, “everything you see is yours and you may eat of every tree in the garden, but of this tree here, the tree of knowledge of good and evil, you shall not eat. For the day you eat thereof, you shall surely die.” What a kind God, who, with this one exception, placed everything into their hands. How He must have loved these first people, and how happy they must have been!

\subsubsection{Penetration} (association, comparison, and abstraction). Where did the first people live? Where was Paradise? What—you do not know? Neither do I. Many learned men have thought that it was in Asia Minor, between the Euphrates and Tigris rivers. (Show map.) They may be right, but no one knows for sure, and we can not seek it, for it is no longer upon earth. Since when has it disappeared from the earth?—At any rate, since our first parents fell into sin there has been no Paradise here below. Much more important is another question, namely, how Paradise looked and what it was like. And to this the name “Paradise” gives the answer. What is the meaning of the word Paradise? Who has
\quellseite{589}
noted it? Paradise, or Garden of Eden, means garden of pleasure. There are pleasure gardens today, for many a rich man has one made for himself and many cities have parks which could be called pleasure gardens. But why call them pleasure gardens? Because the people find pleasure and enjoyment in them. When they become weary, they can rest there. When the sun is hot, what can they find there? Especially who can play and have a good time there? You children, of course. But all the pretty parks and gardens, with which one can they not compare? Now tell me all that was to be found in the pleasure garden called Paradise. What was there besides the beautiful flowers and the trees? What can you tell about the streams that flowed through and watered the garden? Because of this, what was the condition of the grass and all plants and trees in this garden? It is often so different now. What happens so often in summer to the grass and plants? Yes, sometimes the trees lose their leaves in summer when it has been hot and dry and no rain or dew falls. It was not so in Paradise. Everything was as fresh and green as in spring-time. And just think of it, children, there was no hot summer to make everything dry; no autumn when the leaves fall from the trees; and no winter, when all nature seems dead and snow and ice cover all the earth. To be more correct, these four seasons existed also in Paradise, but none of them brought those disadvantages they bring now. Oh, then, what pleasure and joy to live in that garden! Therefore, how do we rightly call it? And whom did God place in that garden to live? What should Adam according to God’s will have there? Surely, pleasure, satisfaction and joy. Truly happy did God want to make Adam, therefore He placed him into this garden. Let us all say this together: \emph{God wanted to make the first human being happy, therefore He placed him into Paradise}. (Write on blackboard.)

Was Adam happy in Paradise? Tell us about his happiness. What, besides the trees and the flowers gave him the most pleasure? To be sure, the birds and the other animals; they had life and could look at him. Tell me how trustful they were toward him. What did he give to each bird and animal? Adam had closely noted their way of living, everything was so new to him, and God had given him keen understanding. If we now knew the names which Adam gave to the animals we would certainly wonder at their fitness. The animals in time learned the names Adam gave them.
\quellseite{590}
How is this shown? But when the dog came at Adam’s call, and the lion lay at his feet, so that Adam could fondle his mane, and when he spoke to them and called them pet names, what could they not do? How do we feel without some one to whom we can speak? Yes, we feel lonesome and forsaken. Just fancy yourself alone, with no one to speak to. How unhappy you would be. How would Adam therefore have felt in spite of the beautiful garden in which he lived? He would certainly have been without something very necessary to his happiness. Man is happy only when he can express himself and finds that his thoughts are understood and answered. The reason why it is so hard in prison is not, because one must work, but because one must be alone and dare not speak to others. What, then, did Adam soon wish for? To whom was his wish satisfactory? What did God say to Himself? What did God immediately do with Adam’s wish? Why would He act upon it immediately? Certainly, He wanted Adam to be happy. Could He not have given Eve to him at the very beginning? He could have done so, but why did He not do it? If no one knows, I shall have to tell you. The Lord wanted Adam to learn what he yet needed. If we have something from the beginning, then we do not yearn for it or value it highly. On the other hand, if we receive something which we need and for which we have long wished, we are much happier when we get it. So it should be with Adam. Tell me how God gave him some one to whom he could speak. How did Adam feel when he awoke from his sleep and suddenly saw Eve before him? How did he express his wonder? He said with these words that here at last was a suitable companion. Where, so far, had he found no suitable companion? It was immediately clear to him that man and woman belong together. With what wonderful words did he show this? What do we say of a man and woman when they leave father and mother and live together? And by what name do we call that life when a man and woman are united? You do not know it, and therefore I will tell you: They are married, and such living together we call married life. Who gave Adam his wife? God Himself thereby instituted the marriage relation and showed that it is pleasing to Him when man and woman enter into the same. If God Himself instituted marriage, what should husband or wife not do?—What is done with the marriage relation when husband or wife separate or do not love each other and are not true to each other?—Yes, they break the marriage tie,
\quellseite{591}
for they break the word, or promise, they have given to each other. That is a great sin. Tell me the Sixth Commandment.—But we must return to our story! Did God make Adam happy when He gave Eve to him? And now Adam had some one with whom he could do what? And what could Eve do when Adam spoke to her? Tell me in a few words how happy they were. Adam and Eve could then say what King David later declared. When he thought of all that God had given him to make him happy, he rejoiced in these words: “The Lord is my shepherd; I shall not want”, etc. Did Adam and Eve have to work? But why did this work not disturb their happiness? Correct. Joyful work is a pleasure and no burden. Now, who gave Eve to Adam, thereby making him happy? Let us repeat together: \emph{God made Adam happy, because He gave Eve to him, with whom he could speak and rejoice}. (Write this on the blackboard.) Mary, you may tell us this. Now Frank may repeat it. John, what was the first thing we heard that God did to make Adam happy? Yes, God made Adam happy when He placed him in the pleasure garden. Louise, what did we hear further?

However, we know still more about the happiness of the first people. God not only placed Adam in Paradise, gave him Eve to be his playmate and companion, but also made them rulers over the animals in the garden and over the whole earth. Who can repeat the words with which God gave them this right? Yes, those are the words. Many people should come from Adam and Eve, so that the entire earth should be filled. And wherever these people should go, they should rule over all the animals, and all the powers of the earth should serve them. Tell me, Louis, how do you think the people should make the beasts serve them? How do we today compel the horse to serve us? How does the cow serve us? The chickens? So should all animals, even the lion and the elephant, serve Adam and Eve, each one according to his powers. I have also spoken of the powers of nature. What had I in mind? Is there none who knows? I was thinking of the air, the wind, gas, electricity. Are not those great powers with which man accomplishes wonders? How do we today see that we control the air? How do we use the wind? For what purposes do we use gas and electricity? And now think of the treasures beneath the earth. Can you mention some of them? Certainly, coal, iron, and the other metals, silver and gold. The good Lord set Adam and Eve
\quellseite{592}
over all these. They should use them all for their comfort, good, and pleasure. How rich God made them! The lords and kings of the earth also have control over many animals, powers and hidden treasures, but they rule over only one land, while Adam and Eve were rulers over the entire earth. What title then can we rightly give them? Perfectly correct. We can call them kings, for they were kings and rulers over the whole earth, with all that was in or upon it. People often follow an earthly king because they must and because they will be punished if they do not. How did the animals and the powers of nature obey Adam? And what made it easier for him to control the animals? Recall the picture you saw the last time. Which beasts did we see peacefully lying side by side? Could that be done today? Why not? What must be done with the lion or the serpent? They would not only hurt the other animals, but whom also would they attack? Yes, man's life would not be secure. Then what was not present among the animals of that time? And between whom also was there no enmity? That is right. As the flowers, large and small, did not wither or fade, so was there no enmity, no death among the animals. They did not prey on each other, but all lived together in peace and harmony. It was good to be king and ruler there. Who made the first people rulers over the whole earth? How did He want to make them? Yes, He desired to make them happy. Let us say together: \emph{God made the first people happy, because He made them rulers over the whole earth}. (Write on the blackboard.) In how many ways have we thus far learned that God made the first people happy? Tell me the three ways.

Was there anything more that God could do to make the first people happy? One might think that they had all that was necessary to make them happy, and yet the most important thing was lacking. Picture yourselves living in the most beautiful garden there ever was, where you had brothers and sisters to speak and play with; where you were rich and ruled over all things, but where you had no father or mother. Would you be really and truly happy? And how would it be if father and mother were living, but they were angry with you and wanted to have nothing to do with you? How would you then feel? Certainly you would be very unhappy. Of whom would you be constantly thinking? And of what especially would you be thinking? If you were good children, all joy would depart if you had to say: Father and mother
\quellseite{593}
want to have nothing to do with us. They are angry with us. Now, who was Adam’s father? Why was God Adam’s and Eve’s father? How would the first people certainly have been, if their heavenly father, the dear Lord, had never come to them and wanted to have nothing to do with them? What, then, would God have done to make the first people perfectly happy? Correct. God came to them in the garden and spoke to them as a father with his children. Could they see God? Yes, children, I believe they could see Him as we some time shall behold God in heaven. What did God have to do, if Adam and Eve were to see Him with their natural eyes? Of course, God must make Himself visible. God is a spirit, but He can take on human form. And what happiness must have been theirs when Adam and Eve saw God come to them in human form. What do you think they did when they saw Him coming to them? And what do you think they did when He walked through the garden with them? Tell me a few questions which they probably asked Him? There were many wonderful things in this beautiful garden which they could not explain. Which was the first question, in your opinion, they asked? Yes, I too, think that they asked the Lord who made the beautiful garden, the whole earth, and the heavens above. O children, those must have been blessed hours when God spake with them and they told Him all that filled their hearts. Now what could not but be impressed upon them as they talked with God? It was out of love to them that God made everything so beautiful; out of love He came to them and communed with them like a father with his children. And I should greatly wonder if Adam and Eve, when they were alone, did not often say to each other: How dearly God must love us! Oh, we can not do otherwise than love Him in return. Of what Scripture passage does that remind you? Yes, so would they have said to each other: “We love Him, because He first loved us”. Now we know what made Adam and Eve happiest. What was it? Let us repeat together: \emph{God made the first people happy, because He came to them and walked with them as a father with his children}. (Write this on the blackboard.) What was the first thing that God did to make Adam and Eve happy? What was the second, the third, the fourth? Yes, God had done so much for them, and it is no wonder that they were happy, perfectly happy. And why did He do so? Right. He did all this purely out of love to them.

\quellseite{594}
And notice, children, this also, that it was because of His love that He gave them only one commandment. He did not wish to plague them with many commandments. They should not constantly have to say: Yes, God has forbidden this and that. What is the commandment God gave them? \emph{You may tell me—and you—and you}. Why did God call this tree the tree of the knowledge of good and evil? Do you not know? Well, the word “knowledge” comes from the word “know”. This tree should help them to know what was good and what was evil. Does it make any difference whether we know what is good and what is evil? Yes, indeed, much depends upon it. You only need to compare yourself with father and mother to see it. They know much better than you what is good and what is bad. They have often enough experienced how bad it is when one does wrong, and how blessed it is when one does right. Hence they do not easily permit themselves to be persuaded to do wrong. And was it not necessary for Adam and Eve also to learn what was good and what was evil, so that they might always choose the good? And what was to be the result if they did the good? Correct. Their happiness was to increase. Are you happy and joyful when you have done wrong? How do you feel? Yes, you feel ashamed and uneasy. It should not be so with Adam and Eve. \emph{They should learn most precisely what was good and what was evil, in order that they might always do good and be happy and blessed}. (Write on the blackboard.) That was what God wanted, and therefore He placed the tree with the forbidden fruit in the garden. What, therefore, were God’s intentions in regard to them? His intentions were good because He loved them. Now one should have expected that Adam and Eve would have done what? Yes, they should have loved God. They said many times that they loved God; however, they should not only say it, but also prove it by their deeds. How could they have best proved that they loved God? Perfectly correct. That is the best proof of our love to God, that we keep His commandments. In which Scripture passage did you learn that? We shall see if Adam and Eve loved God so well that they kept His commandments. That would surely have been the best thanksgiving to God for having made them so happy.

\subsubsection{Application} Children, does God love you and has He made you happy? At any rate, where has He not placed you, that you might live there? Why would that not do? But is not this world in which
\quellseite{595}
you are permitted to live a beautiful one? When do you notice most that this world is beautiful and glorious? In spring-time, of course, when everything blossoms and grows. It is sometimes so lovely that we hear the people say: “It is like heaven.” And as Paradise once was the most beautiful spot on the earth, so, I think, God has also assigned to you a place upon this wide earth where you have it best. Which place do I mean? I am thinking of your home. That is for good children the best place. Ask the grown-ups, who have been among strangers for a while, or those who have no home, if there is a better place on earth than home. It is true, one can be richer and freer elsewhere; but, as a rule, it is nowhere better than where father and mother are. Think of the prodigal son. What did he think at first? What did he learn afterwards? And after what did he finally yearn with his whole heart? Yes, father’s hand holds closest, mother’s eyes beam kindliest, and mother’s heart understands best. Now, who placed you in this beautiful world and gave you father and mother and home? What state of mind should this work in you? Did He leave you without any one to speak to? You certainly would be unhappy children if you had no one with whom you could talk, for it is very difficult for you to sit quiet and silent for but an hour. But whom did God give you to speak to at your pleasure? And who is it that always listens to you and never becomes weary of your prattle? See there, how happy God has made you! and when you become older, you will more and more understand that something of that kingly position which God once bestowed upon Adam and Eve, still remains for you. And if you had no home, no father, no mother, no brothers and sisters or friends, why would you still not be forsaken? Yes, indeed, God in heaven would still be your father and you would be His children; and that is the greatest happiness that one can have. Through what has God become your father and you His children? In which passage of Holy Writ have you learned this? Correct. The great thing in Baptism is that it makes us God’s children. What does a true father do for his child? What does God therefore do for you, because He is your father and you are His children? How does He care for you? In what other ways? And still others? We would not forget the most important thing. The dear Lord sees to it that you hear His word and that therein you are shown the way to true happiness. He therefore provides not only for your body, but also for what? He would make you
\quellseite{596}
happy in body and soul. I know a hymn in which it is beautifully told how good we have it with the good God, with our dear Saviour Jesus Christ. All of you know it; we have sung it often together. It speaks of the Lord Jesus as the shepherd and of us as His lambs. Now you know it; repeat it, John. But the finest thing is that you as baptized children can also speak with God like children speak with their father. How may you call God? Who has taught us that we can call God by the name—father? Where? What are the opening words of the prayer Jesus taught us? How does our Catechism explain these words? God would give us encouragement and pleasure in praying. That is why He tells us especially to call Him father, for His words are a tender invitation to pray to Him. How should we pray to Him? How do children ask their father? Yes, they are not timid and fearful, but what confidence do they place in him? Now is it not something unspeakably great to be able to pray to God in heaven with all confidence and assurance as dear children ask of their dear father? We must call that true happiness. Who has given you the right to speak to God at all times and tell Him all you think and need? How has God then also made you? Let us say it together. \emph{God has made also us happy, for He has made us His children through Baptism}, who can entreat Him as dear children entreat their dear father. And, children, whatever may be lacking to make our happiness complete, He will give us above in heaven; for if we die His children, He will take us home into the heavenly paradise. There our happiness will be full. The last stanza of a certain hymn learned by us recently, fits in here. Will you say it? How must God be disposed toward you if He showers so many blessings upon you? How should you feel in your hearts since God has made you so happy? What kind of children would you be, if you did not feel so? Yes, indeed, you would be ungrateful to God, you would be thankless, evil children. And how must you show that you love God? Yes, that is love to God when we keep His commandments; and the better we keep them the greater will our happiness be.

\subsection[6. Presentation of the Bible History]{6. Presentation of the Bible History in the Sunday School}

(Here the most important parts of the penetration and application must be combined with presentation):

\subsubsection{How the first people sinned}

\emph{Preparation}. How happy God made the first people! He let Adam live in Paradise, and Paradise was the most beautiful garden the earth has ever seen. Here Adam could have all that his heart
\quellseite{597}
desired. Here God created Eve for him, that he might have some one with whom he could speak and share his joy. Oh, how happy both of them were when they wandered through the garden hand in hand and called each other’s attention to the beautiful nooks they had found. I wish that I, too, might have been there. Then, besides, they were masters over all that was in the garden—over the beautiful flowers, which they might pluck as they wished; over all the sweet fruits, which they might eat as they desired; over the birds of the air, which fluttered about them and showed no sign of fear; over all the beasts of the field, even the mighty lions and giant elephants, which obeyed their slightest word. Yes, God made Adam and Eve rulers over the entire earth. Indeed, they were happy people. And yet, you know that their greatest happiness consisted in something higher, much higher. It consisted in this that God came to them in the garden and spoke with them as a father with his children. What blessed hours those must have been when they so freely and joyfully associated with God as you children do with your parents, and when without fear of the almighty and holy One they could look into His eyes and walk at His side! He did not burden them with many commands, so that they had to watch every step, lest they transgress one of His commandments. He gave them a single command, and that was for their good. What was that commandment, Frank? That is correct. They might eat of all the trees in the garden save one. It should be both church and school for them, for from it they should learn what was good and evil, that they might always do the good and thereby grow in happiness and blessedness. For what was God’s purpose in all that He did for them? That they might be truly happy.

What should these first people have done with this command of God, Mary? Yes, they should have gladly kept it. And if they had kept it, their happiness would never have ended—it would have lasted for ever. But they did not keep it. Today we shall learn how the first people disobeyed God’s command, how they sinned.

\emph{Presentation}. How may it have happened that the first people transgressed God’s command and sinned against Him when He had nothing in mind but their happiness? One would think that it could not have been possible for them to sin against God, so well had He arranged all for them. At first, Adam and Eve did keep God’s command. How long that lasted, we do not know; but it
\quellseite{598}
could not have been such a very short time. More than once they thought of the command and said to one another: Of that tree yonder we may not eat, lest we die. I think that when they looked at the forbidden tree, a sort of horror stole over them. I should not wonder but that they came after a while to call the tree “the tree of death” and passed it at a distance rather than go near and touch its fruit. They evidently thought not only of God’s command and the death penalty that would follow disobedience, but also of His love. And they probably said to one another: Since God loves us so, and has made us so happy, we will love Him in return and try our best to keep His command. But how did it come that they finally transgressed and ate of the forbidden tree?

The evil thought did not spring from their hearts. But there was one who could not bear to see them love God and obey His command. That was the devil. He is the great enemy and adversary of God who is always striving to destroy God’s works. As often as he saw these people, so pure and holy and obedient to God’s will, he became angry and planned how to destroy this beautiful work of God, rob the pair of their happiness, make them impure and unholy,—yes, as evil as he himself was. There are boys who, when they are dirty, like nothing better than to make others unclean also. It is just so with the devil. He wanted to make these first people as evil and unclean as himself, so that God could no longer use them, could no longer be their Father, but would have to cast them out of Paradise. He said to himself: If I have once torn them from God, then they are mine and will serve me all their lives; then I shall always be their master and they shall ever be my slaves.

The devil evidently considered long how he might accomplish his purpose. He asked himself: How shall I begin? How can I lead Adam and Eve astray, so that they disobey God and become mine forever? It was clear to him that he could gain his end more easily and surely if he tried Eve instead of Adam; for the woman, being weaker than the man, could more easily be led astray. He also told himself: As I am, I may not show myself. I dare in no way betray that I am the enemy of God and their enemy. I must appear to be their friend who is seeking their best interests. Finally his plan was perfected. The devil knew what he would do.

Craftily and cunningly he began. How craftily and cunningly, the Bible tells us at length. In Paradise there were all kinds of
\quellseite{599}
serpents, large and small. They glided swiftly upon the ground; they coiled themselves in the bushes or wound themselves about the trunks of the trees and so made their way to the branches. Adam and Eve were not afraid of them; for they could not and dared not harm them. They had frequently watched them as they played with the other animals, and wondered at their cunning, in which they surpassed the others by far. It was one of these serpents that the devil used. It should be his mouthpiece, the agent through whom he would deceive Eve and bring her to transgress God’s command.

Eve generally avoided the forbidden tree, but one day, as she drew near to it, she saw an especially pretty serpent coiled about the trunk of the tree and looking in her direction. It seemed to Eve that the serpent was looking directly and steadily at her. She stood there and gazed at the serpent and the tree. And as she looked, the serpent suddenly began to speak. Eve was filled with wonderment, for though she had often before seen the serpent, she had never heard it speak. She thought that no one could speak like she, her husband and the good Lord. Can serpents speak, Mary? That is right; they can not speak. Then how did it happen that the serpent on the forbidden tree could speak, Louise? We know the answer: it was really not the serpent, but the devil who spoke through it. Eve should not know who was speaking to her, therefore the devil spoke through the serpent. And what did he want to do? Who remembers? You are right, John. The devil wanted to deceive Eve and induce her to transgress God’s command.

And now hear what the devil said to Eve: “Did God really say, You shall not eat of every tree of the garden?” He knew well that God had said something different, namely, that God had given Adam and Eve permission to eat of all the trees with but one exception; but he wanted to find out if Eve still remembered God’s command. He wanted to make her uncertain about it. He said to himself: If she is no longer sure of what God told her, I can the more easily get her to disobey and sin against God. We already begin to see how cunningly he started to carry out his plans. But this time his trickery did not help him; for Eve knew what God had commanded, and was not in the least deceived. She replied: “We may eat of the fruit of the trees of the garden; but of the fruit of the tree which is in the midst of the garden, God hath said, Ye shall not eat of it, neither shall ye touch it, lest ye die.” She evidently meant, What kind of foolishness are you speaking; the very opposite is true. God has by no means
\quellseite{600}
said that we may eat of none of the trees of the garden. He is not so cruel that He would put us in this garden with all the many fine fruit trees, and then tell us that we should not eat of them. He is too loving and good for that. He plainly told us that we might eat of all the trees of the garden except this tree about which you have coiled yourself. The dear Lord has His good reasons for making an exception of this tree. Whoever touches the tree or eats of its fruit must die. God does not want us to die. He wants us to go on living and to remain in this beautiful garden forever. What is the matter with you, you foolish serpent? You are generally so clever and know more than all the other animals.

In this manner Eve answered the serpent and the devil who spoke through it. The devil might then have left, for he had lost his case. In spite of all his craftiness, Eve had put him to shame. She not only knew the command exactly, but she also knew that God had given it for her good. But the devil did not consider himself beaten. If he cannot accomplish his purpose one way, he tries another. So he thought quickly and formed another plan by which he could gain his end and bring Eve to disobey. Stop, he said to himself, Eve thinks that God gave the command out of love, because He is seeking her highest good. This faith I must take from her. I must induce her to think that God gave this command out of envy and not out of love—not because He means well with her, but because He does not wish her to have the best. Therefore he quickly replied to the woman: “You shall surely not die if you eat of the fruit of this tree. God has lied. He well knows that the day you eat of it, you shall become like God, knowing good and evil. God does not want you to be like Him, therefore He gave that command. He did so not out of love, but out of envy and jealousy.” What the devil told her was a big lie, and he knew that he lied. But he did not care for that, if only he could bring the woman to believe him and eat of the forbidden fruit.

Did she see the lie? She should have done so; for she knew God; she had been with Him long enough, and had often enough beheld how holy and truthful He was and how He loved her with all His heart. She should have replied to the devil: Depart from me. I will not listen to you. God is not as you say. Would He have placed Adam and me in this beautiful garden and have made us so happy if He had been envious and jealous of us and instead of wishing us well? Thus she should have answered and thereby have smitten the devil a second time and have remained God's dear child. Such thoughts evidently arose in her mind, but other thoughts were now becoming
\quellseite{601}
active. The serpent had told her something she could not forget. It was this: You shall be as God. She thought that must be something great, to be like God and to know both good and evil. And the more she thought of this the more her present happiness did not seem worthwhile. It was with her as sometimes occurs with you children. Bread and butter tastes very good to you until you see a nice, fresh cake on the table; then you suddenly become dissatisfied with the bread and you want the cake. Or, how satisfied you are to be in the house with mother, until you see the children play outside! Then the room suddenly becomes too narrow and small, and you think: We should be really happy if we could play outside. It was something like that in the case of Eve. Paradise was suddenly not good enough. She wanted more—to be like God, and like Him to know good and evil. She forgot entirely to ask if Satan was telling the truth or if he was lying again. She forgot that God is certainly not envious and that He cannot lie. She heard and thought only one thing: to be like God. Her whole heart was suddenly full of that desire. But that was an evil desire; it came into her heart from the devil and not from God. But she did not ask about good and evil, she yearned after but one thing—to be like God. And then, children, Eve lifted up her eyes and looked at the fruit of the forbidden tree. Oh, if she had only not looked, she might still have considered the right and not have given her will over to evil desire. But she looked up and examined the fruit on all sides, and because it looked so good and sweet and promised to taste so good; because she believed the devil’s word that it would make her like God, think, O children! she finally stretched forth her hand, took the forbidden fruit, ate of it, and gave of it to her husband. O children, the angels in heaven must certainly have wept, and it must have cut God to the heart; for now the first people had believed the devil more than God! Now the command was transgressed and the first sin committed upon earth! Now were come upon earth death and all manner of misfortune. “By one man sin entered into the world, and death by sin; and so death passed upon all men, for that all have sinned.” The devil was the cause, for he deceived the first people; but they were not blameless, because they believed the devil more than God. If they had resisted the devil, he would have fled, but they listened to him and did his will. Therefore: Be not deceived! Evil communications corrupt good manners. Resist the devil and he will flee from you. Children, when evil companions entice you, follow them not. There are many sad days, but the saddest of all, as dark as midnight, is that day on which the first people sinned.

\quellseite{602}
\subsection[7. Review and drill]{7. Review and drill at the end of a longer section:}

\emph{At the conclusion of the history of Abraham’s life. A.) Bible story: 1. Headings of main and subordinate divisions of history.—2. Character sketch of Abraham:} a) toward God: pious, god-fearing, obedient, believing; he fears, loves, trusts God above all things; is grateful and humble. b) toward men: peaceable, unselfish, merciful, helpful, courageous, valiant, noble, hospitable, polite, friendly; an interceder, a faithful husband, a providing father, a good, kind master to his servants, a faithful head of the home who teaches God’s word to all his household. \emph{3. The promises given him}, which were already fulfilled; the three great promises, great nation, land, blessing to all people.—\emph{4. Character sketch of Eliezer.—5. Character sketch of Sarah and Rebekah.—B) Doctrine}: \emph{1. God’s essence and attributes}: almighty, truthful, faithful, long-suffering, patient, merciful, but also righteous, jealous, holy, impatient of evil. \emph{2. Human Virtues and Weaknesses}.—Which virtues did Abraham exhibit toward God and men? Which Eliezer, Rebekah? Which did Lot show; which Sarah? On what occasion did they show them? What should a master, the head of a family, a husband, a wife, a servant, be like?—Especially the conceptions of idolatry, unselfishness, magnanimity, worldliness, and immodesty, still largely alien to the child mind, should be diligently reviewed and explained. \emph{3. Religious concepts}: Altar, covenant, oath. \emph{4. Parts of Catechism as derived from the several Bible stories. 5. Stanzas of hymns. 6. A review of the geographical material and of that pertaining to cultural history should not be forgotten}.—From F. Zange, \emph{Leitfaden fuer den ev. Religionsunterricht}, vol. i and ii, 1906. A simpler arrangement would be to assign central position to the subject: Abraham the Father of the Faithful, and to group all material around it.

As a basis for a review of the entire history of Israel, Isaiah 5: 1–7, is particularly adapted; the class should show in detail 1. What God did for His people since their deliverance out of Egypt; 2. How sour instead of sweet grapes grew on the vine of His planting; 3. How the Israelites incurred His wrath.

\subsection[8. Insertion of the geographical element]{8. Insertion of the geographical element and of that of cultural history at the proper place for a broader survey}

a) Israel entered into Canaan and the Lord gave His people rest from all their enemies and showed them all the kindness He had promised. The teacher adds the following: The Israelites fared very well in their new land. The shepherds in the land east of the Jordan and in the
\quellseite{603}
mountains of Judah lived in tents, but among the rest of the Israelites each family had its own house, and an arbor of grape vines around it. The vines grew rank and luxuriantly, climbing all over the arbor and covering the flat roof of the house. In summer and autumn there was a roof of grapes from the arbor to the house. What advantage was there in that for the people? (Plenty of shade.) It was their delight to spend their time there in the shade. Only when they could not remain outside, did they go into the house. When was that? (At night and when it rained.) Of course, their houses were not so comfortably arranged as ours; for they had only one room, with one door and no windows. You already know what kind of roof they had. What kind was it? You would not care to live in such houses, but the Israelites were well satisfied. Why? They knew of no other kind, and they needed nothing different, for they could be outside all day.—Summary, House.

It was delightful out of doors. About the towns and villages stretched the \emph{fields and meadows}. What grew on the meadows? What grew in the fields? (Wheat and barley.) On the hillsides were forests of oak and pine trees, groves of palms, orange, and citron gardens. What fruit was there in profusion? Besides, there were fine apples, plums, apricots, peaches, almonds, and the like. There also grew the fruit of which the spies brought samples. Do you recall what they were? Already in January, when we have snow and ice, the almond tree was blooming there, and the fields were covered with violets, tulips, lilies, etc. At the end of April the barley harvest began, and at the end of May, the wheat harvest. Then came the fig harvest, and from August on the grape and olive harvest. You know what olive oil is; for what purpose is it used? This oil is made from olives. Olives grow on trees about the size and color of our plum trees, and are of different shades. There are blue, black, red, white and green olives. They are pressed, and the oil that comes from them is valuable. Olive-trees are also called oil-trees. Upon what mountain to which Jesus often went with His disciples were such trees? From what did this mountain get its name?—Summary, fruits of that land.

The Israelites \emph{worked diligently} in this beautiful land. Seedtime came in October. What did they do then? That was done as soon as the first rain fell. From the end of October until the beginning of March there was much rain. What time of the year is that with us? But there was also \emph{plenty} to do during the rainy
\quellseite{604}
season, especially when heavy rains damaged the vine-yards. What damage? (Ground washed away, vines washed out, vine-yard walls undermined and thrown down, etc.) What did they then have to do?—The last part of March and the beginning of April was the short rainy period. This they called the late rain. After that the sun shone hot from the cloudless sky, and the grain soon ripened. Which ripened first? Then began a season of continual labor. What had to be done with the ripe grain? Where do we put the sheaves from the field? The Israelites had at that time no barns, and so they threshed the grain upon large flat places in the fields, which they called threshing floors. These floors were circular in shape and were surrounded by a low wall. How do we thresh the grain? With flails and threshing-machines. The Israelites filled the threshing floor with grain and then drove oxen and cows around over it all day. How did the grain get out of the ears? (It was trodden out by the animals.) Then the grain had to be separated from the straw. In the evening, when the sea wind blew, the men took a shovel and threw the trodden grain in the air. What then happened with the heavy grain? (It fell upon the floor.) And what happened to the lighter straw? (It was blown away.)—Summary, Seed-time and harvest.

\emph{Harvesting} was joyful work. By what do we know this? (Their singing.) At this season all the people remained in the fields day and night, masters as well as laborers. Why? (To guard against thieves and fire.) They also prepared their meals there; but, naturally, they were very simple. Ears of grain were roasted at the fire, then rubbed out and eaten. Or the grain was ground between stones. What was made of the grain? (Flour, really groats.) The meal was then mixed with water into dough and thin cakes were then baked of it. And how did they quench their thirst? (With water.) They mixed the water with vinegar. Why? (Such a drink quenches thirst much better.) Where did they carry the grain? (Into barns.) There was little wood in the land of Canaan, and to build barns was too costly. So they made great caves in the earth.—Summary, Harvest.

The \emph{grape harvest} was still more joyful. During the grape season, from August to November, the entire family lived in the vineyard. In nearly every vineyard there was a tower. Why? (To stay at night and in case of rain.) Why did the people stay so long in the vineyard? (There was much work to be done there.)

\quellseite{605}
The grapes grew much more luxuriantly there than they do for us. Many a grape-vine was as thick as a tree and bore grapes without number. Such a vine bore as many as a thousand clusters of grapes. What did they make out of so many grapes? (They cooked them into grape honey and prepared them for wine.) The wine they made in the wine-press. That was a place sunk in the ground, or hewn out of a rock, about 2 yards long, 2 yards wide and a half yard deep. (Show the size.) The men tramped upon the grapes in this wine-press. What flowed out of them? (The juice.) Before every wine-press there was a deep place hewn out of the rock, and a trench led from the press to this basin. In the press and at the mouth of the trench there was a grating. Why? (To serve as a sieve and hold back the grape-skins.) What were the Israelites as long as they were in the vineyards? (Vine-dressers or wine-growers.) Why were these wine-growers so happy at their work? (Good harvest, juicy grapes, sweet must.) How did they show their joy? (By song, music with harp and zither, cymbals, dancing.) --Summary, Wine harvest.

\subsubsection{b) Geographical description of a specific place}

\emph{Nazareth}. The home of Jesus was a small, unimportant place, but beautifully situated in the midst of the province of Galilee. You can see this from the picture of the present-day Nazareth (En Nasira; Picture by Langl). A description of this picture shows the homes of the nearly 7,000 inhabitants scattered over the side of an almost unscalable mountain. Today the majority of the people are Christians. The place is surrounded by gardens of olive trees, fig trees, palms, and cypresses. The houses are of stone and have flat roofs, as is the custom in the East. (Reminder of David.) In the foreground we see a woman carrying a water jar upon her head. She is coming from the tree-shaded Virgin-spring which bubbles in the valley before the town. It is the only spring in Nazareth; and it is here that Mary, the mother of Jesus, may, with the child Jesus, have often drawn water. Farther down, but still on the heights around the town, are the orchards, with cultivated fields between them. Farther in the background are the less fertile heights which serve as pasture for the sheep and the goats. Beyond the height stretches the road to the fruitful Plain of Jezreel, watered by the Kishon (Reminder of Elijah), and surrounded in the blue distance by the ranges of the Carmel mountains. If we could climb to the summit of the mountain with those two Franciscan monks, whose
\quellseite{606}
cloister, however, is not visible on the picture, we should see to the far South the mountains of Gilboa. (Reminder of Saul’s death.) We should farther behold the steep heights of Tabor in the East, the mountains of Galilee to the North, and beyond them the peaks of Lebanon and Hermon, snow-covered even in summer, and the Mediterranean Sea in the far distant west.—Bittdorf, \emph{Methodik}.

\section{2. Instruction in the Catechism}

\subsection[1. A Catechesis to Begin Instruction]{1. A Catechesis with Which to Begin the Instruction of the Catechumens}

The catechist sings with the children the first three stanzas of the hymn, “\emph{Little Children Come to Jesus}”, reads Mark 10:13–16, and says a short prayer. Then he begins his catechization in the following manner:

\emph{Aim: Dear Children! We would hear today that the Saviour wants to have also you children with Him to bless you.}

\subsubsection{1. Developing the truth from the intuitional material} Among all the stories of the Bible I know one which should be especially dear to you and all children. It is so beautiful that the picture of it should hang in every Christian school. I know of none which shows more clearly that the Lord Jesus dearly loves children. Of what am I thinking? Yes, the story of Jesus blessing the little children. That is a true children’s story, and one can imagine nothing more lovely. We want to study it more closely. It was upon a day when many people had again followed Jesus and He had been ceaselessly teaching them and healing their sick. It was, therefore, a day of extreme exertion for the Lord. Finally came the mothers, pressing through the crowd and bringing their children to Jesus. The larger ones they led by the hand; the smaller ones they carried in their arms. What was the Lord Jesus to do for these little ones? He was to touch and bless them. That was a pious wish. Those mothers had seen that the sick were healed; that the blind received their sight, and that the lame walked when Jesus laid His hand upon them. They had also seen that it was good to be in the presence of Jesus, and had beheld how great a blessing came from Jesus upon all whom He touched. For whom did they wish such a blessing? For their children. Was there something lacking, or were they sick? No. Then what kind of blessing did these mothers want for their children if they were not sick? A blessing for their
\quellseite{607}
souls. Such it would have been, even if these mothers had not known what they wished for their children. They surely wished their children to become good, pious children. But Jesus had labored the whole day and was weary. What does one desire when he is weary? He seeks rest. What would Jesus then gladly have done? He would gladly have rested. Yes, we can well think so: for although our Saviour was true God, He was at the same time true man; and, while here on earth, He often became weary like other men. Who else must have thought that the Saviour wished to rest? The disciples. What did they therefore do when they saw these mothers coming with their children? They forbade them. Yes, they motioned to them with their hands that they should stay away. And as that did not keep them back, they spoke to them with hard words. And when that did not stop them, the disciples threatened them with force. What, under the circumstances, did it seem impossible for these mothers to do? To get near to Jesus. It seemed as though they would have to leave without accomplishing their purpose. Who, then, took the matter in hand? The Lord Jesus. How did He feel when He saw what was happening? He was much displeased. With whom? With the mothers who were coming to disturb His rest? Oh, no, with the disciples. Why? Because they would not let the mothers come with their children. What did Jesus say to His disciples? Suffer the little children to come unto Me, and forbid them not. He knew that these mothers were eager to bring their children to Him; and what must He have wished when He rebuked His disciples? He wanted to have the children with Him. But was He not weary in consequence of His preaching and healing of the sick? Certainly He was weary. If He was weary and yet wanted the children to be permitted to come to Him, which was to Him most important, to rest or to have the children with Him and to rejoice with them? To have them with Him and to rejoice with them. Yes, children, we have such a Saviour who forgets all weariness when souls come to Him and desire His blessing. Even when it is who that comes? Even when it is a child. Therefore we learn from this history: \emph{Jesus wants the children to be with Him, even when He is weary. He always has time for them.} Let us say that together. Again! Again! Anna, say it alone. Frank, repeat it. Now I shall write it on the blackboard.

The disciples did not wish to let the mothers come to Jesus with their children, because He was tired and needed rest. Can you give
\quellseite{608}
another reason why they probably wanted to turn these mothers away? Because they thought that Jesus could do nothing for these children. That He could be a Saviour for grown folk they had learned from His dealings with them and others, but how could He help these little ones, who could not even understand what He said? However, who must have trusted that He could give something to these children? The mothers must have so trusted. Otherwise what would they certainly not have done? They would not have tried so hard to force their way to Him with their children. Upon whose side did Jesus place Himself, that of the disciples or that of the mothers? Upon the side of the mothers. What did He say that He could and would give the children? The Kingdom of Heaven. What are His exact words? For of such is the Kingdom of Heaven. Not only grown folk, but also children can be His disciples and therewith members of the kingdom of heaven. No one—no one will be shut out unless he shuts himself out. The Lord Jesus has expressed Himself emphatically with regard to this, in His words to His disciples. Who knows what these words are? “Verily I say unto you, whosoever shall not receive the kingdom”, etc. Like whom must therefore the grown folks become before Jesus can use them in His kingdom? Like the children. The Lord Jesus did not thereby mean that the disciples and all adults must again become small and weak like the children, but He means that the grown-ups must become as open and receptive for the good and heavenly as children are in comparison with their elders before He can use them in His kingdom. The disciples must have opened their eyes wide at these words. They thought that Jesus could do nothing for whom? For the children. And who must have rejoiced at these words of Jesus? The mothers. Yes, indeed, for they perceived that Jesus could give the children much more than they thought. How did Jesus prove before all people that He can be something to the children? What did He do with them? He blessed them. More exactly! He took them up in His arms, put His hands upon them, and blessed them. There you see the true friend of children, who casts no one out, who also presses the smallest and least to His heart and would give them all the blessings of His heavenly kingdom. Let us note: \emph{Jesus can and will give also children His heavenly gifts.} Let us repeat that together a number of times. John, say it alone. Louis, do the same. And now I shall also write this upon the blackboard. How many things have we thus far learned from our history? Two. What is the first? Jesus
\quellseite{609}
wants the children to be with Him, even when He is weary. What is the second? Jesus can and will give also children His heavenly gifts.

\subsubsection{2. Comparison} The Lord Jesus always showed that He had a place in His heart for the children. Take your New Testaments and turn to the Gospel of Matthew, chapter 18. Read verses 1—4. About what were the disciples at strife? Which one of them was to be greatest in the Kingdom of Heaven. Where did they therefore all expect to be together? In the Kingdom of Heaven. They felt certain of that; but about what did they contend? Which of them was to be the greatest—the chief, in the Kingdom of Heaven? That one or the other of them might never get into heaven was not questioned by any of them. They had forsaken all and had followed Jesus, and how could they be shut out when Jesus should as king set up His kingdom? What did the Lord Jesus do to awaken the disciples from their dream of security? He took a child that happened to be in the house and placed it in the midst of the disciples. And what did He say to them? Read it in verses 3 and 4. “Verily I say unto you . . . . . . . greatest in the Kingdom of Heaven.” About which question should the disciples therefore not argue? Who of them should be greatest in the Kingdom of Heaven. But of what should they strive to be certain? That, above all things, they should enter the Kingdom of Heaven. And on what condition only were they to enter it? That they become converted and become as little children. As small and weak and ignorant as little children? No, as willing, simple, and humble as they. Who, therefore, will go into the Kingdom of God before them, according to Jesus’ words? The children. Hence, Jesus must have determined to receive also whom into His kingdom? The children. Certainly. What He here says is the same as when He says: Suffer the little children to come unto Me; for of such is the Kingdom of Heaven.—Later, when Jesus rode into Jerusalem on Palm Sunday, He declared that the children had understood Him better than many grown people, and that they gave Him great cause for joy. Open your Testaments at Matthew the 21st chapter. John, read verses 15 and 16. “And when the chief priests . . . . . perfected praise.” You know that when Jesus rode into Jerusalem on Palm Sunday a great procession of people praised Him. And that was so contagious that also who joined in? The children. What did they shout? “Hosanna to the Son of David.” They did right, for Jesus desired to set up His kingdom, in which they too should be welcomed and receive all that they needed. Who only
\quellseite{610}
stood by and grudged Jesus this joyous reception? The chief priests and scribes. What, according to their opinion, should Jesus have done? He should have silenced the children. But what did Jesus reply? “Have ye never read . . . . . . . perfected praise?” Yes, He took the children under His protection and said that they were doing precisely what they should do; and that what the Psalmist once said of the babes and sucklings, they were doing, namely, praising God for setting up His kingdom. Thus we again see how much Jesus loved the children and how He had determined to receive them into His kingdom and to bless them therein. That which we learned from the history of Jesus blessing the little children is entirely correct. What was it? 1. Jesus wants the children to be with Him, even when He is weary. 2. Jesus can and will give also children His heavenly gifts.

\subsubsection{3. Valuation} In what did Jesus desire to give children a part? In the Kingdom of Heaven. Is that anything so very great and precious? We will consider this more closely. Tell me, is it something great to be the child of a king? Yes, it is something great. The child of a king has everything it needs. There are many servants who wait upon it and fulfil every wish; and when the child has grown up, sooner or later, he will help rule the land. The greater the kingdom, the greater is the glory of such a king. When it was expected that a son would be born to Emperor Napoleon, they had a golden cradle ready, in which they afterward laid the child. Whose kingdom is greater than that of any king or emperor on earth? The Kingdom of God. Who will therefore have it better and more glorious, one who has part in an earthly kingdom or one who has part in the Kingdom of Heaven? He who has part in the Kingdom of Heaven. Where is the Kingdom of Heaven?—You can not tell? The Kingdom of Heaven is where Jesus Christ is. He who has Jesus Christ as his friend and helper is in the Kingdom of Heaven, and he is far better off than the son of the richest king upon earth. Jesus can bless, protect, and comfort, better than any king upon earth. What did Peter say when He saw Jesus upon the Mount of Transfiguration? It is good to be here. Peter was right, children! Nowhere is it better than with Jesus. Therefore we sing in church, “\emph{Who is, Jesus blessed, Like to Thee, sweet Rest?}” Ask the disciples; ask John, who lay upon His breast; ask the children whom He embraced and kissed; ask Mary the sister of Martha, who knew nothing sweeter than to be with Jesus! The person who is with Jesus is blessed; for he has all that he needs, even when he becomes sick and poor, even in trouble and under the
\quellseite{611}
cross. Instead of saying, The Kingdom of God is where Jesus is, we can also say, The Kingdom of Heaven is there where we have God as our Father and are become His children. When is a child happiest and most blessed,—when it is with whom? When it is with father and mother. Why were Adam and Eve so happy and blessed in Paradise? Because God was a father to them and they could speak to Him as children. When shall we be happiest? When God is our father and we are His children. Yes, children, that is the true blessedness. Even Heaven itself would be of no benefit to us, if God would not be there as our father and we as His children.

\subsubsection{4. Application} Of those children in the story whose mothers brought them to Jesus, He said: “Of such is the Kingdom of Heaven.” and we have just learned that to be in heaven means nothing less than to be with whom? With the Lord Jesus. And to have whom as father? To have God as father. And what is man's condition when he is with Jesus and has God as his father? He is blessed, then and then alone. Have you been brought to Jesus? Yes. When were you brought to Jesus? When I was baptized. Who brought you to Jesus? My parents. What must they have thought of Baptism if they had you baptized? They must have thought that Baptism was good and beneficial. Yes, they knew that, in Baptism, the Lord Jesus would bless you as He blessed those children who were brought to Him. They had learned from the words of Jesus that He wanted also whom with Him? The children. And that He could and would help also whom? The children. They had learned that He wanted to give the children a share in what? In the Kingdom of Heaven. They well knew that to be in the Kingdom of Heaven meant the same as what? As being with Jesus and having God as their father. And it was also known to them that you would be in what condition, too, if you were with Jesus and had God as your father? That we should be happy and blessed. Yes, that is why your parents brought you to Jesus in Holy Baptism. They wanted you to be with Jesus, to have God as your father, and thus to be happy and blessed children. See how well they meant it! But is it really true that, by Baptism, one gets God as his father and becomes His child? Yes. The Apostle Paul was strongly convinced of this truth; for when he writes in his letter to Titus (open New Testament at place): “According to His mercy He saved us by the washing of regeneration, etc.”, he means nothing else than this. Before he believed on Jesus and was baptized by Ananias in Damascus, he was not with Jesus, God was not his father, and he certainly
\quellseite{612}
was not blessed. He was very unhappy and had no peace with God. But when he believed and was baptized, he was as if new-born, happy, and blessed; for then he knew that God was his father and he was God’s child, and could say: Now, with all my sins and in every need I can go to Him with all assurance and confidence. Paul also thought of this when he wrote to the Galatians (find place in Testament): “Ye are all the children of God, etc.” All who believe in Christ and are baptized are God’s children, and Jesus surrounds them all with His presence and protection, as the clothing we put on covers the body. Paul’s baptism was considered by him to be the great turning-point in his life. Before that, he was without Jesus, without the Father, and therefore, instead of being blessed, lost and condemned; after being baptized, he was with Jesus and the Father, and therefore blessed. So, then, when did Paul’s blessedness begin? With his baptism. Not only when, therefore, did he consider himself blessed? After his death and in heaven. Hence he did not write: God shall save us, but; God hath saved us; He has already done so. And as it was with Paul, so, dear children, is it with us, with you. Also you are saved through Baptism and have received full salvation; for you have been brought to Jesus. He bestowed His kingdom upon you in that He became your Saviour and Friend, and in that God became your Father. The salvation which you shall have in heaven grows out of that salvation which was given you in Baptism, as the apple-tree grows from the seed. Whoever is not saved on earth will not be saved in heaven; for to be saved means nothing else than to be with Jesus, to have God as Father and to be His child. Let us repeat that! Again! By what means were you brought to Jesus, get God as your Father and become His children? Through Baptism. Oh, then Baptism is something great and glorious. A certain French king once said: “The three handfuls of water with which I was sprinkled in my baptism are more precious than the royal crown which I now wear upon my head.” A baptized child is like the child of a king in its cradle; salvation is just as certain to him as is the crown to the king’s son. Nay, he is much more fortunate, for the king’s child must wait many years after its birth before it receives the crown, while the child that is baptized receives salvation immediately in Baptism. The Lord Jesus takes it and blesses it, makes of it a child of God, which can always with all confidence go to God as a child to its Father. So you who are baptized are not the losers thereby. You have received in Baptism greater salvation than those children
\quellseite{613}
who were brought to Jesus by their mothers. And if you have been brought to Jesus by Baptism, have God as your Father and are in truth saved, what should, for that reason, be your greatest care? That we hold fast to this salvation and not lose it. That means that you remain what? That we remain God’s children. Yes, children, remain with Him, with Jesus, with God your father. And to this end may these hours of religious instruction help—that you remain the children of God!

Now open your catechisms and read, page 27: (“That . . . I may remain a child of God”) (the lesson to be assigned). Also read the three Scripture passages found in that connection. You can all understand this now very well. Memorize the Scripture passages for tomorrow. But we first want to repeat each of them several times together. Read carefully that which is printed in large type. I shall ask you about it tomorrow. And now, since we have been brought to Jesus by Baptism, have God as our Father, and therewith are in possession of salvation, let us sing a hymn on Baptism. I will read the first verses and then you may read them. (Here the catechist should paraphrase the difficult words.) And finally we will sing them together.

\subsection{2. Second Catechesis (outlined)}

\emph{Aim: The Holy Scriptures the Book of all Books.}

\subsubsection{1. Developing the truth from the intuitional material at hand:}

From the Scriptures Jesus shows the sad disciples on the way to Emmaus, who believed everything to be lost and who could nowhere find comfort and new hope, that they have no reason to mourn. He shows them the way God took to redeem mankind and make them His children, and also the ways they now must go. Their hearts rejoice and their eyes become bright. They now understand the ways of God.

\subsubsection{2. Comparison:} Philip also uses the Scriptures when he would show the eunuch of Ethiopia the same, and the latter is thereby led to Baptism and joyfully goes his way.—The Bereans search the Scriptures for themselves; and each new discovery makes them the more certain that Jesus is the Redeemer through whom alone they can become and remain God’s children. The grandmother of Timothy instructs him in the Scriptures from childhood, and Paul directs him to them when he has reached manhood, as the book that can make him wise unto salvation.

\quellseite{614}
\subsubsection{3. Valuation:} Therefore the Scriptures are greater and more precious than all other books. Therein we are instructed concerning the highest and most important things. Therein it is not mere men, but God who speaks to us. 2 Pet. 1:21. Ask the psalmist and the prophets (Scripture verses); ask Luther; ask many great ones in the world, and many a humble mother. You will find that they have all learned this: By means of the Bible God Himself speaks to us, instructs, comforts, admonishes as none other. If we would show the heathen the way they can become the children of God, we take the Bible in hand. If we ourselves want to learn this fact anew, we listen to the preaching from the Bible, and also search the Word ourselves. What do our hymns say?

\subsubsection{4. Application} And now we also want to learn from the Scriptures how we became God’s children and how we may remain such.—Search the Scriptures.—Show how that is done.—Compass—if we would go the right way, we must always—give individual cases—examine them.

\subsection{3. A Catechesis on the First Commandment}

The catechist sings with the class several stanzas of the hymn, “\emph{If thou but suffer God to guide thee}”, etc., reads Exodus 14, and says a brief prayer. Then he reviews the previous lesson, whereby the present one is introduced.

\emph{Aim: We would learn that children of God trust in God above all things.}

\subsubsection{1. Developing this truth from the intuitional material}

To what did God call Moses when He appeared to him on Mount Horeb and spoke to him out of the burning bush? Moses was to lead Israel out of Egypt. Yes, he should be their deliverer. That was a difficult task. Why was it a difficult task? Because Pharao did not want to let them go. And why was it for Moses especially a risky matter to return to Egypt and lead Israel out? Because he was once compelled to flee from Pharao. What did Moses say when God assigned him this difficult task? He did not want to go and raised many objections. What did he finally say to God? Send whom Thou wilt. Moses was not altogether wrong. Had he gone to Pharao in his own strength and attempted by his own wisdom to lead Israel out of Egypt, he had better not have undertaken it. He would have failed miserably. When did Moses truly experience this? When,
\quellseite{615}
forty years before, he had tried it by his own wisdom and strength. But now he wanted to do it under God’s commission. That was entirely different. What could he rightly expect when he was acting by God’s order? That God would help him. Certainly God also promised him that He would be with him. That was a glorious promise, and Moses could say to himself: Yes, if God will go with me, I will make the venture. He is mightier than Pharao and everything else. He is the creator of heaven and earth. He can show me how to begin and how to accomplish the work. I could have no better companion. If he is for me, who can be against me? When God said to Moses, I will go with thee, He after a manner stretched out His hand and said: I will lead thee. And what did Moses finally do? He laid his hand in God’s hand that He might lead him. And of what could Moses be certain when God was leading him? That he would be rightly led. Yes, Moses trusted God; for when one is given a hard task and he places his hand in God’s hand and trusts that He will lead him aright, that he will succeed with God’s help, we say of him that he trusts God. We will write that on the blackboard: \emph{He who trusts God, places his hand in God’s hand that He may lead him, and is certain that with God’s help he can accomplish the most difficult task.}

As Moses went to Egypt, which words of the Lord must have been constantly ringing in his ears? The words: I will be with thee. Yes, they were his comfort, and always gave him fresh courage. When his brother Aaron met him on the way he was very glad; but upon whom did he rely more than upon Aaron? Upon God the Lord. When would Moses later especially have relied upon God’s assistance? When he had to go to the children of Israel and tell them that God had appeared to him and commissioned him to lead them out of Egypt. But on what occasion to an even greater degree? When he had to appear before Pharao. Yes, all kinds of thoughts must then have arisen in Moses’ mind, such as these: Suppose he recognizes you, and still remembers why you fled before, and imprisons you; or has you killed! With such thoughts as these it would not have been surprising if he had lost courage and turned away. Why did he not do so? Why did he go calm and unafraid to Pharao to fulfil his mission? Because he trusted in God. He heard God saying again and again, I will be with thee, and he gripped His hand the tighter. And what did he trust God to do? To help him and lead him aright. At any rate he was sure that as long as God stood by him he could
\quellseite{616}
not fail. He must have said to himself, “Commit thy way unto the Lord; trust also in Him, and He shall bring it to pass.” Did God stand by him and champion his cause? Yes. True, in the beginning it did not seem so. What reply did Pharao give when Moses demanded of him that he let Israel go? I will not let Israel go. And what did Pharao say when Moses called upon God the Lord? Who is the Lord that I should obey his voice? I know nothing of the Lord. And when Pharao refused again and again, and finally forbade Moses to appear before him, it must have been difficult for him to trust God. All seemed to be in vain. Did he then give up his trust in God? No. What rather did he do? He prayed to God. He kept his hand in God’s and continued to hope that in the end God would give him success in his undertaking and lead Israel out. Even though he might have to wait long and Pharao might refuse again and again, he said, “Nevertheless I am continually with Thee: Thou hast holden me by my right hand. Thou shalt guide me with Thy counsel, and afterward receive me to glory.” With such thoughts as these he strengthened his courage and assured himself that since God had bidden him take this way, He would see to it that he did not fail. We note: To trust God is to put one’s hand in God’s hand, to leave it there and to hold it ever fast, even though one hindrance come upon the other and make the performance of our duty more and more difficult. (Blackboard).

After Pharao had let the Children of Israel go, where was Moses trust in God again put to the test? At the Red Sea. Other trials of his trust had not been wanting, but the hardest came at the Red Sea. When God showed him the road he must travel, he had to trust Him. How did God show the road? He sent a pillar of cloud before the Children of Israel. Did He in this way guide them to Canaan by the nearest way? No. At first it seemed as though God would lead them straight to Canaan, but in which direction did He suddenly guide them? Toward Egypt. Correct. He made a sharp turn southward and so led them farther and farther into the country of the enemy. To whom especially must strange thoughts have come? To Moses. It was probably hard for him to understand why they were being led that way, for Moses knew the way to Canaan well. What would he have done if he had followed his own understanding? He would have said, we are being led the wrong way, and he would have let go of God’s hand. What would he then, of course, not have done? He would have no longer trusted God. What does it mean to trust God? To
\quellseite{617}
place your hand in God’s hand and let Him lead you. Yes, and to be sure of what? That His way is the best. But Moses did not cease trusting God. He could not understand why God was leading them this way; to his mind it was all wrong. But he leaned more upon God than upon his own understanding. And what must Moses, more than ever before, have said to himself? Think of the Scripture passages you have already learned. Commit thy way unto the Lord; trust also in Him and He shall bring it to pass. And if he had been acquainted with our hymnal and learned our hymns, which one would he have thought of? “\emph{Commit thy way confiding, When trials here arise.}” And when the morning revealed the fact that they had travelled southward the whole night and were going farther into the land of the enemy, and his reason was telling him more and more plainly that this was wrong,—what did he then do with God’s hand? He still held fast to it. Correct. He clasped it so much the tighter and said, “Nevertheless, nevertheless, I am continually with Thee”, etc. But you said before that Moses’ trust was given its severest test at the Red Sea. Whither did God finally lead the Israelites? To the Red Sea. Yes, the Red Sea lay before them, and what appeared to the right and left? High mountains. And who came behind them? Pharao and his army. Israel’s case seemed desperate, shut in on all sides like a mouse in a trap or a fish in a net. There seemed to be no way of escape. And who was it that had thus led Israel? God himself. And had God led them wrongly and brought them into misfortune? It appeared so. Who began to believe this? Many of the people. What did they say to Moses? “Why did you lead us out of Egypt? Were there not graves in Egypt, that we should perish in the wilderness?” But who could not believe that God had led them wrongly? Moses. He had beheld too many of God’s wonderful ways, and could not now cast his trust in Him away. It is true that his spirit was deeply grieved and that he could find no way out of the difficulty, but he did not give up hoping in God. God had called him to deliver the people of Israel out of Egypt; He had led Him in the past, and he could not believe that God would now permit him and his people to perish. He held fast to the belief that God knew the way out of the trouble. And when his soul trembled and feared, he said to it: “Why art thou cast down, O my soul? and why art thou disquieted within me?” etc. Here we first clearly see what true trust in God is. He who in truth trusts God, is also
\quellseite{618}
certain that God’s way is the right one even when? Even when it leads into deepest distress and one can find no way of escape. Repeat this a number of times. \emph{Who trusts God is even then entirely sure that God’s way is the best even then when it leads into deepest distress and one can find no way of escape.} I will also write it upon the blackboard. Moses had strong trust in God. What did he do in the strength of such confidence? He cried unto God. What did he say? Lord, help us! Certainly, he reminded God of His omnipotence by which He can save us out of every trouble, and of which promise did he undoubtedly remind Him? That He would be with him. He reminded God of His faithfulness, according to which He must now keep His word. “I have trusted in Thee; Thou canst not and wilt not let me be put to shame”—thus he must have called to God. And what was God’s reply? “Lift thou up thy rod, and stretch out thine hand over the sea, and divide it; and the Children of Israel shall go on dry ground through the midst of the sea. And the Egyptians shall know that I am the Lord.” Oh, that was a new trial of his trust. He, a weak human being was to divide the waters of the sea and make a way through for his people. Did his faith stand this test? Yes. What did Moses do? He stretched out his hand over the sea. Behold him standing there by the side of the sea, with his hand holding the rod stretched out over the waters. His reason told him: Moses, you are making yourself ridiculous; never before has a man parted the waters of the sea! But Moses stood firm in this confidence: \emph{He trusted that God could even work a miracle to save His people.} That was the highest point of his trust. (Blackboard.) And see how contagious it was! What did the people do when they saw the parted waters standing like walls on either side and the way open through the midst of the sea? They went through. What trust in God! They, like Moses, laid their hands in God’s, that He might lead them. And what did they trust? That God’s leading would be the best even though it led through the midst of the sea. Children, that is trusting in God above all things, when we, like Moses, first do what? When we lay our hand in God’s hand, that He may lead us and are sure that with God we can perform the most difficult task. In the second place, when we do what? When, like Moses, we leave our hand in God’s hand and only hold the tighter when one hindrance comes upon another to make our task more difficult. Thirdly, when we do what? When we are certain that
\quellseite{619}
God is the best guide, even though we come into deepest distress and can find no way out. Yes, and what is now the last thing this trust in God clings to? That God can work a miracle if He will, to save His flock.—It was not strange that God was pleased with Moses; for he trusted God as but few trust Him.

\subsubsection{2. Comparison}

We see, however, that still other men and women trusted God. Name a king of Israel who trusted God above all things. David. When did David show his trust in God? When he went forth against Goliath the giant. What did he, accordingly, say to Goliath? “You come out to me with sword and spear, but I come in the name of the God of Israel.” Here also there was a great task to perform. In order to perform it, David laid his hand in God’s and trusted that He would help him slay Goliath and so fulfil his duty. When did David trust God in suffering? When Saul persecuted him for years. There he laid his hand in God’s hand and was sure that God would do what? That He would lead him aright. Yes, and even though he was forced to wait for years and years before obtaining the throne which had been promised him, and was forced to flee from place to place, what did he still firmly believe? That God was leading him aright. What, therefore, did he sing in the twenty-third Psalm? “He leadeth me in the paths of righteousness.” When did Abraham show his trust? When he left his homeland. Yes, when he left his own land to go into a land the name of which he did not even know. Then he laid his hand in God’s hand, that He might lead him. And of what was he certain? That God would be the best guide for him. Of what hymn stanza are we very strongly reminded when we see Abraham thus going forth into the strange, unknown land? “\emph{Commit thy way confiding}”, etc. What did Abraham thereby show? His trust in God. Here again we see that it dealt with the performance of a great task, which Abraham undertook, upheld by true faith in God. Which one of Abraham’s descendants also trusted God through long, hard trials? Joseph. When did his trials begin? When he was sold into Egypt by his brethren. And when did they end? When he became ruler over Egypt. How many years was that? More than 13 years. Oh, that was a long time of trial. What thoughts must certainly have often come to him? That God had forgotten him. When especially did such thoughts come to him? When Potiphar had him cast into prison. Again when? When the cup-bearer was freed and forgot his promise. But what did Joseph always do with such thoughts? He always overcame them. Yes, he trusted God. Even in the greatest trials he did not let go of God’s hand; but what did he do? He held it all the
\quellseite{620}
tighter. What confidence in God did he thereby show? That He had not forgotten him, appearances notwithstanding, and that He was the best guide nevertheless. Thoughts such as Paul expresses in the passage, “We know that all things work together for good to them that love God”, must often have passed through his soul. And thus he showed that the true children of God should not throw away their trust even when? Even when the distress is great and long continued. When we now see that Abraham, Joseph, Moses, David trusted in God above all things—and we could mention many others, especially our dear Dr. Martin Luther—who should then trust God? All who would be God’s children. Yes, to be God’s child and to calmly and trustfully place one’s hand in God’s that He may lead one, and to be sure that—even in the greatest trials—He is the best guide: these things necessarily belong together. May God’s children then trust no one but God? Oh, yes. In whom does the child trust? Its parents. In whom the pupil? His teacher. In whom the king? His people. In whom the sick? The physician. But the true children of God, whom do they trust more than any one else? God. Yes, they trust in God above all things. They regard Him as their only refuge and help. They are certain that others can help them only when who permits? When God permits. Hence, what kind of father do we call God in the First Article, because we Christians trust in Him that He will always lead us aright and help us in every need? The Father Almighty.\footnote{Here it would be better to stop and postpone the Valuation and Application until the next time. Such important fundamental truths as are here dealt with must be thoroughly considered, which may well take half a week. If the catechist has time, he may thoroughly go through the hymns “Commit thy way confinding” and “If thou but suffer God to guide thee”. Much else in the explanation of the Commandments can be briefly treated—the negative more briefly than the positive.}

\subsubsection{3. Valuation} Did Moses receive any benefit from his trust in God? Yes. Who went with him when he trusted in God’s help and went to Egypt? God went with him. How did that show itself? By the signs and wonders which he did before Pharao. By whose power had he forty years before tried to free the Children of Israel? By his own power. What did he then accomplish? Nothing. What was he then forced to do? He was forced to flee. But what was he now able to do, when, trusting God, he undertook the great work? He could lead Israel forth. What effect did his trust in God have upon him? It made him strong. Yes, trust in God makes one strong and
\quellseite{621}
courageous. We see this also in David. How large was David, compared with Goliath? He was small and weak and wholly without experience in handling weapons of war. He was so weak that who thought it impossible for him to conquer Goliath? His brethren and King Saul. But with what success did he nevertheless meet? He killed the giant Goliath. How did his trust in God make him? Strong and courageous. And fearless, we may add. Fearlessly Moses went before Pharaoh, fearlessly David went against Goliath, fearlessly Abraham went forth into the unknown country. When, therefore, can one undertake a task with courage and strength, and without fear? When one trusts God above all things. (Of course, there is one thing we must not forget. Who commanded Abraham to go into that strange land? God. Who sent Moses to Egypt? The Lord God. Who bade David to go forth against Goliath? God told him, for He sent him to the camp at the time when he had to hear how Goliath blasphemed God. By whom, then, must the task be imposed if one would assume it in God’s name and be sure that God will make one strong, courageous and fearless? By God. One has no right to push oneself where one is not called and then expect God to help him. To thoughtlessly thrust one’s life into danger and then expect God to help, is not trusting God, but tempting Him, and that is sin.) Upon what then does all depend when God gives us a difficult duty to perform? Upon this, that we undertake it trusting in God. When Lincoln sounded the call for troops after the shelling of Fort Sumter, he had to contend not only against the war-like Confederates but also against a strong opposition in his own country. Well might he have trembled. But fearless, strong, and courageous he went to war; for it was not he that had provoked the war. He could trust in God. His good conscience and trust in God made him courageous, fearless, and strong. Now let us repeat together: \emph{Trust in God makes us fearless, courageous and strong.} We will say it again. (Blackboard.) Of this the prophet once spoke, who himself always depended upon God for all things. It was the prophet Isaiah. Read Isaiah 40:31. Anna, read the verse, “But they that wait on the Lord”, etc. O children, all good people who trusted in the Lord have experienced how they received new strength. More than once they became so weary that it seemed as though they could not go a step farther; but what did God give those who trusted in Him? New strength. What could they therefore do? They could mount up like an eagle. How does the eagle fly? With power and to a great height. Yes, he seems to be absolutely tireless. And what else can they do? Run without becoming weary, walk without becoming
\quellseite{622}
faint. We see this in Elijah when he ate the food which God gave him; in Moses, who forgot all his weariness when he heard that he should stretch out the rod in his hand and lead Israel through the sea dry and safe. Who would not gladly trust God when he thereby becomes fearless, courageous and strong?

If we look at Joseph, we shall learn another truth about the blessing of trusting God. How long did Joseph have to suffer trials? More than 13 years. That is a long time, as long, children, as you have lived. And when things go wrong, time passes, oh! so slowly. A single day seems very long, and how long would 13 years seem. Think of a sick person who must lie still for 13 years or more! Joseph might easily have become discouraged if he had not done what? If he had not trusted to God that He would eventually help him and restore him to honor. And because he trusted God, he would think like that other pious man who was also compelled to endure much, but who repeatedly said to himself, “I will trust in Thee”, etc. What did Joseph learn because he was sure that the end would be good? He learned to wait until God’s hour came. And that was no impatient waiting. When it has rained for a week, we often can not wait until the sun shines again. What kind of waiting is that? That is an impatient waiting. Joseph waited patiently and calmly. What do we generally do when we wait impatiently? We grumble and complain. Yes, sometimes our fault-finding and complaining is very much in evidence. On the other hand, what was Joseph like, since he did not murmur or complain? He was quiet and calm. And why could he be so quiet? Because he trusted in God above all things. Find the 62nd Psalm. John, read the first verse: “Truly, my soul waiteth”, etc. Who also could speak like the Psalmist? Joseph. His soul waited upon the Lord, of whom he was certain that He would help him when the proper time was at hand. What, then, does trust in God tend to make us also? It makes us patient and calm. Correct. And it does so because it makes us cheerful. Only that person who is cheerfully confident, who without doubt and fear looks forward to the end, can wait, be patient and still. Let us say that together: \emph{Trusting in God makes us cheerful, patient, and calm.} (Blackboard.) Children, that is one of the greatest virtues, to be cheerful, patient, and calm in trial and suffering. To undertake a difficult task in a strong, courageous, and fearless spirit means much; but to remain cheerful, patient, and calm amid long drawn-out trials—that is much more. And in whom only shall you find both? Only in him who trusts God above all things. What, then, must we learn? True trust in God.
\quellseite{623}
What did God do with Abraham, who trusted in Him? He led him into the Promised Land. What did He do with Joseph? He finally exalted him and made him ruler over all Egypt. What success did Moses have because he trusted in God? He let him lead Israel out of Egypt. What did He let David do? He let him conquer Goliath. Therefore, what did He always do? He always helped. Yes, He has always done all things well, and brought things to a successful end, for whom? For those who trusted Him above all things. Provided they always did what? Provided that they laid their hand in God’s hand that He might lead them, and provided they were sure that He is always the best guide. Trust in God, accordingly, makes us not only fearless, courageous, and strong when we have to undertake a hard task; likewise, not only cheerful, patient, and calm when we must suffer long and much; but what else does it do? It leads to a successful end. To what does it lead when we have a difficult task? That we may accomplish it. From what does it keep us under long, hard trial? From murmuring and complaining. Yes, from murmuring and complaining against God and man; from doubt of God’s help. And to what good end does it lead? That we shall be finally delivered out of all trials and once more have happiness and joy, health and success. Let us say that together! Trust in God always leads to a good end. Always? Children, is that not saying too much? Think of Lazarus! He certainly trusted God, and how did matters turn out with him? Badly. How badly? His body was covered with sores. Was there any other misfortune? He was poor and had nothing to eat, but the crumbs from the rich man’s table. He was poor and sick; and, children, we should not forget this: he had to see others about him strong and healthy and that one before whose door he lay living in mirth and splendor every day. That made his sickness and poverty all the more bitter for him. You say that Lazarus trusted God, and that trusting God always leads to a successful end. Did it lead to a successful end in the case of Lazarus? Yes. But how long did he remain sick and poor? Until his death. Did he then come to a successful end in this life? Do you call that a successful end when one must remain in sickness and poverty until his death, while others are healthy and rich? No. How can you then say that the trusting of Lazarus led to a successful end? Because this life does not end all. What further happened to Lazarus? His soul was carried up into Abraham’s bosom. What was its condition there? Unspeakable happiness. Yes, God so highly honored him and surrounded him with so much happiness that he no longer thought of what? Of
\quellseite{624}
his sickness and poverty on earth. But only when did matters come to a successful end? Not until after his death. Then is it correct that trusting in God leads to a successful end? Yes, it is correct. But how should we not understand it? As if God already in this world always lets things turn out well. Oh, no, children. God often does as in the cases of Abraham, Joseph, Moses, and David. When does He often permit the one who trusts Him to enjoy success? Here in this life. And what does He often do with trials if one trusts Him? He brings them to an end and again bestows health or riches or honor. But He does not always do this. Many a one who trusted in Him was not permitted to see the success of his endeavors; more than one pious Christian has died in poverty and wretchedness. But when would He always help him, give him honor and bring things to a successful end? After death. Where does He take those who in this life have trusted Him above all things? To Himself, into everlasting joy and glory. We must therefore add something to the sentence “trusting God leads to a good end”, in order that it may not be misunderstood. What words would you add? In this life or in the life to come. Yes, we will so note it and repeat it together a number of times: Trusting God leads to a good end, if not in this life, then in the life to come! (Blackboard.) When, especially, will the true children of God have to prove their trust in Him? When God leaves them in distress until the end of their lives. For what especially will they then trust God? That after death He will bring them to honor and lead them to joy and glory. In what words will they thereupon take new comfort? Turn to Psalm 73. Frank, read verse 23. Yes, so they will say, now all the more. And if their spirits become disquieted, they will say to themselves as we read in Psalm 42:11. What do you read there, John? “Why art thou cast down”, etc. Or with what the Lord says in Isaiah 49:15. What is said there? “Can a woman forget her suckling child”, etc. (Here have the hymns read, “Commit thy way confiding”, and “If thou but suffer God to guide thee”, discuss them briefly. Also Psalms 23 and 91.) Yes, children, so great and important is trust in God. What does it do in the first place? In the second? In the third? Should you not also trust God above all things, place your hands in His that He may lead you, and trust Him that He is always the best guide? Is His guidance good even when the task He has assigned you is great and hard; when the trials He sends last long; even when He lets you die in wretchedness and poverty? Yes, even then.

\quellseite{625}
\subsubsection{4. Application} Such great tests of trust in God as we find in the lives of Moses, Abraham, Joseph, or David, are not demanded of you children. You are not required to forsake friends and acquaintances and go alone into a strange land, the name of which you do not know, as did Abraham. You will not be sold into slavery like Joseph and, though innocent, have to lie for years in prison like he; you are not required to deliver a people out of Egypt and to lead them through the sea like Moses; nor need you, like David, go out to fight a giant. Though in later days hard tasks and severe trials may confront you, they are as yet small and insignificant. But, nevertheless, what shall you soon have to show? True trust in God. If you have a difficult problem to solve or an examination is before you, and you are anxious and fearful, what should you do? Trust in God and hope in His help. Here, certainly, also another thing is necessary. What do you first have to do, if you know that an examination will shortly be held? Study diligently and work hard. Of course, for God has no pleasure in lazy people; He will rather let them fail dismally, so that they may finally learn to be diligent. But often, though you had been diligent, not only occasionally but regularly, from day to day, yet, how did you feel when the examination approached? You felt afraid. Yes, your heart beat and leaped violently. How, then, can you quietly and serenely go to the examination? If you say to God: Dear Lord, be with me and help me to give the correct answers. For what does trust in God make one? It makes one fearless, courageous, and strong. A youth was working for a farmer and was sent into the field with the horses and machine. He was only sixteen years old and had just recently come from Germany. He had but little experience with such young, mettlesome horses, and had seen the machine in operation but once. How would that work out? His heart was heavy as he rode out to the field. The young horses were worse than usual, for the flies were plaguing them. And when he tried to set the machine, he could not, for he had forgotten how it was done. He tried again and again, but without success; the horses refused to stand any longer, and he could not discover how to operate the machine. He wished he were back again at the house or with his parents in Germany. Then he happened to remember that in the First Commandment it says that we should trust God above all things. Do you know what he did? Reins in hand, he knelt down beside the machine and prayed: “Lord, help me! Thou hast often helped my mother when she called upon Thee. Help me also!” Then his heart grew calm and
\quellseite{626}
his head clear. The next time he tried it he succeeded, and before the day was over he had accomplished considerable work. See, boys, that is the blessing of trusting God. You will soon be confirmed and more than one of you will be going forth from your homes to the service of others or to another school. That may not be pleasant and your hearts may be troubled and heavy. But what did Abraham do when he left for the strange land? He trusted God. And how did he then feel? He became fearless, courageous, and cheerful. What can you do that you likewise may be of good courage? I can trust God. Yes, you can place your hand in God’s hand that He may lead you, and you can be certain that He will lead you aright. Does that mean that it will always go well with you? No. Often people will not be satisfied with you, however much you may strive to please them. Of whom should you then think? Of Joseph. Yes, he did his best and was sent to prison. But what did he not do, despite all this? He did not forget God. Still more. What did he do with God’s hand when things went wrong with him? He held it the closer. And of what was he confident? That God was the best guide. And that it would lead to what kind of an end? To a successful end. So you must think when you have done your duty and men are not satisfied with you. Or perhaps your parents become ill and die. Upon whom will you then rely? Upon your wisdom and cleverness, upon your strong, healthy body, upon your friends and acquaintances, or upon God your heavenly Father? Upon my heavenly Father. Why not upon your own wisdom and strength? Because they will not suffice. Why not upon friends and acquaintances? They may die, and even though they live, they can not always help. Who can always help, however great the need? Our heavenly Father. What kind of father is He? An almighty father. If He should let you become ill and even die, while the others about you are healthy and happy, what would be difficult for you? To trust in God. Certainly. But if you pray to Him He will preserve your trust in Him, and you will hold firmly that even then He is guiding you aright. At death pious children of God more than ever lay their hands in God’s and say: “Abide with me; fast falls the eventide”, etc. Even in the valley of death, accordingly, they are not alone; but who is with them? God is with them. What words of David’s Shepherd Psalm may they, accordingly, apply to themselves? “Yea, though I walk through the valley of the shadow of death”, etc. Why is, also then, God
\quellseite{627}
the best guide? Because the time to enter heaven has come. Yes, children, that is truly the best guide who leads us to heaven. If God then leads us so well, what should be and shall be our chief care? We will then trust in God above all things.—By all means, children, we will do that!—Open your catechism, page 31. There read what stands under the heading “trust in Him”; and again on page 32. Also read the Scripture passages plainly and with proper emphasis: Frank, No. 15, Anna, No. 16, Emil, No. 17, etc. Now, study the lesson carefully for tomorrow. (Passages which were quoted in the first part of the catechization have already been learned between the First and the Second Parts; Cf. page 598, note.)

\subsection[4. The Nature of Faith]{4. Outline of a Catechesis on the Nature of Faith}

Aim: We want to learn what it means to believe.

\subsubsection{1. Developing the nature of faith from the intuitional material:} From the story of the departure of Abraham from his fatherland the class is to be shown how faith has to deal with the invisible (unknown land), but only with such invisible things of which God has spoken to us (God spoke with Abraham about the unknown land); that faith clings to these as real and firmly relies upon them—so firmly, indeed, that it will give up everything visible and material for them.—

\subsubsection{2. Comparison} The martyrs, who gave up the earthly, tangible, perceptible, in order to gain the invisible, everlasting life. Heb. 11.—

\subsubsection{3. Valuation} Proceeding from Heb. 11, especially the life of Abraham, and terminating in Jer. 5:3; Heb. 11:6; Gen. 15:6; Hab. 2:4.—

\subsubsection{4. Application} Experiences from the life of youth.

\subsection[5. The First Part of the Third Article]{5. Outline of a Catechesis on the First Part of the Third Article}

(For the Juniors and Seniors of the Academic Department of a College.)

Aim: We cannot by our own reason or strength come to faith in the Lord Jesus. \subsubsection{1. Paul supplies the intuitional basis} The catechist, with the co-operation of the class, develops a sketch of Paul’s life up to the time of his journey to Damascus in a manner in keeping with the age and knowledge of those in the grades mentioned. He leads them to Tarsus into the house of Paul’s father, who himself had been a Pharisee; gives them an idea of the instruction on the Law which Paul here received; has them accompany Paul,
\quellseite{628}
the youth, to Jerusalem and sit beside him at the feet of Gamaliel; shows them by particular examples his zeal in the fulfilling of the Law (Phil. 3:1–11 to be read in this connection), making it clear to them, on the basis of Gal. 2:16, how only one thought filled the mind of Paul—his seeking after righteousness. He then presents comprehensively how knowledge and zeal were not wanting; how Paul outstripped the majority; how reason and strength were there, and yet total inability to know Christ in spite of them and to attain righteousness by his natural strength. No wonder he wrote 1 Cor. 2:14 and 1 Cor. 12:3. If Christ Himself had not intervened, and had not called to him on the way to Damascus; if He had not enlightened his mind and spirit while his eyes were blinded; if He had not revealed Himself to him (Gal. 2:11, 17), i.e., removed the scales so that he could see and know Jesus; if He had not let him be baptized by Ananias and thereby become justified and sanctified, he would never have come to know Christ as his Lord, and thus, in spite of all his knowledge and strength, remained a poor, lost soul. But when he arose out of the washing of regeneration he could confess: “I believe that I cannot by my own reason or strength believe in Jesus Christ, my Lord . . . . . . . sanctified and kept me in the true faith.” Now he could write Eph. 2:1 and Titus 3:5 and Rom. 3:28.—

\emph{2. For Comparison} the scene at Caesarea Philippi may be used, where the Master asks the question: “Whom do men say that I the Son of Man am?” and, upon the confession of Peter, rejoins: “Flesh and blood hath not revealed it unto thee, but my father which is in heaven”; the Bible reading and lack of understanding on the part of the eunuch of Ethiopia until Philip opens the Scriptures to him. One may also remind the pupils of Athens, which, with all its wisdom, could not find God, built an altar to the unknown God, and regarded the message of the risen Christ as ridiculous.—

\subsubsection{3. Valuation} Is it a piece of good fortune to be led to Jesus, to be brought to faith in Him? Ask the Greeks who desired to see Jesus; ask the converted Paul; ask the baptized eunuch who went his way with joy; ask Stephen, face to face with death by stoning; go out into the mission field and ask the converted heathen; listen to the hymns of Jesus which the Church sings!—

\subsubsection{4. Application} To know Latin and Greek, and to have other knowledge is good, but to know Jesus Christ is better than all knowledge.
\quellseite{629}
All wisdom and power do not lead to God; thank God that you have been baptized, thereby sanctified and brought to Christ. Are you still with Him, and can you say with Luther, “and preserved me”?

\subsection[6. The Conclusion of the Third Article]{6. Outline of a Catechesis on the Conclusion of the Third Article\footnote{From Berndt, \emph{Methodik des Unterrichts in der evang. Religion}, 1909, page 105 f. According to the analytical method, to which Luther’s explanation is the sum of the entire catechization, p. 523 f.}}

\emph{I believe in the resurrection of the body and the life everlasting.}

\subsubsection{1. With what Christians may comfort themselves at the grave}

a) How did Jesus comfort sorrowing humanity? The son of the widow of Nain, the raising of Lazarus.—b) How did the Lord confirm to them His comfort? “I am the resurrection and the life”, etc., certainty of the eternal life of Jesus; His resurrection the pledge of ours. He who is in Jesus has passed through death unto life.—c) How has God written this comfort already in the work of nature? Autumn, Spring, 1 Cor. 15.—d) How did pious Christians comfort themselves? 1 Cor. 15.—Paul (2 Cor. 5)—Ignatius of Antioch and others—The church hymns; “Jesus my Redeemer lives”, etc.;—Luther at the deathbed of his Lena.—e) How did D. Martin Luther proclaim this comfort? “In which Christian Church He . . . . . . at the last day will raise up me and all the dead.”—Resurrection of the body = all people.

\subsubsection{2. In what Christians may rejoice in the hour of death}

a) What does the Lord say about our joy in heaven? Parables.—b) Why is eternal life a great joy for the Christian? Rev. 21:1–7; Communion with Christ, God, all the redeemed, freedom from all inner imperfection and outer restraint.—c) How pious Christians would rejoice in eternal life. Stephen—Paul in his letter to the Philippians—Augustine and Monica (\emph{Confess.} XI, 10)—“Jerusalem thou City Fair”—and like hymns.—d) How does Luther speak of this joy? “And give unto me and all believers in Christ everlasting life.”

\subsubsection{3. How Christians conduct themselves on their way to heaven}

a) The earnest sermon preached to us on this subject by Christ. Parable of the Wise and Foolish Virgins, and others.—b) How the faithful disciples followed His admonition. Peter (Legend by Kinkel)—Paul.—c) How Christians proceed to heaven. “Blessed are they that mourn, for they shall be comforted.”—Appropriate hymns.

\quellseite{630}
\subsection{7. Questions for Review and Drill}

What can God demand of us since He has become through Baptism our God? When do we hold Him to be our God?—We should fear Him, i.e., always hold Him before our eyes as whom? We should love Him as who loved Him?—And we shall trust Him as who trusted Him?—When did Moses trust God above all things? When again? And especially when?—What did Moses do when he was commissioned to go to Egypt to save Israel? (He placed his hand in God’s that He might lead him.) What trust did he thereby repose in God?—And what did Moses not immediately do with God’s hand when Pharao would not let Israel go?—But what did he do all the more?—What did he further trust Him to do?—When did God seem to be an erring guide?—What did Moses then do instead of casting aside his trust in God?—How did he speak to his disquieted soul? Recall the Scripture passages you have learned! Frank, you, too, may repeat them.—By what can we see that Moses was right when he refused to cast aside his trust in God?—Was Moses ever sorry that he had laid his hand in God’s that God might lead him?—Show how David trusted God.—Show how Abraham trusted God.—From whose life especially have we seen that we should trust God even amid trials?—How long did Joseph’s humiliation last?—How did he comfort himself when it seemed as though God had forgotten him?—In what passage of Scripture has God Himself said that He will not forget His children?—What did Joseph say to himself when all evidence seemed lacking that God was thinking of him? Think again of your Scripture passages!—When, therefore, should we not let go of God’s hand, but hold it the closer?—And what trust shall we have in Him even in the longest trials?—For these two things belong to true trust in God: 1. Place our hand in God’s hand that He may lead us; 2. Be confident that He is the best guide, even when we have a hard task to perform or when long and severe trials rest upon us. In which Scripture passage are we admonished to commit all our ways unto Him and to trust in Him?—Which hymn also contains the same admonition?—We should trust God above all things also because a great blessing accompanies such trust. How did Moses practice this?—What did David do when he went out against Goliath?—Or Abraham, when he was told to go out into the unknown land?—In which Scripture passage are we told that trust in God makes one strong, courageous, and fearless?—Peter, you also may repeat
\quellseite{631}
it. Anna, likewise!—What effect did Joseph’s trust in God have upon him in his long trial?—But I know of another thing which should make it easy for us to trust in God. Who can tell it to me?—But does trust in God always lead to a successful end?—Even if we must die in poverty and misery?—From whose life have we seen that one may be pious and trust God and still have to suffer in wretchedness until death?—But why can we say that Lazarus’ trust in God led to a successful end?—What should we therefore add to the sentence: ‘Trust in God leads to a successful end’?—What must we truly do if trust in God is to make us fearless, strong, and courageous to fulfil our task; if it makes us patient and cheerful in trial, and if it always, in this or at least in the life to come, leads to a successful end?—In which psalm has David declared that God has always led him aright?—In which psalm would Moses help us to find joy in trusting God?—Which two hymns especially are calls upon us to trust in God? (When they have been learned and explained, other questions regarding them are in order). Show me how also you can exhibit true trust in God.—Louis, you also may show this.—Emil, likewise.—Yes, children, to trust God in all things is one of the main points in Christianity. It makes us courageous in life, cheerful, calm, and patient in trial, confident in death. May God Himself awaken this within our hearts and preserve it until our end!

\subsection{8. An Excursus Through the Catechism}

Thanking and serving. According to L. Schultze, \emph{Katechet. Bausteine}.\footnote{There is much to be improved before it is really fit to be used.}

It is a good thing to give thanks unto the Lord, dear children! Therefore we find so much about thanksgiving in our Catechism. Where do we find the word “thank” for the first time? In the First Chief Part, in the Second Commandment: Call upon His name . . . and worship Him with . . . thanksgiving.—For what shall we thank God? For everything.—Where is that written? In the First Article of the Christian Creed. In which words? “For all of which I am in duty bound to thank and praise . . . obey him.” Tell what that “all” comprises which is here mentioned. That God has made me and all my goods.—How can that “all” with which God preserves body and life, be expressed in two words?—Daily bread.—Who has taught us so to call it? Jesus in the Lord’s Prayer.—In which
\quellseite{632}
petition? In the fourth.—What do we therefore, according to this petition, owe Him also for our daily bread? That we receive our daily bread with thanksgiving.—We have already had three Chief Parts in which the giving of thanks is dealt with. Why is thanksgiving nevertheless so often lacking? Give me the entire explanation of the Fourth Petition. “God gives daily bread . . . to know and to receive our daily bread with thanksgiving.” What, according to this, should precede the giving of thanks? “That He would lead us to know it.”—That is the point. Therefore a thankful person is one who sees and knows.—Who never forgot to give thanks? The Lord Jesus.—Prove this from the Catechism. “Our Lord Jesus Christ . . . took bread, gave thanks, brake it,” etc.—So we have another Chief Part that mentions the giving of thanks! Which one is it? The Fifth.—Did Jesus thank only for the bread? No; “after the same manner also He took the cup, when He had supped, gave thanks,” etc.—For what did Jesus give thanks? For everything.—But to thank with words or with the tongue is not all of thanksgiving. How should our gratitude express itself? Think of the First Article. For all of which I am in duty bound to thank and praise, to serve and obey Him.—Service! That is what you owe also to the Savior. For what purpose has He purchased you? “That I might be His own and live under Him in His kingdom and serve Him.”—You also owe this to men. The two parts, “serve and obey,” stand beside each other elsewhere in the Catechism. Where? In the Fourth Commandment. “Serve and obey them, hold them in love and esteem.”—Should you obey everybody? No; only parents and masters. But “serve”? Does that only apply to parents? No; to all men. We should also be of service to our neighbor. Where are we told this? In the Ninth Commandment: “But help and serve him in keeping it.”—To what does this here refer? To his inheritance or house.—Which other Commandment emphasizes this duty? The Seventh: “But help him to improve and protect his property and living.”—When especially are you bound to help him? In every need.—Where is that written? In the Fifth Commandment: “but help and befriend,” etc.—But what if the person be your enemy? “So we on our part will heartily forgive and readily do good to those who sin against us.”—Now we have learned something about thanking and serving from four Chief Parts. Which is lacking? The Fourth Chief Part—of Holy Baptism.—You will not find the word “thank” there, but the idea and thought is there of the very best thanksgiving, which is at the same time a daily serving. Which is that? “That the old Adam in us by . . . . . . . . . putrify forever.”

\quellseite{633}
\section{3. Instruction in Hymns}

The hymn: “\emph{Wake, arise! The voice is calling.}” (According to Schueren, page 119 ff.)

Twice this week the bell in the tower was tolled. Our fathers experienced times when the death knell was sounded every day in the year. Every day the bell tolled its message of death. In the century in which Luther lived, three times a fearful pestilence stalked through the German land—a pestilence so dreadful that in one visitation more than half of the inhabitants in many places died. Those were the unfortunate years 1529, 1572, and 1597. When the last pestilence swept through, there lived in Unna, a small city in the province of Westphalia, a pious Lutheran preacher named Nicolai. He had already suffered much in many places at the hands of the enemies of the followers of Luther. Then came the great pest. Daily he saw from 20 to 30 corpses carried past his house. He spent his time from early morning until late night among the dead and dying. In a short time 1400 persons were buried. Of what could he think but of death and the grave? It was midnight also for his congregation. The Bridegroom was near at hand to lead the wise virgins to the heavenly courts. All things indicated this. He himself was a watchman placed by God upon the towers to remind the congregation of the Bridegroom’s coming. And he was faithful to his trust. He admonished to repentance and then comforted the suffering with the prospect of the heavenly Jerusalem, the city of God with the pearly gates, and put it all in the form of a beautiful hymn. (Teacher recites the hymn.) Now open your books. Sing the first verse. How beautiful the melody and how appropriate to the words! Let us study the words of this hymn!

I will read the first verse of the hymn once more. A capable child then reads it, and then all in chorus. Of which parable does this verse remind you? It reminds of the Parable of the Ten Virgins.—Repeat this parable!—To what are the virgins in this verse admonished? They are admonished to be ready.—For what should they be ready? To meet the bridegroom.—How must the virgins be in order to meet the bridegroom? They must be awake.—What, then, do the watchmen do who are placed by God upon the towers? They call: Awake!—To whom do they thus call in this verse? They call to the city of Jerusalem.—Who is meant thereby? The people are meant thereby.—We shall see what people are meant. Who is the bridegroom? Christ is the bridegroom.—Whom does He wish to lead to the heavenly palace? Those who believe in Him and follow Him.—How do we call
\quellseite{634}
those people who do this? We call them Christians.—What are we therefore to understand by Jerusalem? The Christian congregation.—What is said to it? To awake and be ready.—Why? Because the bridegroom is coming.—When will He come? He will come at midnight.—When will that be? When may the Lord call you away? He may call me away at any hour.—Which hour, then, is the midnight hour? Every hour is the midnight hour.—When should I, therefore, be ready? I should always be ready.—What is meant by being ready and having one’s lamp filled and burning, you have already learned.

I shall read the first stanza again.—You read it!—And you!—And you!—Now the soprano! The alto! Right; now the second stanza should be easier. Treatment of the two following stanzas as before, but more briefly. In the second stanza attention must be called to the fact that the second part is a joyful response to the first part; and in the middle of the third stanza the teacher reads the corresponding Scripture verse from Revelation 21.—At the close the teacher reads the whole. Then the children read it together. The singing of the third stanza closes the lesson. It will hardly be necessary to tell the children to learn the hymn at home. Many already know it by memory, others will learn it without being commanded. During the review one can tell of the blessings which the hymn has brought into the life of many a person.

